% ALT 2027 submission. Keep alt2027.cls, jmlr.cls and jmlrutils.sty next to
% this file. Compiler: pdfLaTeX (Overleaf default).
\documentclass[anon]{alt2027} % submission (anonymized)
%\documentclass[final]{alt2027} % camera-ready (de-anonymized)

% Loaded by the class: amsmath, amssymb, natbib, graphicx, url, xcolor,
% hyperref, algorithm2e (environment algorithm2e), and theorem, lemma,
% proposition, corollary, definition, remark, conjecture, example and proof
% environments. The class refuses amsthm.
\usepackage{times}
\usepackage[T1]{fontenc}
\usepackage{booktabs}
\usepackage{multirow}
\usepackage{microtype}
\usepackage{enumitem}

% The class numbers theorems, lemmas, propositions, corollaries, definitions
% and remarks with one shared counter. Give lemmas, propositions,
% corollaries, definitions and remarks their own counters, so that the
% numbering reads Theorem 1, 2, 3, Lemma 1, 2, ... Delete this block to go
% back to the shared counter.
\makeatletter
\newcommand{\altrenumber}[2]{%
  \expandafter\let\csname #1\endcsname\relax
  \expandafter\let\csname end#1\endcsname\relax
  \newtheorem{#1}{#2}}
\makeatother
\altrenumber{lemma}{Lemma}
\altrenumber{proposition}{Proposition}
\altrenumber{corollary}{Corollary}
\theorembodyfont{\upshape}
\altrenumber{definition}{Definition}
\altrenumber{remark}{Remark}
\theorembodyfont{\itshape}
\providecommand{\qed}{\hfill\jmlrBlackBox}
% allow slightly loose lines rather than lines that run into the margin
\setlength{\emergencystretch}{3em}

% The class puts a theorem's title inside the list label, which cannot break
% across lines, so long titles ran into the margin. Same look, but the title
% is set as ordinary text and can wrap.
\makeatletter
\renewcommand*{\@opargbegintheorem}[3]{%
  \ifdef{\@jmlr@currentthm}%
  {%
    \letcs{\jmlr@this@theoremheader}{jmlr@thm@\@jmlr@currentthm @header@font}%
    \letcs{\jmlr@this@theorembody}{jmlr@thm@\@jmlr@currentthm @body@font}%
    \letcs{\jmlr@this@theoremsep}{jmlr@thm@\@jmlr@currentthm @sep}%
    \letcs{\jmlr@this@theorempostheader}%
       {jmlr@thm@\@jmlr@currentthm @postheader}%
  }%
  {%
    \let\jmlr@this@theorembody\@theorembodyfont
    \let\jmlr@this@theoremheader\@theoremheaderfont
    \let\jmlr@this@theoremsep\@theoremsep
    \let\jmlr@this@theorempostheader\@theorempostheader
  }%
  \trivlist
   \item[\hskip\labelsep{\jmlr@this@theoremheader #1\ #2}]%
   {\jmlr@this@theoremheader (#3)\jmlr@this@theorempostheader}\ %
   \mbox{}\jmlr@this@theoremsep
   \jmlr@this@theorembody
}
\makeatother

% Blackboard bold sets
\newcommand{\R}{\mathbb{R}}
\newcommand{\N}{\mathbb{N}}
\newcommand{\Z}{\mathbb{Z}}
\newcommand{\Q}{\mathbb{Q}}
\newcommand{\C}{\mathbb{C}}

% Labels
\newcommand{\alg}{\mathrm{alg}}
\newcommand{\meta}{\mathrm{meta}}
\newcommand{\fix}{\mathrm{fix}}

% Probability
\newcommand{\Prb}{\mathbb{P}}
\newcommand{\E}{\mathbb{E}}
\DeclareMathOperator{\Var}{Var}
\DeclareMathOperator{\Cov}{Cov}

% Operators
\DeclareMathOperator{\diag}{diag}
\DeclareMathOperator{\supp}{supp}

\newcommand{\argmax}{\operatorname*{arg\,max}}
\newcommand{\argmin}{\operatorname*{arg\,min}}
\newcommand{\ind}{\mathbf{1}}

% Inner product
\newcommand{\inner}[2]{\left\langle #1, #2 \right\rangle}

% Calligraphic sets
\newcommand{\cA}{\mathcal{A}}
\newcommand{\cS}{\mathcal{S}}
\newcommand{\cT}{\mathcal{T}}
\newcommand{\cF}{\mathcal{F}}
\newcommand{\cE}{\mathcal{E}}

% RL notation
\newcommand{\SSS}{\mathcal{S}}
\newcommand{\AAA}{\mathcal{A}}
\newcommand{\PP}{\mathcal{P}}
\newcommand{\RRR}{\mathcal{R}}
\newcommand{\gammaRL}{\gamma}

% Confidence bounds
\newcommand{\olQ}{\overline{Q}}
\newcommand{\ulQ}{\underline{Q}}
\newcommand{\olV}{\overline{V}}
\newcommand{\ulV}{\underline{V}}

% Misc
\newcommand{\eps}{\varepsilon}
\newcommand{\dt}{\mathrm{d}t}
\newcommand{\given}{\,\bigl\vert\,}

% Fractions
\newcommand{\half}{\tfrac12}
\newcommand{\thalf}{\tfrac32}

% Notes
\newcommand{\note}[1]{\paragraph{Note.} #1}



\newcommand{\todo}[1]{\textcolor{red!70!black}{[\textbf{TODO:} #1]}}
\newcommand{\Cold}{\mathsf{C}}
\newcommand{\Exp}{\mathsf{E}}
\newcommand{\Hyb}{\mathsf{H}}
\newcommand{\wid}{w}

\title[Information Value Is Not Policy Value]{Information Value Is Not Policy Value:\\ When LLM Warm-Starts Help Linear Contextual Bandits}

% Hidden by the [anon] option. Fill in for the camera-ready version.
\altauthor{%
 \Name{Author Name1} \Email{abc@sample.com}\\
 \addr Address 1
 \AND
 \Name{Author Name2} \Email{xyz@sample.com}\\
 \addr Address 2%
}

\begin{document}
\maketitle
\begin{abstract}
Routing the early rounds of a linear contextual bandit to a frozen LLM and
then handing control to LinUCB can beat both, sometimes by a large factor,
and the LLM's accuracy does not predict when. We explain why. The expert's
routed rounds reach the learner in two ways: through \emph{coverage}, the
directions its actions explore, and through the \emph{transfer}, what its
observed rewards say beyond that coverage. Coverage alone cannot produce
large gains when LinUCB explores as the theory requires: on any instance with
a fixed finite arm set in which the best arm is always available, the saving
it certifies over cold LinUCB is capped by the ratio of the prices the two
pay for coverage, and the cap is attained. Under linear rewards with Gaussian
noise, experts of equal coverage are interchangeable. The transfer is not
capped. When LinUCB explores less than the theory requires, as it does in
practice, we construct two experts with identical coverage, cost and number
of right answers, the losing one possibly the more accurate, whose
multipliers over cold LinUCB differ in most runs by a factor that grows with
the horizon, and in expectation once the expert gives about $\log T$ right
answers. On AmazonCat-13K, three LLM experts of similar, low accuracy
multiply the router's accuracy over cold LinUCB by $1.71$, $2.72$ and $3.29$,
cutting its regret by $6.5$ to $21\%$. Coverage alone saves about $2$ to
$3.5\%$ of cold LinUCB's regret for each; the transfer, which steers the
learner to the most frequent category, carries the rest and every
significant difference between experts, and hiding the rewards of a handful
of right answers on that category, while keeping their coverage, removes
their gain. Information value is not policy value.
\end{abstract}

\begin{keywords}%
linear contextual bandits, LinUCB, warm start, large language models, exploration, regret
\end{keywords}


\section{Introduction}
\label{sec:intro}

Large pretrained language models carry prior knowledge that can guide the
early decisions of a sequential learner
\citep{brown2020language,ouyang2022training,achiam2023gpt,alamdari2024jumpstarting,abbasiyadkori2024believe},
in recommendation \citep{li2010contextual,agarwal2016making}, healthcare
\citep{tewari2017ads,svensson2023sequential} and dialogue \citep{li2016deep}.
A striking empirical finding is that combining a frozen LLM with an adaptive
linear bandit can strictly dominate both, sometimes by a large factor
\citep{chen2024llmdecision,huangcontextual}: routing the early rounds to the
LLM and then handing control to LinUCB beats the LLM alone and LinUCB started
from scratch. The literature on warm starts and model selection
\citep{bastani2021mostly,zhang2019warm,agarwal2017corralling} reports nothing
comparable, and the LLM's accuracy is a poor guide: on AmazonCat-13K, three
LLM experts that are right on $3.6$--$5.3\%$ of contexts raise the accuracy of
the same learner by factors from $1.71$ to $3.29$ (cutting its regret by $6.5$
to $21\%$), and the least accurate is not the least useful. This paper
explains where the gain comes from and why accuracy cannot predict it.

\paragraph{Policy value and information value.}
The \emph{policy value} of an expert is the reward its own actions earn; its
\emph{information value} is how much its actions teach a learner that keeps
their data. The two can be wildly different. On the instance of
Proposition~\ref{prop:three-way}, a constant policy that is wrong on every
round produces data under which LinUCB's regret falls from linear to
constant: no policy value, enormous information value. That data, however, is
collected before the evaluated horizon and not paid for. When the expert's
rounds count toward the regret, as in a handoff, the data it leaves behind is
worth nothing on the same instance, for every expert: after the handoff
LinUCB does exactly what cold LinUCB does (Proposition~\ref{prop:one-pass}). The question
is when an \emph{online} handoff can multiply, and by how much.

\paragraph{The accounting.}
Write $\Cold$ for cold LinUCB, $\Exp$ for the frozen expert, and $\Hyb_\tau$
for the hybrid that follows $\Exp$ for $\tau$ rounds, feeding every
observation to LinUCB, and then hands over. Then
\[
R_T(\Cold)-R_T(\Hyb_\tau)=A_\tau+I_\tau,\qquad R_T(\Exp)-R_T(\Hyb_\tau)=B_\tau
\]
(Proposition~\ref{prop:decomp}), where $A_\tau$ is the expert's advantage on
the routed rounds, $I_\tau$ the regret LinUCB saves afterwards because of the
data it inherited, and $B_\tau$ the learner's advantage over the expert on
the tail. The inherited data enters LinUCB through two objects: the design
matrix, which sets every confidence width, and the response sum, which sets
the estimate. The first is \emph{coverage}; beyond coverage, the second
contributes a single expert-dependent term, the \emph{transfer}.

\paragraph{Coverage has a price.}
With a valid confidence radius, LinUCB's regret is limited by the confidence
widths of actions it would otherwise mistake for the best, and widths shrink
only with coverage. Coverage is never free: the \emph{price} of a direction is
the least regret any action pays per unit of coverage of it, and cold LinUCB
buys the same coverage on its own, at the price of the arms its optimism
forces it to pull. Theorem~\ref{thm:general-price} holds for LinUCB with a
valid confidence radius, on any instance with a fixed finite arm set in which
the best arm is always available: the multiplier of any policy over cold
LinUCB is at most the ratio of cold LinUCB's price to the expert's price on
the arms the warm start certifies, plus cold LinUCB's regret on the arms it
does not certify. It thus caps the certified saving rather than the
multiplier outright, and a large certified saving requires an action of small
gap that covers a direction cold exploration can buy only dearly
(Corollary~\ref{cor:diverge}). The cap is attained (Theorem~\ref{thm:price}),
and a catalog of eleven instances (Appendix~\ref{app:catalog-statements})
shows that removing any one of recurrence, cheapness, correct targeting, or
the expert's suboptimality on the tail can remove the gain. Under valid
exploration, then, coverage multiplies only through cheap informative
actions, and it cannot explain experts of similar accuracy and similar
coverage whose multipliers differ severalfold.

\paragraph{The transfer decides (main result).}
Under linear rewards with Gaussian noise, and for an expert whose choices do
not depend on the noise, the response sum is determined in law by the Gram
matrix (Proposition~\ref{prop:response-sum}), so experts that cover the same
directions are interchangeable (Corollary~\ref{cor:coverage-complete}). This
fails in practice for two reasons: rewards are not linear in the features,
and an LLM reads text the features only summarize, so it can prefer an action
on exactly the rounds where that action pays. Either way the estimate at
handoff acquires a transfer term (Definition~\ref{prop:transfer}). Our main
result (Theorem~\ref{thm:transfer-separation}) shows that when LinUCB explores
less than the theory requires, as it does in practice, the transfer can
decide everything: two experts with identical coverage, cost and number of
right answers on the routed rounds can have multipliers over cold LinUCB that
differ, in most runs, by a factor growing with the horizon, according to
where their right answers fall, and the losing expert can be the more
accurate one. A handful of right answers on an action that pays often settles
the learner on it, with a failure probability that decays exponentially in
their number; about $\log T$ of them make the separation hold in expectation.
With a valid radius the effect disappears. The two theorems therefore
describe two regimes: when exploration is valid, coverage governs and its
value is capped by prices; when exploration is reduced, as in practice, the
transfer can exceed that cap.

\paragraph{What the data shows.}
On AmazonCat-13K the three experts above cover the same directions to within
$5\%$, and coverage alone saves about $2$ to $3.5\%$ of cold LinUCB's regret
for each; every significant difference between them disappears when the
transfer is removed. The transfer acts as Theorem~\ref{thm:transfer-separation}
describes: one category is present on $30\%$ of contexts, a learner that
settles on it ends near $28\%$ accuracy against cold LinUCB's $8\%$, and
whether it settles is decided by a few of the expert's routed answers; adding
$15$ correct answers on that category to the weakest expert doubles its
accuracy multiplier, while the same answers with their rewards hidden from the
learner, which keeps their coverage, leave it unchanged. On OneShotWikiLinks, where no entity covers more than
$0.36\%$ of contexts, the Flan-T5 experts gain almost nothing (accuracy
multipliers $0.93$ to $1.05$).

\paragraph{Contributions.}
\begin{itemize}[itemsep=1pt, topsep=1pt, leftmargin=*]
\item \textbf{The transfer can decide everything}
(Theorem~\ref{thm:transfer-separation}): with reduced exploration, two
experts equal in coverage, cost and number of routed right answers can differ
by a factor of order the horizon, and the more accurate one can lose.
\item \textbf{Coverage is capped by prices}
(Theorems~\ref{thm:general-price} and~\ref{thm:price},
Corollary~\ref{cor:coverage-complete}): with a valid radius the certified
saving of any policy at any cutoff is bounded by a price ratio, the bound is
attained, and under linear rewards with Gaussian noise coverage is all there
is; on the canonical instance an offline warm start multiplies by a factor
of the horizon while an online handoff inherits nothing
(Propositions~\ref{prop:three-way} and~\ref{prop:one-pass}).
\item \textbf{Experiments that isolate the transfer}
(Section~\ref{sec:experiments}): on two datasets and twelve configurations,
with every tuning parameter chosen on seeds disjoint from the evaluation
seeds, coverage gives a small gain that is the same for every expert, the
transfer carries most of the gain and every significant difference between
experts, and right answers on the most frequent category double the
multiplier only when their rewards reach the learner.
\end{itemize}
Section~\ref{sec:related-main} discusses related work,
Section~\ref{sec:setup} fixes notation, Section~\ref{sec:decomp} gives the
accounting, Sections~\ref{sec:price}--\ref{sec:two-channels} the theory, and
Section~\ref{sec:experiments} the experiments. The catalog and all proofs
are in the appendix.

\section{Related Work}
\label{sec:related-main}

\textbf{LLMs and bandits.} \citet{chen2024llmdecision} combine an LLM policy
with a contextual bandit through an online model-selection master and report
reward multiplication; \citet{huangcontextual} calibrate LLM-derived priors to
set exploration bonuses; \citet{jumpstart2026} warm-start a linear bandit from
LLM-generated preferences through a single prior-error term; and
\citet{calibgated2026} inject LLM pseudo-observations into LinUCB with a
calibration-dependent weight. We separate the design matrix from the response
sum, bound what the first can deliver by a price ratio, and give the
expert-dependent remainder in closed form; our decomposition and bound hold
for any routed set, so they apply to these rules as well.
\textbf{Warm starts and model selection.} Warm-start analyses
\citep{bastani2021mostly,zhang2019warm} bound the regret of a learner started
from logged data through the quality of that data; they are offline in our
sense, and our contribution relative to them is the online accounting, in
which acquiring the data is charged and can cost more than it is worth. In
hybrid and offline-to-online reinforcement learning
\citep{song2022hybrid,wagenmaker2023leveraging} offline data helps when it covers what
the online learner needs, and with demonstrations its value depends on the
demonstrator's quality \citep{hao2023leveraging}; here the data is acquired
online, and its value depends on which of the expert's answers are right
rather than on its accuracy. Conservative bandits \citep{wu2016conservative} keep a learner close to a baseline policy; here the
baseline is a source of data.
Corralling and model selection
\citep{agarwal2017corralling,pacchiano2020model,marinov2021pareto} compete
with the best of several base learners; we ask when a combination beats both
components, and our bound applies to a CORRAL master as to any other policy.
\textbf{Geometry of exploration and misspecification.} The price of
Section~\ref{sec:price} is the single-direction shadow of the allocation
programs behind instance-dependent lower bounds
\citep{graves1997asymptotically,lattimore2017end}; our use is comparative,
with two policies paying different prices for the same coverage. Unlike a
hint about $\theta^\star$ \citep{cutkosky2022leveraging}, whose value grows
with its accuracy, a warm start is data in the design matrix, and its value
is not a function of accuracy. Work on misspecified linear bandits treats the
part of the reward the linear model misses as a cost to be controlled
\citep{ghosh2017misspecified,lattimore2020learning,foster2020adapting}; here that part, carried by the transfer,
is the main source of the gain. Appendix~\ref{sec:related} gives an extended
discussion.

\section{Setup}
\label{sec:setup}
\label{sec:preliminaries}

We work in a standard linear contextual bandit
\citep{abbasi2011improved,chu2011contextual,lattimore2020bandit}. At each round
$t=1,\dots,T$, the learner observes a finite action set $\mathcal A_t$ where
each $a\in\mathcal A_t$ has a feature vector $x_{t,a}\in\mathbb R^d$ with
$\|x_{t,a}\|_2\le L$, selects $A_t\in\mathcal A_t$, and observes
$r_t=\langle\theta^\star,x_{t,A_t}\rangle+\eta_t$, where $\theta^\star$ is
unknown with $\|\theta^\star\|_2\le S$ and $\eta_t$ is conditionally mean-zero
and $\sigma$-sub-Gaussian. Write $\mu_t(a)=\langle\theta^\star,x_{t,a}\rangle$,
$a_t^\star\in\arg\max_a\mu_t(a)$, and $\Delta_{t,a}=\mu_t(a_t^\star)-\mu_t(a)$
for the suboptimality gap of $a$ at round $t$. For a policy $\pi$ and an
interval $I\subseteq[T]$ of rounds, $R_I(\pi)=\sum_{t\in I}\Delta_{t,\pi_t}$ is
its pseudo-regret on $I$, and $R_T(\pi)=R_{[T]}(\pi)$. The language aspect of
the applications we target enters only through the feature map: context and
action embeddings from pretrained models are combined into $x_{t,a}$, and the
frozen expert may be an LLM. Nothing in the theory uses language structure.

\paragraph{LinUCB and its warm start.}
LinUCB maintains $V_t=V_1+\sum_{s<t}x_{s,A_s}x_{s,A_s}^\top$,
$c_t=c_1+\sum_{s<t}r_sx_{s,A_s}$, $\widehat\theta_t=V_t^{-1}c_t$, and plays
$A_t\in\arg\max_a\langle\widehat\theta_t,x_{t,a}\rangle+\beta_t\|x_{t,a}\|_{V_t^{-1}}$,
where $\|x\|_{V^{-1}}=\sqrt{x^\top V^{-1}x}$ is the confidence width of $x$ and
$\beta_t$ is the confidence radius. \emph{Cold-started} LinUCB, written
$\Cold$, uses $V_1=\lambda I_d$, $c_1=0$. A \emph{warm start} $(M,c)$ uses
$V_1=\lambda I_d+M$ and $c_1=c$ with $M=\sum_sx_sx_s^\top\succeq0$ and
$c=\sum_sr_sx_s$ the Gram matrix and response sum of some logged data. We say
the warm start is \emph{offline} if that data is not counted in the evaluated
horizon and \emph{online} if it is.

\paragraph{Frozen expert and handoff hybrid.}
A frozen expert $\Exp$ proposes $a_t^{\Exp}\in\mathcal A_t$ at each round by a
rule that is never updated (for an LLM, the action whose embedding is closest
to the LLM's output; see Section~\ref{sec:exp-setup}). For $\tau\in\{0,\dots,T\}$
the \emph{handoff hybrid} $\Hyb_\tau$ plays $a_t^{\Exp}$ for $t\le\tau$, updates
$(V_t,c_t)$ with every played action and observed reward including those
rounds, and plays LinUCB for $t>\tau$. Thus $\Hyb_0=\Cold$ and $\Hyb_T=\Exp$,
and the LinUCB phase of $\Hyb_\tau$ is LinUCB with the online warm start
$(M_\tau,c_\tau)=\bigl(\sum_{t\le\tau}x_{t,A_t}x_{t,A_t}^\top,\sum_{t\le\tau}r_tx_{t,A_t}\bigr)$.
The deterministic router used in our experiments
(Appendix~\ref{app:meta-algorithm}) chooses $\tau$ adaptively and may
interleave the two; the theory below is for a fixed $\tau$, which is the
right object for understanding why a multiplier appears; Appendix~\ref{sec:handoff}
comments on adaptive $\tau$.

\paragraph{Key terms.}
Six quantities recur. The \emph{confidence width} of an action with features
$x$ under a design matrix $V$ is $\|x\|_{V^{-1}}$; LinUCB plays a suboptimal
action only while $2\beta_t$ times its width is at least its gap
(Lemma~\ref{lem:width-limited}). The \emph{coverage} of a direction $u$ by a
sequence of played actions is $\sum_s\langle x_s,u\rangle^2$, the growth of the
design matrix along $u$; widths along $u$ fall only as its coverage grows
(Lemma~\ref{lem:loewner}, Appendix~\ref{app:widths}). The \emph{price} of $u$
is the least regret any available action pays per unit of coverage of $u$
(Section~\ref{sec:price}). A design
matrix \emph{certifies} an action once $2\beta$ times its width is below its
gap, so that the confidence bound alone rules the action out. The
\emph{acquisition regret} of a warm start is the regret of the routed rounds
that produced it. Finally, the \emph{multiplier} of a policy $\pi$ over cold
LinUCB is $R_T(\Cold)/R_T(\pi)$; the experiments also report the ratio of
accuracies used in the literature (Section~\ref{sec:experiments}).


\section{Where the Gain Comes From}
\label{sec:decomp}

The catalog and the experiments are organized by one identity.

\begin{proposition}[Decomposition identity]
\label{prop:decomp}
For any instance, any expert $\Exp$, any $\tau$, and any realization of the
noise,
\begin{align*}
R_T(\Cold)-R_T(\Hyb_\tau)
&=\underbrace{R_{[1,\tau]}(\Cold)-R_{[1,\tau]}(\Exp)}_{A_\tau}
 +\underbrace{R_{[\tau+1,T]}(\Cold)-R_{[\tau+1,T]}(\Hyb_\tau)}_{I_\tau},\\
R_T(\Exp)-R_T(\Hyb_\tau)
&=\underbrace{R_{[\tau+1,T]}(\Exp)-R_{[\tau+1,T]}(\Hyb_\tau)}_{B_\tau}.
\end{align*}
\end{proposition}

The identity is elementary; its use is to name the three terms.
$A_\tau$ is the expert's \emph{direct} advantage: positive if the expert earns
more than cold LinUCB would have on the routed rounds, negative if the expert's
actions cost more regret than cold exploration. $I_\tau$ is the
\emph{information} value of the routed rounds: after round $\tau$ the hybrid
runs LinUCB from the warm start $(M_\tau,c_\tau)$ that the routed rounds left,
while cold LinUCB runs the same algorithm from the state its own first $\tau$
rounds left, so $I_\tau$ compares the two starting states on the same
future. $B_\tau$ is the adaptive learner's advantage on the tail and is large
only when the expert remains substantially suboptimal late while the
warm-started learner does not. We write $R_{[1,\tau]}(\Exp)=\sum_{t\le\tau}\Delta_{t,a_t^{\Exp}}$
for the \emph{acquisition regret} of the warm start: the price, counted in the
online horizon, of the data in $M_\tau$.

\begin{remark}
\label{rem:decomp}
The identities hold pathwise in the reward-table model, in which the $k$-th
pull of an action always receives the same noise draw, so that two policies
can be run on the same noise \citep[\S4.6]{lattimore2020bandit}, and
therefore also in expectation. They also hold when the routed rounds are not
an initial segment: for any set $\mathcal T\subseteq[T]$ of routed rounds,
fixed in advance or chosen adaptively (as by the router of
Appendix~\ref{app:meta-algorithm}), replacing $[1,\tau]$ by $\mathcal T$
gives the same identities, because each side is a difference of sums over a
partition of $[T]$. If the warm start $(M,c)$ is built from data outside the
horizon (an \emph{offline} warm start $\Hyb^{\mathrm{off}}$, LinUCB started
from $(\lambda I+M,c)$), the routed rounds are not charged and the identities
read $R_T(\Cold)-R_T(\Hyb^{\mathrm{off}})=:I^{\mathrm{off}}$ and
$R_T(\Exp)-R_T(\Hyb^{\mathrm{off}})=:B^{\mathrm{off}}$. A multiplier over both
base policies means that $A_\tau+I_\tau$ is almost all of $R_T(\Cold)$ and
$B_\tau$ almost all of $R_T(\Exp)$; every diverging ratio in this paper is
read along a sequence of instances.
\end{remark}

\paragraph{What a warm start can change.}
\label{sec:width}
Fix the confidence event $\mathcal E_\delta$ on which
$\|\widehat\theta_t-\theta^\star\|_{V_t}\le\beta_t$ for all rounds $t$ at
which LinUCB acts, where $(V_t,c_t)$ are built from all data so far including
any warm start $(M,c)$ formed from the same kind of data; by the
self-normalized bound of \citet{abbasi2011improved} this event has probability
at least $1-\delta$ with
$\beta_t=\sigma\sqrt{2\log(\det(V_t)^{1/2}\det(\lambda I)^{-1/2}/\delta)}+\sqrt\lambda S$.
Define, for a positive definite $V$,
\begin{equation}
\label{eq:phi}
\Phi_t(V)\;:=\;\max_{a\in\mathcal A_t}\;\min\Bigl\{\Delta_{t,a},\;2\beta_t\|x_{t,a}\|_{V^{-1}}\Bigr\}.
\end{equation}

\begin{lemma}[Width-limited regret]
\label{lem:width-limited}
On $\mathcal E_\delta$, at any round $t$ at which LinUCB acts with design
matrix $V_t$, the instantaneous pseudo-regret satisfies
$\Delta_{t,A_t}\le\Phi_t(V_t)$. Moreover $\Phi_t$ is monotone: $V\preceq V'$
implies $\Phi_t(V')\le\Phi_t(V)$ when the radius $\beta_t$ is held fixed.
\end{lemma}

A competitor with a large gap but a larger width is as dangerous as its gap,
one with a tiny gap is cheap however wide it is, and the only way to lower the
regret LinUCB can incur on a round is to shrink the widths of competitors
whose widths exceed their gaps. Widths shrink through a warm start and through
LinUCB's own observations, both by exact linear algebra
(Lemmas~\ref{lem:loewner} and~\ref{lem:transfer},
Appendix~\ref{app:extra-theory}). One consequence is that the
determinant-ratio bound behind $\widetilde O(\sqrt{rT})$-type guarantees
cannot certify a multiplier: a multiplier that a width-limited bound can
certify must come from competitors whose widths the warm start pushes
\emph{below their gaps}
(Proposition~\ref{thm:handoff}, Appendix~\ref{sec:handoff}).

\section{The Price Bound}
\label{sec:price}

\paragraph{Why a price.}
By Lemma~\ref{lem:width-limited}, the routed rounds can lower LinUCB's later
regret only by shrinking widths, and widths shrink only where the routed
actions add coverage. Coverage is never free: covering a direction means
playing actions that point along it, and those actions have gaps. Cold
LinUCB buys the same coverage on its own---optimism makes it pull each arm
whose width exceeds its gap until that width falls below the gap---so the
expert and cold LinUCB buy the same thing, and the question is what each pays
per unit. That exchange rate is the price. The theorem below says that no
policy beats cold LinUCB by more than the ratio of the two prices on the
arms it certifies, plus what cold LinUCB pays on the arms nobody certified.

\paragraph{Prices.}
For a direction $u\in\R^d$ and an
action set $\mathcal A_t$, write
\begin{equation}
\label{eq:price}
\kappa_t(u)\;:=\;\min\Bigl\{\frac{\Delta_{t,a}}{\langle x_{t,a},u\rangle^2}:\;a\in\mathcal A_t,\ \langle x_{t,a},u\rangle\neq0\Bigr\},
\qquad \kappa(u):=\inf_t\kappa_t(u),
\end{equation}
the least pseudo-regret any policy must pay at round $t$ to add one unit to
the coverage $\sum_s\langle x_{s,A_s},u\rangle^2$ of $u$. Coverage is what
reduces widths (Lemma~\ref{lem:loewner}), so $\kappa(u)$ is the exchange rate
between regret and width reduction in direction $u$. The following lemma
says that no policy, expert or otherwise, gets coverage for less.

\begin{lemma}[No free coverage]
\label{lem:no-free-coverage}
For any policy $\pi$, any interval $I$ of rounds, and any direction $u$,
$\sum_{t\in I}\langle x_{t,\pi_t},u\rangle^2\le R_I(\pi)/\kappa(u)$.
Consequently, if $\Hyb_\tau$ hands off with design matrix $V_{\tau+1}=\lambda I+M_\tau$, then
for every unit vector $u$,
$\|u\|^2_{V_{\tau+1}^{-1}}\ \ge\ 1\big/\bigl(\lambda+R_{[1,\tau]}(\Exp)/\kappa(u)\bigr)$.
\end{lemma}

The lemma is one half of the theorem; the other is that cold LinUCB
pays its own price for each arm and stops once the arm's width falls below
its gap.

\paragraph{Setting for the theorem.}
In words: we fix a finite set of arms in which the best arm is always
offered, compare cold LinUCB with an arbitrary policy $\pi$ up to a cutoff
$t_0$, and measure how much coverage $\pi$ bought, at what price, and which
arms that coverage lets the confidence bound rule out. Let
$\mathcal X\subset\R^d$ be a finite set of arm vectors with $\|x\|_2\le L$,
let $\mathcal A_t\subseteq\mathcal X$ for every $t$, and suppose an arm
$x^\star$ with $\mu(x^\star)=\max_{x\in\mathcal X}\mu(x)$ belongs to every
$\mathcal A_t$. The gap of an arm is then the same at every round,
$\Delta(x):=\mu(x^\star)-\mu(x)$; write $\mathcal S:=\{x\in\mathcal X:\Delta(x)>0\}$
and fix $\beta\ge\beta_t$ for every $t\le T$ (one exists because
$\|x\|_2\le L$ bounds $\det V_t$).

\emph{What cold LinUCB pays.} For $x\in\mathcal S$ put $u_x:=x/\|x\|_2$ and
\begin{gather*}
\kappa_x:=\frac{\Delta(x)}{\|x\|_2^2},\qquad
C^\star_x:=\frac{4\beta^2\|x\|_2^2}{\Delta(x)^2},\\
n^\star_x:=\Bigl(\frac{C^\star_x-\lambda}{\|x\|_2^2}+1\Bigr)\ind\{C^\star_x\ge\lambda\},\qquad
N_x:=\#\{t\in[T]:x\in\mathcal A_t\}.
\end{gather*}
Here $\kappa_x$ is the price of covering $u_x$ by pulling $x$ itself,
$C^\star_x-\lambda$ is the least coverage of $u_x$ at which the confidence
bound can exclude $x$, and $n^\star_x$ bounds how often cold LinUCB pulls $x$
before its own pulls exclude it.

\emph{What $\pi$ pays.} Write $x^\pi_t$ for the arm $\pi$ plays at round $t$
and, for a set $U\subseteq\mathcal S$ of arms,
$\mathrm{cov}_U(x):=\sum_{y\in U}\langle x,u_y\rangle^2$. The price $\pi$ pays
for coverage of $U$ up to $t_0$, and the coverage of $U$ it obtains for free
by playing optimal arms, are
\begin{gather*}
\kappa^\pi_U(t_0):=\min\Bigl\{\frac{\Delta(x^\pi_t)}{\mathrm{cov}_U(x^\pi_t)}:
\ t\le t_0,\ \Delta(x^\pi_t)>0,\ \mathrm{cov}_U(x^\pi_t)>0\Bigr\},
\\
C^{\mathrm{free}}_U(t_0):=\sum_{t\le t_0:\,\Delta(x^\pi_t)=0}\mathrm{cov}_U(x^\pi_t),
\end{gather*}
with the minimum over an empty set equal to $+\infty$. Finally
$U_{\mathrm{cert}}(t_0)$ is the set of arms $x\in\mathcal S$ that $\pi$'s
first $t_0$ rounds certify: $2\beta\|x\|_{(V^\pi_{t_0+1})^{-1}}<\Delta(x)$, with
$V^\pi_{t_0+1}=\lambda I+\sum_{t\le t_0}x^\pi_t (x^\pi_t)^\top$. For the
handoff hybrid with $t_0=\tau$, $\kappa^\pi_U(\tau)$ is the expert's price and
$R_{[1,t_0]}(\pi)$ its acquisition regret, but nothing below uses this.

\begin{theorem}[Price bound for the online multiplier]
\label{thm:general-price}
In the setting above, on the confidence event of cold LinUCB, let $\pi$ be
any policy, $t_0\in\{0,\dots,T\}$ any cutoff and $U\subseteq
U_{\mathrm{cert}}(t_0)$ any set of certified arms, and write
$\kappa^{\Cold}_U:=\max_{x\in U}\kappa_x$ ($:=0$ if $U=\emptyset$) and
$P_U:=\kappa^{\Cold}_U/\kappa^\pi_U(t_0)$ for the price ratio. Then
\[
R_T(\Cold)\ \le\ P_U\,R_{[1,t_0]}(\pi)\ +\ \kappa^{\Cold}_U\,C^{\mathrm{free}}_U(t_0)\ +\ \sum_{x\in U}\Delta(x)
\ +\sum_{x\in \mathcal S\setminus U}\Delta(x)\min\{n^\star_x,N_x\},
\]
and consequently, whenever $R_T(\pi)>0$,
\begin{equation}
\label{eq:general-price}
\frac{R_T(\Cold)}{R_T(\pi)}\ \le\ P_U\ +\
\frac{\kappa^{\Cold}_U\,C^{\mathrm{free}}_U(t_0)+\sum_{x\in U}\Delta(x)+\sum_{x\in \mathcal S\setminus U}\Delta(x)\min\{n^\star_x,N_x\}}{R_T(\pi)} .
\end{equation}
For $\pi=\Hyb_\tau$ and $t_0=\tau$ this bounds the multiplier of the handoff
hybrid by the price ratio between cold exploration and the expert on the
arms the warm start certifies, plus a residual.
\end{theorem}

\emph{Proof sketch.} The proof (Appendix~\ref{app:proof-general-price}) has
two steps. First, cold LinUCB pays its own price, once per arm: on its
confidence event it plays a suboptimal arm $x$ only while
$2\beta\|x\|_{V_t^{-1}}\ge\Delta(x)$, and its own pulls of $x$ shrink that
width, so it pulls $x$ at most $\min\{n^\star_x,N_x\}$ times
(Lemma~\ref{lem:cold-price}). Second, certification costs at least the price:
by Cauchy--Schwarz, $\pi$ certifies $x$ only once its coverage of $u_x$
exceeds $C^\star_x-\lambda$, and it buys that coverage at a price of at least
$\kappa^\pi_U(t_0)$ per unit (Lemma~\ref{lem:cert-cost}). Splitting cold
LinUCB's bill into certified and uncertified arms and paying for the
certified ones at the rate $\kappa^{\Cold}_U$ gives the first display; since
$R_T(\pi)\ge R_{[1,t_0]}(\pi)$, dividing by $R_T(\pi)$ gives the second.

Only the confidence event of cold LinUCB is used, and
$\pi$ is arbitrary, so the bound covers the fixed-$\tau$ hybrid, the adaptive
router (take $t_0$ to be its last routed round), a CORRAL-style master, and
cold LinUCB itself. It bounds the ratio of the two policies' regrets on the
same instance, not one upper bound against another: two policies with the
same asymptotic rate can still differ by a diverging factor, so a statement
about the multiplier must bound one policy's regret in terms of the other's,
as the theorem does. Appendix~\ref{app:price-template} evaluates the bound on the
instance of Theorem~\ref{thm:price}.

\paragraph{The two parts of the bound.}
The price ratio covers the arms the policy paid to certify. The
residual is what nobody paid for, and it has three pieces: the coverage the policy obtains for free by playing optimal arms
(zero when the optimal arm is orthogonal to the certified directions, as on
every instance of the catalog), one pull of each certified arm,
and cold LinUCB's regret on the arms whose widths the warm start does not push
below their gaps. Only the last can be large, and it is counted arm by arm,
so it can be loose when many arms share information. The theorem therefore
controls the \emph{certified saving}, the part of cold LinUCB's regret on arms
the policy paid to certify, rather than the multiplier outright. On those arms the hybrid faces
the same confidence bound as cold LinUCB, so no coverage argument can say
which of the two does better there, and the term cannot be dropped
(Corollary~\ref{cor:diverge}). When the arms cold LinUCB would
otherwise pay for are certified, the residual is small and the price ratio is
the bound, and Theorem~\ref{thm:price} shows that it is then attained up to
constants.
Whether the residual can be large without a large price ratio is open. The
bound holds whether the saving is collected during the routed rounds
($A_\tau$) or after them ($I_\tau$), and the confidence radius cancels in
$P_U$ because it scales $C^\star_x$, the coverage both policies must reach,
identically.

\begin{corollary}[Diverging online multipliers require vanishing prices]
\label{cor:diverge}
Consider a sequence of instances, policies $\pi$ and cutoffs $t_0$ along
which $R_T(\Cold)/R_T(\pi)\to\infty$ and, for some choice of
$U\subseteq U_{\mathrm{cert}}(t_0)$ on each instance, the numerator of the
residual in \eqref{eq:general-price} is at most a fixed multiple of $R_T(\pi)$.
Then $P_U\to\infty$: the price $\kappa^\pi_U(t_0)$ that $\pi$ pays for the certified
coverage vanishes relative to the price $\kappa^{\Cold}_U$ cold LinUCB pays for it, so
$\pi$ plays at least one action whose gap is negligible against the coverage it supplies.
For the handoff hybrid this means the expert has access to such an action.
On the canonical instance of Proposition~\ref{prop:three-way} no online expert
certifies any arm that is offered after the handoff, so all of cold LinUCB's
regret on those arms is in the residual; that is the price-theoretic content of
Proposition~\ref{prop:one-pass}. Offline, where the acquisition is not
charged, the same instance multiplies by a factor of the horizon.
\end{corollary}

The proof is in Appendix~\ref{app:proof-diverge}. Two relaxations of
Theorem~\ref{thm:general-price} and a discussion of when cheap informative
actions exist are in Appendix~\ref{app:extra-theory}.

\section{The Bound Is Attained, and Information Value Can Vanish}
\label{sec:catalog}

Theorem~\ref{thm:general-price} is attained up to constants, and the
information value it controls can vanish. On the canonical separating
instance an offline warm start multiplies by a factor of the horizon while the
data an online handoff inherits is worth nothing; on a modification of it the
price ratio is attained.

\paragraph{The canonical instance: offline against online.}
\label{sec:canonical}
The canonical instance, which we construct to model the phenomenon reported
by \citet{chen2024llmdecision}, uses
horizon $n$, a single reward direction $g=e_1$, and $n$ orthogonal ``bad''
arms $b_t=Le_{t+1}$ with $L>1$.

\begin{proposition}[Offline three-way separation]
\label{prop:three-way}
Fix $n\ge1$, $L>1$, $\lambda>0$. There exist $d\ge n+1$, $\theta^\star$, and
two-action sets $\{\mathcal A_t\}_{t=1}^n$ such that, for any sub-Gaussian
noise: (1) a constant policy $\pi_{\mathrm{const}}$ has pseudo-regret exactly
$n$; (2) cold LinUCB also incurs regret $n$, selecting the same action as
$\pi_{\mathrm{const}}$ each round; (3) for $m$ large enough, LinUCB
warm-started with the design matrix that $\pi_{\mathrm{const}}$ generates from
$mn$ \emph{offline} rollouts incurs $O(1)$ regret with probability $\ge1-\delta$.
\end{proposition}

Cold LinUCB fails because at each round the fresh arm $b_t$ has width
$L/\sqrt\lambda$ whatever was played before, while $g$ has width
$1/\sqrt\lambda$ and estimated value $0$, so optimism picks $b_t$ every time.
The offline design matrix $M=mL^2\sum_{j\ge2}e_je_j^\top$ does not teach
LinUCB that $g$ is good; it teaches it that the \emph{alternatives} are not
uncertain, since $\|b_t\|^2_{(\lambda I+M)^{-1}}=L^2/(\lambda+mL^2)$ falls
below the gap once $m$ is moderate. Acquiring that data is not charged.
Online it would be, and the instance is hostile to online acquisition: every
round introduces a direction that no earlier observation touches.

\begin{proposition}[No online gain on the canonical instance]
\label{prop:one-pass}
Consider the instance of Proposition~\ref{prop:three-way} with horizon $T$
and confidence radius $\beta$, and let $\zeta:=\sigma\sqrt{2\log(2T/\delta)}/\sqrt\lambda$.
If $\beta(L-1)/\sqrt\lambda>1+\zeta$, then with probability at least
$1-\delta$, for \emph{every} expert policy $\Exp$ and every $\tau\in\{0,\dots,T\}$,
the handoff hybrid $\Hyb_\tau$ plays $b_t$ at every round $t>\tau$. Consequently
$R_{[\tau+1,T]}(\Hyb_\tau)=R_{[\tau+1,T]}(\Cold)=T-\tau$ and $I_\tau=0$.
\end{proposition}

No online router realizes the offline separation on this instance: a bad
expert's data has $I_\tau=0$, and a good expert has $B_\tau\le0$ because it is
already optimal (Corollary~\ref{cor:cat-oracle},
Appendix~\ref{app:catalog-statements}). Online multiplication needs an expert
whose routed actions are cheap \emph{and} cover directions that recur.

\paragraph{The price ratio is attained.}
Theorem~\ref{thm:general-price} is tight up to constants. We modify the
canonical instance so that directions recur and, for a while, a cheap
informative action is available: a \emph{probe} of gap $\eps$ that covers a
bad direction at rate $\gamma^2$ per round. The hybrid that routes the probes
and then hands off beats cold LinUCB by a factor between
$\gamma^2/(2220\,\eps L^2)$ and $1+\gamma^2/(\eps L^2)$
(Theorem~\ref{thm:price}, Appendix~\ref{app:template-thm}), the ratio of cold
exploration's price $1/L^2$ to the probe's price $\eps/\gamma^2$. A catalog of
eleven instances built the same way (Table~\ref{tab:catalog},
Appendix~\ref{app:catalog-statements}) shows that removing any one of
recurrence, cheapness, correct targeting, or the expert's suboptimality on
the tail can remove the gain: across these linear instances a gain over cold
LinUCB appears only when the expert buys the directions cold LinUCB would
otherwise pay for, more cheaply and early enough for them to recur.

\section{Coverage and the Transfer}
\label{sec:two-channels}

A warm start $(M,c)$ enters LinUCB through two objects: the design matrix
$V_1=\lambda I+M$, which fixes every confidence width, and the response sum
$c_1=c$, which fixes the estimate $\widehat\theta_1=V_1^{-1}c$.
Theorem~\ref{thm:general-price} bounds the value of a warm start through the
first. We now show that in the model of Section~\ref{sec:setup}, with an
expert that sees only what the learner sees, the first is the whole story,
and that the second matters only when the logged rewards depart from that
model, in which case it has an explicit form.

\begin{proposition}[The response sum is determined by the design matrix]
\label{prop:response-sum}
Let the warm start be built from pairs $(x_s,r_s)_{s\le\tau}$ with
$r_s=\langle\theta^\star,x_s\rangle+\eta_s$ and $\eta_s$ i.i.d.\
$\mathcal N(0,\sigma^2)$, independent of which actions were logged and of
which rounds were routed (as for a fixed handoff time and an expert whose
choices depend only on the features). Then
\[
c=M\theta^\star+\xi,\qquad \xi:=\sum_{s\le\tau}\eta_sx_s\sim\mathcal N(0,\sigma^2M),
\qquad
\widehat\theta_1-\theta^\star=-\lambda V_1^{-1}\theta^\star+V_1^{-1}\xi .
\]
Hence the law of $(V_1,c_1)$ given $M$ does not depend on which policy
generated the logged actions: two experts, or two routed sets, with the same
Gram matrix leave warm starts with the same law, and the LinUCB phases of the
corresponding hybrids have the same law; in particular they have the same
$\E[I_\tau]$. ($B_\tau$ can still differ, since it also involves the expert's
own actions after the handoff.)
\end{proposition}

For $\sigma$-sub-Gaussian noise the distributional identity is replaced by the
observation that the self-normalized bound of \citet{abbasi2011improved}, the
only property of $\xi$ used in this paper, depends on the features only
through $M$; every bound on the LinUCB phase in
Sections~\ref{sec:decomp}--\ref{sec:catalog} is therefore unchanged when one
expert is exchanged for another of equal coverage, provided neither the
experts' choices nor the routed rounds depend on the noise.

\begin{corollary}[Coverage is complete under linear rewards]
\label{cor:coverage-complete}
Under the model of Section~\ref{sec:setup} with i.i.d.\ Gaussian noise, and
with the logged actions and routed rounds independent of the noise, two
experts whose routed rounds leave the same Gram matrix give hybrids whose
LinUCB phases have the same law, and hence the same $\E[I_\tau]$. A systematic
difference in $I_\tau$ between experts of equal coverage therefore comes from
a departure from the linear model, from non-Gaussian noise, or from choices
(of the expert or of the router) that depend on the noise.
\end{corollary}

\begin{proof}
Immediate from Proposition~\ref{prop:response-sum}: equal Gram matrices give
warm starts $(V_1,c_1)$ of the same law, and the LinUCB phase is a
deterministic function of $(V_1,c_1)$ and of the future contexts and noise,
which are independent of the warm start. Cold LinUCB does not depend on the
expert, so $\E[I_\tau]=\E[R_{[\tau+1,T]}(\Cold)]-\E[R_{[\tau+1,T]}(\Hyb_\tau)]$
is the same for both.
\end{proof}

Corollary~\ref{cor:coverage-complete} is a prediction that can be checked on
data: experts whose routed rounds cover the same directions should leave
warm starts of the same value. Section~\ref{sec:experiments} checks it.

\paragraph{When the transfer appears.}
Suppose $r(x)=\langle\theta^\star,x\rangle+\varepsilon(x)+\eta$ for a fixed
reference $\theta^\star$---for instance the best linear predictor---and a
residual $\varepsilon$; the binary Precision@1 reward of
Section~\ref{sec:experiments} is of this form for every $\theta^\star$, and
badly so, since no linear function of $\phi(x,a)$ reproduces an indicator.
Regret is still measured against the true best action, with
$\bar r:=\langle\theta^\star,\cdot\rangle+\varepsilon(\cdot)$ and
$\Delta_{t,a}=\bar r(x_{t,a^\star_t})-\bar r(x_{t,a})$. A second departure concerns the expert
rather than the reward. Proposition~\ref{prop:response-sum} assumes that the
noise is independent of which actions were logged. An LLM reads the raw
context, which the features only summarize, so it can prefer an action on
exactly the rounds where that action pays; its logged rewards are then
selected, and the response sum is no longer determined by the Gram matrix,
even when $\varepsilon\equiv0$.

\begin{definition}[The transfer]
\label{prop:transfer}
Fix a reference $\theta^\star\in\R^d$ (under misspecification, the best linear
predictor). The \emph{transfer} of a warm start built from
$(x_s,r_s)_{s\le\tau}$ is
\[
m:=V_1^{-1}\sum_{s\le\tau}\bigl(r_s-\langle\theta^\star,x_s\rangle\bigr)\,x_s,
\qquad\text{so that}\qquad
\widehat\theta_1=\theta^\star-\lambda V_1^{-1}\theta^\star+m
\]
exactly, since $c=M\theta^\star+V_1m$ and $V_1^{-1}M=I-\lambda V_1^{-1}$.
Given $M$, the transfer is the only part of the estimate at handoff that
depends on which rounds were logged and what they returned.
\end{definition}

Under the linear model with Gaussian noise, and with the logged actions and
routed rounds independent of the noise, $m=V_1^{-1}\xi$ and its law is fixed
by $M$ (Proposition~\ref{prop:response-sum}). It departs from that law
through a residual $\varepsilon(x)=\E[r\mid x]-\langle\theta^\star,x\rangle$
(misspecification); through an expert whose choices depend on information
about the reward that the features do not carry, as an LLM reading the raw
text does, which makes $\E[r_s-\langle\theta^\star,x_s\rangle\mid x_s]\ne0$ on
the routed rounds even when $\varepsilon\equiv0$; and through routed rounds
chosen using past rewards. On AmazonCat the reward is a fixed function of the
context, so these sources cannot be told apart there, and we measure their sum.

The transfer can help or hurt: Proposition~\ref{prop:transfer-regret}
(Appendix~\ref{app:extra-theory}) bounds its effect on a round, and the bound improves when it moves the
estimate toward a reference whose residual is smaller on the best action than
on the actions LinUCB would otherwise play.

\paragraph{How far the transfer can go.}
Proposition~\ref{prop:response-sum} says that two experts of equal coverage
are interchangeable when the logged noise is independent of the logged
actions. The next result shows how far this fails once the expert sees more
than the learner and LinUCB explores less than the theory requires, as in our
experiments. It modifies the canonical instance of
Proposition~\ref{prop:three-way}: the good action is right only on a fraction
$p$ of the rounds, which rounds is visible to the expert but not to the
learner, a second action, right on a smaller fraction $q$, is added, and the
fresh arm is right on a fraction $f$.

\emph{The instance.} Fix $k\ge1$, $0\le q<p$ and $0\le f<p$ with $p+q+f\le1$,
$L>1$, $\lambda>0$, a horizon $T>4k$, and $d=T+2$. At round $t$ the learner is
offered $D=e_1$, $R=e_2$ and a fresh arm $b_t=Le_{t+2}$. Each round has a type
$y_t\in\{D,R,F,\varnothing\}$ that the learner does not see: on a round of
type $D$ or $R$ the action of that name has reward $1$, on a round of type $F$
the fresh arm has reward $1$, and every other action has reward $0$. The
first $\tau:=4k$ rounds form a window with $k$ rounds of type $D$, $k$ of type
$R$ and $2k$ of type $\varnothing$, in any order; later types are i.i.d.\ with
$\Prb(D)=p$, $\Prb(R)=q$, $\Prb(F)=f$. After the window the learner's expected
reward is linear in its features, $\mu(x)=\langle\theta^\star,x\rangle$ with
$\theta^\star=pe_1+qe_2+(f/L)\sum_{t}e_{t+2}$, and regret is measured against
it throughout: the gaps are $0$ for $D$, $p-q$ for $R$ and $p-f$ for every
$b_t$. Two experts see the type. $\Exp_A$ plays $D$ on rounds of types $D$,
$R$ and $F$ and plays $R$ on rounds of type $\varnothing$; $\Exp_B$ plays $R$
on rounds of types $D$ and $R$, the fresh arm on rounds of type $F$, and $D$
on rounds of type $\varnothing$. After the window $\Exp_A$ is right with
probability $p$ and $\Exp_B$ with probability $q+f$, so $\Exp_B$ is the more
accurate expert whenever $f>p-q$. LinUCB uses regularization $\lambda$ and a
constant radius $\beta>0$; we write $\omega:=\beta L/\sqrt\lambda$ for the
index of an untried fresh arm, and break ties in favour of the fresh arm.
The two experts are built so that everything coverage can see is the same:
each plays $D$ and $R$ exactly $2k$ times in the window, pays the same
regret, and is right exactly $k$ times. Their wrong answers are padding that
makes these totals equal, so the only difference left is which action their
right answers fall on.

\begin{theorem}[Equal coverage, unequal value]
\label{thm:transfer-separation}
On this instance the following hold.
\begin{enumerate}[itemsep=1pt,topsep=1pt,leftmargin=*]
\item[(i)] \emph{(Same coverage, same cost, same right answers when routed.)}
$M_\tau^{A}=M_\tau^{B}=2k(e_1e_1^\top+e_2e_2^\top)$, both acquisition regrets
equal $2k(p-q)$, and each expert is right on exactly $k$ of its $\tau$ routed
rounds. The two hybrids therefore hand off with the same design matrix, hence
the same widths, certified arms and prices, and their estimates differ
only through the transfer:
$\widehat\theta^{A}_{\tau+1}=\frac{k}{\lambda+2k}\,e_1$ and
$\widehat\theta^{B}_{\tau+1}=\frac{k}{\lambda+2k}\,e_2$.
\item[(ii)] \emph{(Cold, and coverage alone.)} Cold LinUCB plays the fresh arm
at every round, so $R_T(\Cold)=(p-f)T$. The same holds after round $\tau$ for
either hybrid if the routed rounds contribute their covariance but not their
rewards.
\item[(iii)] \emph{(The expert whose right answers are rare.)}
$\Hyb^{B}_\tau$ never plays $D$ after round $\tau$, so
$R_T(\Hyb^{B}_\tau)\ge(T-2k)\min\{p-q,\,p-f\}$.
\item[(iv)] \emph{(The expert whose right answers are frequent.)} Suppose the
exploration is small, $\omega<p$, and $h:=k(1-2\omega)-\omega\lambda>0$. With
probability at least $1-e^{-8(p-\omega)h}$, $\Hyb^{A}_\tau$ plays $D$ at every
round after $\tau$, and then $R_T(\Hyb^{A}_\tau)=2k(p-q)$. If also
$\omega\le\min\{p,\tfrac12\}/2$, this probability is at least
$1-e^{-2p(k-2\omega\lambda)}$.
\item[(v)] \emph{(In expectation.)}
$\E R_T(\Hyb^{A}_\tau)\le2k(p-q)+(p-f)T\,e^{-8(p-\omega)h}$. In particular, if
$k\ge\bigl(\ln T/(8(p-\omega))+\omega\lambda\bigr)/(1-2\omega)$, then
$R_T(\Cold)/\E R_T(\Hyb^{A}_\tau)\ge(p-f)T/\bigl(2k(p-q)+p-f\bigr)$.
\end{enumerate}
Consequently, with the probability in (iv),
\[
\frac{R_T(\Cold)}{R_T(\Hyb^{A}_\tau)}=\frac{(p-f)T}{2k(p-q)},
\qquad\text{while always}\qquad
\frac{R_T(\Cold)}{R_T(\Hyb^{B}_\tau)}\le\frac{p-f}{\min\{p-q,\,p-f\}}\cdot\frac{T}{T-2k}.
\]
\end{theorem}

Two experts that buy exactly the same coverage, at exactly the same price,
with exactly as many right answers on the routed rounds, give warm starts
whose multipliers over cold LinUCB differ, in most runs, by a factor of order
$T/k$; when $f>p-q$ the expert that loses is the more accurate one. Nothing
that depends on coverage alone---the design matrix, the widths, the
prices---can tell the two hybrids apart, and accuracy ranks the experts the
wrong way. What separates them is where their routed right answers fall:
$\Exp_A$'s are on an action that is right on a fraction $p$ of later rounds,
$\Exp_B$'s on one that is right on a fraction $q$. The learner settles all or
nothing: once $\Hyb^A_\tau$ leaves $D$ it never returns, so for a fixed $k$ it
fails for good in a fixed fraction of runs, at most $e^{-8(p-\omega)h}$, and
the average multiplier stays bounded. Part (v) shows that about $\log T$ right
answers make that fraction small enough for the separation to hold in
expectation as well, at order $T/\log T$.

That a learner which explores too little can lock onto whatever paid first
is known: greedy and near-greedy linear bandits can succeed or fail
depending on the data \citep{bastani2021mostly}, and the value of
demonstrations to an online learner depends on the demonstrator's quality
\citep{hao2023leveraging}. What Theorem~\ref{thm:transfer-separation} adds is
that, for an online warm start, lock-in makes the value of an expert's data
depend on where its right answers fall rather than on its coverage, its cost
or its accuracy: two experts equal in all three differ by a factor of order
the horizon.

\begin{remark}[Where the theorem sits]
\label{rem:transfer-regime}
The instance violates the hypotheses of both earlier results, and each
violation is needed. The experts see the type, so their logged rewards are not
independent of their actions, which Proposition~\ref{prop:response-sum}
assumes. And the radius is below any valid confidence radius, which
Theorem~\ref{thm:general-price} assumes. A valid radius has
$\beta\ge\sqrt\lambda\,\|\theta^\star\|_2\ge\sqrt\lambda\,p$, so
$\omega=\beta L/\sqrt\lambda>p$ because $L>1$, and the condition $\omega<p$ of
(iv) cannot be met: as $\Hyb^A_\tau$ keeps playing $D$, its index for $D$
tends to $p<\omega$, so after finitely many rounds it switches to fresh arms
and, its statistics for $D$ no longer changing, never returns. Both hybrids
then behave like cold LinUCB, and both multipliers tend to $1$ as the horizon
grows. With a valid radius coverage governs; with the reduced exploration used
in practice (Section~\ref{sec:exp-setup}) the transfer can decide. The fresh
arms serve only to make cold LinUCB fail and to make leaving $D$ permanent;
$d=T+2$ is what they cost.
\end{remark}

\begin{remark}[The instance as a cartoon of AmazonCat]
\label{rem:transfer-amazon}
Each piece of the instance has a counterpart in Section~\ref{sec:amazon}.
The action $D$ plays the role of the category present on $29.8\%$ of
AmazonCat contexts, and $R$ that of a rarer category. The fresh arms stand for
the thousands of categories cold LinUCB keeps trying because their widths
are large, which is why it settles on the dominant category in only $1$ of
$20$ seeds. The number $k$ of right answers on $D$ corresponds to the number
of rewarded routed answers on the dominant category: about $0.2$ per run for
Flan-T5-small and $3.7$ and $4.0$ for Flan-T5-base and Flan-T5-large. The
all-or-nothing settling of part~(iv) and its steep dependence on $k$ appear
in Tables~\ref{tab:settle} and~\ref{tab:dose}. The instance is not fitted to
the data; it isolates the mechanism the data display.
\end{remark}

Because the transfer is not a function of coverage, an estimate of the
multiplier built from coverage alone cannot tell the experts apart. A dry run
on labeled held-out data ranks them correctly, because it replays the expert's
right answers and so includes the transfer, but it underestimates the gain
about threefold and needs labeled data and the expert's answers on it
(Appendix~\ref{app:diagnostics}).

\section{Experiments}
\label{sec:experiments}

Corollary~\ref{cor:coverage-complete} predicts that experts of equal coverage
leave warm starts of equal value unless the transfer intervenes, and
Theorem~\ref{thm:transfer-separation} predicts how it can intervene. We test
both on real data, running the router, the LLM alone, and cold LinUCB on the
same stream; the setup mirrors \citet{chen2024llmdecision}.

\textbf{Two scales.} The literature reports \emph{reward multiplication},
which with binary Precision@1 is a ratio of accuracies; we call it the
\emph{accuracy multiplier}. The theorems bound a ratio of regrets, and since
every context has a correct action, $R_T(\pi)=T(1-\mathrm{acc}(\pi))$; we
therefore also report the \emph{regret ratio}
$(1-\max\{\mathrm{acc}_{\mathrm{LLM}},\mathrm{acc}_{\mathrm{Cold}}\})/(1-\mathrm{acc}_{\mathrm{Router}})$.
At low accuracies the two differ sharply: on AmazonCat an accuracy multiplier
of $3.29$ from a base of $8.4\%$ is a regret ratio of $1.27$, or $21\%$ less
regret. $A_\tau$, $I_\tau$ and $B_\tau$ are counted in correct answers and are
the same on either scale.

\subsection{Setup}
\label{sec:exp-setup}

\textbf{Datasets.} \textbf{OneShotWikiLinks-311}\footnote{\url{https://www.kaggle.com/generall/oneshotwikilinks}} ($T\!=\!622{,}000$,
$|\mathcal A|\!=\!311$) pairs two model families (Flan-T5, Qwen2.5) with the
same dataset, giving six configurations. \textbf{AmazonCat-13K}
\citep{bhatia2016extreme} ($T\!=\!50{,}000$, $|\mathcal A|\!=\!13{,}330$) pairs two embedding spaces
(BERT-768d, BGE-1024d) with the same categories. Reward is binary
Precision@1.

\textbf{Policies.} Each run tracks three policies on the same stream: the
\textbf{Router} (Algorithm~\ref{alg:confidence-gated},
Appendix~\ref{app:meta-algorithm}), which routes a round to the LLM when the
LLM's width $\|\phi(x_t,a_t^{(1)})\|_2/\sqrt{c_{\mathrm{LLM}}}$ is below LinUCB's width
$\|\phi(x_t,a_t^{(0)})\|_{V_t^{-1}}$ and updates $(V_t,c_t)$ on every round
with the played action and observed reward, so that its LinUCB component is
the online-warm-started learner of Section~\ref{sec:setup}; the \textbf{LLM}
alone, which proposes the action whose embedding is closest to its textual
output \citep{chen2024llmdecision}; and \textbf{cold LinUCB}. We use
$\lambda=0.1$ and the exploration scale $\alpha_t=\max(0.1,0.9985^t)$, which
reaches its floor at $t\approx1{,}500$. This is smaller than the confidence
radius the theory requires, so the confidence event of
Section~\ref{sec:width} is not guaranteed in our runs: we use the theory to
explain the experiments, not to certify them.
\todo{Phase B: add the tuned cold-LinUCB baselines (exploration schedule and
floor, $\lambda$, category-term features) and the Router re-tuned against them.}
A seed fixes the order of the
stream and is shared by all policies, so every comparison is paired. The
confidence parameter $c_{\mathrm{LLM}}$ of each configuration is chosen on seeds $11$--$20$
by mean Router accuracy and evaluated on seeds $1$--$10$, so tuning and
evaluation never share a seed; this selects $c_{\mathrm{LLM}}=100$ for AmazonCat/BERT,
$c_{\mathrm{LLM}}=0.1$ for AmazonCat/BGE, $c_{\mathrm{LLM}}=10^4$ for WikiLinks/BERT and
$c_{\mathrm{LLM}}=10$ for WikiLinks/BGE (Appendix~\ref{app:more-exp}).

\subsection{AmazonCat-13K}
\label{sec:amazon}

AmazonCat-13K has $13{,}330$ categories and $T=50{,}000$ rounds. We pair
BERT-768d features with Flan-T5 experts and BGE-1024d features with Qwen2.5
experts (Table~\ref{tab:amazon-results}).

\begin{table}[h]
\centering
\caption{AmazonCat-13K, $T=50{,}000$, mean $\pm$ standard deviation over
seeds $1$--$10$. Acc.\ mult.\ is the seed-averaged Router accuracy over the
seed-averaged $\max\{$LLM, cold$\}$; the regret ratio is
$(1-\max\{\text{LLM},\text{cold}\})/(1-\text{Router})$, the quantity
Theorem~\ref{thm:general-price} bounds. Cold LinUCB does not depend on the
expert. BERT rows are at $c_{\mathrm{LLM}}=100$ and BGE rows at $c_{\mathrm{LLM}}=0.1$, each selected on
the disjoint seeds $11$--$20$ (Appendix~\ref{app:more-exp}). On BGE the
two larger experts gain significantly (paired $t=3.92$ and $2.59$ against
cold) and the smallest does not ($t=0.01$).}
\label{tab:amazon-results}
\small
\setlength{\tabcolsep}{3pt}
\begin{tabular}{@{}llcccccc@{}}
\toprule
Features & LLM & LLM acc.\ & Cold LinUCB & Router & LLM\% & Acc.\ mult.\ & Regret ratio \\
\midrule
\multirow{3}{*}{BERT-768d} & Flan-T5-small  & $5.20\%$ & $8.39\pm2.46\%$ & $14.37\pm6.12\%$ & $2.8\%$  & $1.71\times$ & $1.07$ \\
& Flan-T5-base  & $3.61\%$ & $8.39\pm2.46\%$ & $22.80\pm9.00\%$ & $3.1\%$  & $2.72\times$ & $1.19$ \\
& Flan-T5-large & $5.27\%$ & $8.39\pm2.46\%$ & $\mathbf{27.64\pm5.56\%}$ & $3.2\%$  & $\mathbf{3.29\times}$ & $\mathbf{1.27}$ \\
\midrule
\multirow{3}{*}{BGE-1024d} & Qwen-0.5B  & $21.38\%$ & $32.59\pm5.00\%$ & $32.62\pm6.33\%$ & $0.04\%$  & $1.00\times$ & $1.00$ \\
& Qwen-3B  & $22.39\%$ & $32.59\pm5.00\%$ & $\mathbf{40.20\pm4.45\%}$ & $0.05\%$  & $\mathbf{1.23\times}$ & $\mathbf{1.13}$ \\
& Qwen-7B  & $27.17\%$ & $32.59\pm5.00\%$ & $38.71\pm6.38\%$ & $0.05\%$  & $1.19\times$ & $1.10$ \\
\bottomrule
\end{tabular}
\end{table}

\begin{figure}[h]
\centering
\includegraphics[width=0.86\linewidth]{two_channel.png}
\caption{Coverage versus the transfer on AmazonCat/BERT-768d, ten seeds, standard
errors over seeds. All three versions of the learner route exactly the same
rounds, the ones the full Router chose, and play the expert's answer there;
they differ only in what a routed round adds to the response sum. The three
experts cover the same directions to within $5\%$. With the transfer present
(dark bars) they span $1.71\times$ to $3.29\times$; with it removed they fall
to $1.16\times$--$1.50\times$ and no pairwise difference survives.
\emph{Transfer removed} keeps the expert-independent part
$M\theta_{\mathrm{or}}$ of the response sum and \emph{no response sum}
discards it entirely; the two do not differ significantly on any expert.}
\label{fig:two-channel}
\end{figure}

\paragraph{Equal coverage, unequal value.}
Held-out accuracy orders the three Flan-T5 experts differently from their
multipliers (Table~\ref{tab:amazon-results}). The decomposition locates the difference:
$A_\tau$ is $47$, $27$ and $53$ correct answers and $I_\tau$ is $2{,}936$,
$7{,}178$ and $9{,}571$ (Table~\ref{tab:amazon-decomp},
Appendix~\ref{app:more-exp}), and $96$--$100\%$ of
each Router's advantage accrues after its last routed round. The experts
cover the directions cold LinUCB is trapped on to within $5\%$, at nearly the
same prices, so their price ratios ($1.73$, $1.59$, $1.58$; an explanatory
statistic measured after the run, Appendix~\ref{app:amazon-transfer}) are
close to one another, ordered opposite to the multipliers, and above every
observed regret ratio. In the linear model with Gaussian noise,
Corollary~\ref{cor:coverage-complete} would make experts of equal coverage
interchangeable; the real rewards are not linear, so we test directly whether
the transfer carries the separation, and Figure~\ref{fig:two-channel} shows
that it does. We replay each run with the routed rounds held fixed: all three
versions of the learner route exactly the rounds the full Router chose and
play the expert's answer there, so they have the same routed data and the same
coverage, and differ only in what those rounds add to the response sum.
Feeding the learner $\langle\theta_{\mathrm{or}},\phi\rangle$ instead of the
observed reward on the routed rounds, which deletes $m$ of
Definition~\ref{prop:transfer} with $\theta^\star=\theta_{\mathrm{or}}$, turns
every significant pairwise difference into a null result (large over small:
paired $t=5.18$ becomes $1.33$); dropping the response sum on the routed
rounds gives the same collapse; and a linear reward makes the experts
interchangeable, as Proposition~\ref{prop:response-sum} predicts. On the
regret scale coverage alone saves about $2$ to $3.5\%$ of cold LinUCB's
regret for every expert, while the full router saves $6.5\%$, $16\%$ and
$21\%$ (Appendix~\ref{app:more-exp}).

\paragraph{The dominant category.} On AmazonCat one category is present on
$29.8\%$ of contexts, so a policy that knew this in advance and always
predicted it would score $29.8\%$. An online learner does not know it, and
finding it is not easy: context-free bandits that try categories without
using the context reach $0.04\%$ (UCB1) and $2.31\%$ (optimistic greedy) over
the same $50{,}000$ rounds and $13{,}330$ categories, seeds $1$--$10$. The
Routers that multiply are the ones that discover this category. LinUCB
settles on one frequent category and predicts it on a large share of rounds,
and its accuracy then tracks that category's frequency: runs that settle on
the dominant category finish near $28\%$, while cold LinUCB settles on a
category present on $11\%$ of contexts, or on none, in $18$ of $20$ seeds.
Whether a Router settles on the dominant category is decided by a handful of
the expert's routed answers. Across $177$ Router runs on BERT features (every
$c_{\mathrm{LLM}}$ and CORRAL, excluding configurations that route most rounds), a run whose
expert gave no correct routed answer on this category reached it in $3$ of
$45$ runs, and a run that received five or more reached it in all $54$.

\paragraph{A controlled test.} Adding $15$ correct answers on this category
to Flan-T5-small ($+0.15$ points of held-out accuracy) doubled its multiplier
($1.33\times$ to $2.62\times$, paired $p=0.002$), while $15$ correct answers on
random categories gave no significant gain ($1.65\times$, $p=0.15$), and the
same $15$ answers with their rewards hidden from the learner, so that they add
coverage along the category but no transfer, gave none either ($1.48\times$,
$p=0.47$; against the same answers with their rewards, $p=0.003$). Only about
two of the $15$ changed answers reach the learner in a typical run, because
the Router consults the expert on few rounds, almost all of them early: two
rewarded answers on the dominant category are enough to double the
multiplier, and the same two answers with their rewards hidden do nothing.

\paragraph{On the evaluation stream.} The same test on the $50{,}000$-round
stream and seeds of Table~\ref{tab:amazon-results} gives the same result:
$75$ added answers on this category raise the multiplier from $1.71\times$ to
$2.56\times$ ($p=0.007$), and with their rewards hidden only to $1.84\times$
($p=0.61$; Table~\ref{tab:inject50k}). This is the pattern of
Theorem~\ref{thm:transfer-separation}. As in its part~(ii), the routed
covariance without the routed rewards buys nothing; every correct answer
carries a residual near $1$, but only a handful on a frequent action push the
estimate in a direction that keeps paying, settling is all or nothing, and
its probability rises steeply with the number of such answers.
Appendix~\ref{app:dominant} gives the details.

\paragraph{BGE features and other routers.}
With BGE-1024d features cold LinUCB is much stronger and settles on the
dominant category by itself in $7$ of $20$ seeds, so there is less to add; the
two larger Qwen2.5 experts still gain significantly while routing $0.05\%$ of
rounds, again almost entirely through $I_\tau$ (Appendix~\ref{app:bge}). The
decomposition and the price bound hold for any routed set, so they apply to
the randomized CORRAL master \citep{agarwal2017corralling} used by
\citet{chen2024llmdecision} as to the gate. Against CORRAL and offline warm
starts on AmazonCat/BERT, no alternative is significantly better on any
expert, and where one comes out ahead in the mean the margin costs $3$--$5\times$
as many expert answers (Appendix~\ref{app:routers}).

\subsection{OneShotWikiLinks-311}
\label{sec:wiki}

On OneShotWikiLinks-311 ($622{,}000$ rounds, $311$ entities) the most
frequent entity covers $0.36\%$ of contexts, so there is no category for a
warm start to steer the learner toward. With Flan-T5 experts the Router
consults the expert on $13$--$15\%$ of rounds and gains $0.00$, $0.58$ and
$0.85$ points over cold LinUCB (accuracy multipliers $1.00\times$,
$1.05\times$ and $0.93\times$; regret ratios within $0.01$ of one): it moves
toward the expert rather than multiplying, and stays below Flan-T5-large
itself. The routed data does reach useful directions ($I_\tau$ is $2{,}392$,
$2{,}859$ and $3{,}098$ correct answers and positive on every seed); for
Flan-T5-small it exactly offsets what the expert loses on the routed rounds
($A_\tau=-2{,}395$). With Qwen2.5 experts the expert is already better than
the learner ($29$--$39\%$ against $19\%$), as in row~3 of
Table~\ref{tab:catalog}: the tuning seeds select the $c_{\mathrm{LLM}}$ at which the
gate routes every round, so the Router is the expert ($1.00\times$,
$I_\tau=B_\tau=0$), and at the largest $c_{\mathrm{LLM}}$ at which the gate hands
back it falls below the two larger experts. Full results are in
Appendix~\ref{app:more-exp}.

\section{Conclusion}
\label{sec:conclusion}

Reward multiplication in LLM--bandit hybrids has an exact accounting, and the
data the expert leaves behind reaches the learner through coverage and
through the transfer. The transfer is what decides the gain. When
exploration is below the valid radius, as in practice, it can separate
experts of identical coverage, cost and number of routed right answers by a
factor of order the horizon, even when the losing expert is the more
accurate. On AmazonCat-13K it carries every significant difference between
experts and most of the regret the router saves, through a handful of right
answers on the most frequent category: with their rewards hidden from the
learner, the same answers buy nothing. Coverage cannot play this role. With a
valid radius the multiplier it can deliver is capped by a ratio of prices,
the cap is attained on explicit instances, and the information value it
controls vanishes on the canonical one; under linear rewards with Gaussian
noise, for an expert whose choices ignore the noise, coverage is all there
is; and on AmazonCat it saves about $2$ to $3.5\%$ of cold LinUCB's regret
for every expert. A warm-start expert cannot be chosen by its accuracy or by
its coverage: what decides its value is which of its answers are right.
Information value is not policy value.

\paragraph{Limitations.}
The price bound assumes linear rewards and a fixed arm set with the best arm
always available; real streams satisfy these only approximately, and the
price ratio we report on them is a post-hoc explanatory statistic. Its
residual, cold LinUCB's regret on arms the warm start does not certify, is
not governed by a price, and we have no lower bound showing that it can be
large without a large price ratio, so the theorem characterizes the
certified saving rather than the multiplier outright. It also counts cold
LinUCB's regret arm by arm, so with many arms that share directions, as on
AmazonCat, the residual is large and the bound is loose: it explains why
coverage cannot produce large gains under valid exploration rather than
predicting the size of a gain. The constructions use a
fixed radius, unoptimized constants and very large instances.
Theorem~\ref{thm:transfer-separation} needs exploration below a valid
confidence radius, as in our experiments, holds with high probability for a
fixed number of right answers and in expectation only when that number grows
like $\log T$, and compares two experts that both see the reward; how an informed expert
compares with one that sees only the features is shown only numerically
(Appendix~\ref{app:proof-transfer-sep}). We have no estimator that predicts
the size of the multiplier before the expert is run. The large AmazonCat
multipliers depend on a dominant category that the learner can be steered
to, and we have not tested other datasets with a skewed label distribution.
Only two real datasets are used, multi-expert behaviour is not stress-tested,
reward models suited to binary feedback may change the split between
coverage and the transfer, and Lemma~\ref{lem:width-limited} would need an
analogue for Thompson sampling.

% ALT strongly encourages an anonymous AI Disclosure section when
% generative AI was used substantively. Uncomment and edit:
% \section*{AI Disclosure}
% We used [tool] to [briefly explain how it was used].

% Hidden by the [anon] option.
\acks{We thank ...}

\bibliography{references}


\appendix

\section{Related Work}\label{sec:related}
This appendix expands on Section~\ref{sec:related-main}. We cover (i) linear contextual bandits and confidence geometry,
(ii) model selection and expert aggregation, (iii) warm starts and prior
information, (iv) the geometry of information gain in linear bandits, and
(v) large language models for sequential decision making.

Sequential decision making with contextual information has been studied
extensively in the contextual bandit literature, where the learner must
balance exploration and exploitation under partial feedback
\citep{langford2007epoch,agarwal2014taming,foster2021statistical,
lattimore2020bandit}. In the linear setting, algorithms such as LinUCB and
OFUL achieve regret bounds that depend explicitly on the geometry of the
accumulated design matrix through self-normalized confidence ellipsoids
\citep{chu2011contextual,abbasi2011improved}. Subsequent work has developed
statistically efficient contextual bandit algorithms under general function
approximation assumptions \citep{agarwal2012contextual,foster2018practical,
foster2020beyond,simchi2022bypassing,zhu2022contextual,rucker2023infinite,
zhang2024efficient}. A central insight of this line of work is that the
conditioning of the covariance matrix determines how quickly uncertainty
shrinks and therefore how rapidly regret decreases. Our work builds on this
perspective by showing that actions selected by a frozen predictor can
reshape this geometry early in learning and accelerate subsequent adaptation.

\paragraph{Model selection and expert aggregation in bandits.}
A closely related direction studies how to combine multiple candidate
policies or base algorithms online. Classical expert-learning methods provide
regret guarantees relative to the best fixed expert in hindsight
\citep{auer2002nonstochastic}. In the bandit setting, CORRAL
\citep{agarwal2017corralling} uses Online Mirror Descent with a log-barrier
regularizer to achieve regret competitive with the best base algorithm;
subsequent refinements tighten the rates and extend the framework to
contextual and stochastic settings
\citep{pacchiano2020model,marinov2021pareto,zhu2022pareto,zhu2022near,
dann2024data}. These approaches are inherently \emph{randomized} and target
robustness: the master performs nearly as well as the best base algorithm
regardless of which one that is. Our objective is different---we want a
\emph{deterministic} rule that exploits the complementarity between a good
early frozen policy and a strong asymptotic adaptive learner to
\emph{strictly outperform} both. Correspondingly, our analysis does not
bound a relative regret against the best base, but instead tracks how early
actions selected by one policy reshape the design matrix and accelerate the
statistical efficiency of another.

\paragraph{Warm starts and prior information in contextual bandits.}
Several works have investigated how prior information can be incorporated
into contextual bandit algorithms through initialization, offline datasets,
or pretrained predictors. Such approaches can substantially improve
short-horizon performance by reducing uncertainty before online interaction
begins \citep{bastani2021mostly,zhang2019warm}. Related directions in
transfer learning and structured bandits further show that leveraging
information across tasks or user populations can accelerate adaptation and
reduce exploration requirements \citep{lazaric2012transfer,gentile2014online}.
Meta-learning approaches similarly learn priors that enable faster
exploration in new environments \citep{sharaf2019meta}. Empirical work on
neural-linear bandit architectures also demonstrates that representation
pretraining can significantly improve early decision quality
\citep{riquelme2018deep}. Our analysis formalizes these intuitions in a
linear bandit framework: informative early actions contract the confidence
ellipsoid on specific competitors (Lemmas~\ref{lem:loewner}
and~\ref{lem:transfer}), and the resulting saving is a difference of
width-limited regrets rather than an improvement of the
determinant-ratio rate, which Proposition~\ref{thm:handoff} shows cannot by
itself exceed a constant factor.

\paragraph{Geometry of information gain in linear bandits.}
A central theme in the analysis of linear bandit algorithms is that learning
efficiency is governed by the geometry of the accumulated design matrix.
Early work on stochastic linear optimization under bandit feedback showed
that regret depends explicitly on the growth of the covariance matrix
associated with the chosen actions and the resulting confidence ellipsoid
\citep{dani2008stochastic}. Subsequent analyses of linearly parameterized
bandits further clarified how the eigenvalue structure of this matrix
determines parameter identifiability and sample complexity
\citep{rusmevichientong2010linearly}. Later work established sharp
self-normalized confidence bounds and finite-time regret guarantees
expressed directly in terms of the design matrix and its inverse
\citep{abbasi2011improved}, while related results connected spectral growth
of the covariance matrix to optimal arm identification in linear settings
\citep{soare2014best}. From an information-theoretic perspective, regret can
also be related to determinant growth and effective dimension of the
covariance operator \citep{srinivas2012information,valko2013finite}. The
price \eqref{eq:price} is the single-direction shadow of the allocation
programs behind asymptotic instance-dependent lower bounds, where regret is
minimized subject to coverage constraints on every direction
\citep{graves1997asymptotically,lattimore2017end}; our use of it is
non-asymptotic and comparative: two policies on the same instance pay
different prices for the same coverage. Prior side information has also been
analyzed as hints or predictable sequences in online learning
\citep{rakhlin2013online,cutkosky2022leveraging}. The comparison with
\citet{cutkosky2022leveraging} is the closest: there the learner receives a
hint about $\theta^\star$ and the regret improves as a function of the hint's
accuracy, with the standard rate recovered when the hint is poor. Our warm
start is not a hint about $\theta^\star$ but data in the design matrix, and
its value is not a function of accuracy: two experts of similar accuracy can
leave design matrices of very different value (Section~\ref{sec:amazon}),
and an expert of zero accuracy can leave one of enormous value
(Proposition~\ref{prop:three-way}). What governs the value is the price ratio
of Theorem~\ref{thm:general-price}. The warm-start analyses of
\citet{bastani2021mostly,zhang2019warm} bound the regret of a learner started
from logged data through the quality of that data; they are offline in our
sense (row~1 of Table~\ref{tab:catalog}), and our contribution relative to
them is the online accounting, in which the acquisition of the data is
charged and can exceed its value. Our
work extends this geometric viewpoint by showing that a fixed policy can
shape the spectrum of the design matrix during early rounds of interaction,
thereby improving conditioning of the confidence ellipsoid and accelerating
the learning rate of subsequent adaptive exploration.

\paragraph{Large language models for sequential decision making.}
Recent work has explored the use of large language models (LLMs) as policies
in contextual bandit environments, demonstrating that frozen pretrained
models can achieve strong short-horizon performance even without online
updates \citep{chen2024llmdecision}. More broadly, LLMs exhibit strong
in-context learning capabilities \citep{brown2020language,xie2021explanation,
garg2022can,akyurek2022learning}, and instruction tuning further improves
their effectiveness in downstream decision tasks
\citep{ouyang2022training,chung2024scaling}. Several recent approaches
interpret sequential decision making itself as a sequence-modeling problem
and adapt transformer architectures to interactive settings
\citep{chen2021decision,zheng2022online,reid2022can,sun2023smart,
raparthy2023generalization,lee2023supervised}.
Closest to our work is the approach of \citet{chen2024llmdecision}, who
combine an LLM policy with an adaptive contextual bandit learner using an
online model-selection master and empirically document the
reward-multiplication phenomenon on large-scale recommendation datasets.
Our contribution relative to that line is theoretical: we give a
regret-based explanation for why such combinations can beat both components,
and replace the randomized CORRAL-style master with a deterministic
confidence-gated rule whose routing decision directly exposes the geometric
warm-start mechanism. A concurrent line of work \citep{huangcontextual} also
uses LLM uncertainty to gate exploration, but does so by calibrating an
embedding-based uncertainty estimate to adjust UCB-style exploration
bonuses, rather than by comparing self-normalized confidence radii of a
frozen predictor and an online learner; their analysis bounds the regret
of a single calibrated agent through the accuracy of the LLM's reward
predictions, whereas we give an exact decomposition of the hybrid's
advantage, bound the multiplier over the cold learner on any finite-arm
instance by a ratio of coverage prices, and show the bound is attained; in
their setting the LLM's value enters through prediction error, in ours
through the design matrix, and the two can disagree by a factor of the
horizon (Proposition~\ref{prop:three-way}).
%% NOTE: the next two citations need references.bib entries:
%%   jumpstart2026   -- "Jump Start or False Start? A Theoretical and Empirical
%%                      Evaluation of LLM-initialized Bandits", arXiv:2604.02527
%%   calibgated2026  -- "Calibration-Gated LLM Pseudo-Observations for Online
%%                      Contextual Bandits", arXiv:2604.14961
Two concurrent works are closer still. \citet{jumpstart2026} warm-start a linear
bandit from LLM-generated preferences and bound the warm-started estimator's
error through a single prior-error term measured in the geometry of the
pretraining design matrix, which treats the design matrix and the response
vector as one object; our analysis separates the two, shows that under linear
rewards the second is determined by the first, and identifies the
expert-dependent remainder in closed form, which is the part that differs
across equally accurate experts on our data. \citet{calibgated2026} inject LLM
pseudo-observations into LinUCB with a calibration-dependent weight on both the
design matrix and the response vector, a routing rule close in spirit to ours,
and analyze it through a bias--variance argument without a regret bound; the
decomposition of Proposition~\ref{prop:decomp} and the price bound of
Theorem~\ref{thm:general-price} apply to any such rule, since both hold for an
arbitrary routed set. Additional evidence that LLMs can
exhibit in-context reinforcement learning behavior
\citep{monea2024llms,krishnamurthy2024can} and that LLM-derived priors can
jump-start exploration \citep{alamdari2024jumpstarting,abbasiyadkori2024believe}
further motivates the study of principled mechanisms for combining frozen
predictors with online learners.


\section{Further Theory}
\label{app:extra-theory}

This appendix collects results used in the main text but stated there only in
passing: the two width identities, the regret of the handoff hybrid, the
width-limited bound with the realized estimation error, what the transfer can
buy, two relaxations of Theorem~\ref{thm:general-price}, and a discussion of
when cheap informative actions exist.

\subsection{How Widths Shrink}
\label{app:widths}

Lemma~\ref{lem:width-limited} is the familiar fact that an optimistic algorithm can only play an
arm whose width exceeds half its gap, restated so that the gap and the width
sit in one expression. It makes three facts visible at once. A competitor with
a large gap but an even larger width is exactly as dangerous as its gap
(Proposition~\ref{prop:three-way}: $\Delta=1$, width $\beta L/\sqrt\lambda>1/2$).
A competitor with a tiny gap is cheap no matter how wide it is, so a heavy
mass of near-ties does not by itself mean large regret. And the only way to
reduce the regret LinUCB can incur on a given round is to reduce widths
\emph{on competitors whose widths exceed their gaps}; reducing the width of
$g$ in Proposition~\ref{prop:three-way} would have achieved nothing.

There are exactly two ways widths get reduced: through data that arrives as a
warm start, and through LinUCB's own observations. Both are exact linear
algebra.

\begin{lemma}[A warm start shrinks widths in the directions it covers]
\label{lem:loewner}
For any $M\succeq0$ and any $x$, $\|x\|_{(\lambda I+M)^{-1}}\le\|x\|_{(\lambda I)^{-1}}=\|x\|_2/\sqrt\lambda$.
If $P$ is an orthogonal projector, $M\succeq mP$, and $x\in\mathrm{Im}(P)$, then
$\|x\|^2_{(\lambda I+M)^{-1}}\le\|x\|_2^2/(\lambda+m)$. If instead $x\perp\mathrm{Im}(M)$,
then $\|x\|_{(\lambda I+M)^{-1}}=\|x\|_2/\sqrt\lambda$: a warm start does nothing
for directions it does not cover.
\end{lemma}

\begin{lemma}[One observation shrinks the width of another direction by an
exact factor]
\label{lem:transfer}
For $V\succ0$ and vectors $x,z$ with $z\ne0$,
\[
\|z\|^2_{(V+xx^\top)^{-1}}=\bigl(1-\rho_V(x,z)\bigr)\|z\|^2_{V^{-1}},
\qquad
\rho_V(x,z):=\frac{(z^\top V^{-1}x)^2}{(1+\|x\|^2_{V^{-1}})\,\|z\|^2_{V^{-1}}}\in[0,1).
\]
At a cold start, $\rho_{\lambda I}(x,z)=\cos^2(x,z)\cdot\|x\|_2^2/(\lambda+\|x\|_2^2)$.
\end{lemma}

Lemma~\ref{lem:transfer} is where ``feature coherence'' enters, and it shows
precisely how much of that notion survives. Observing $x$ reduces the width of
$z$ by the squared cosine \emph{in the $V^{-1}$ geometry}, scaled by how novel
$x$ itself is. The Euclidean cosine is the right quantity only at $V=\lambda I$;
as $V$ becomes anisotropic the relevant angle is measured in $V^{-1}$, and two
vectors with positive Euclidean inner product can have $\rho_V$ near zero (or
vice versa). In Proposition~\ref{prop:three-way} the Euclidean cosine between
any two bad arms is exactly zero, so $\rho=0$ and no observation ever helps a
future round; that is what makes the instance hard, not the ambient dimension
as such. A large ambient dimension is necessary for many mutually
$V^{-1}$-orthogonal directions to exist, but it is the lack of carry-over
between directions, not the dimension, that cold LinUCB pays for.

\paragraph{Dimension, coherence and gaps.}
Effective dimension, feature coherence and gap structure are all readings of
Lemma~\ref{lem:width-limited} together with Lemmas~\ref{lem:loewner}
and~\ref{lem:transfer}: a uniform gap $\Delta_{\min}$ caps $\Phi_t$ at
$\Delta_{\max}$ but saves nothing while widths exceed $\Delta_{\min}/2$; a small
rank $r$ of the span of played directions caps the elliptical potential
$\sum_t\|x_{t,A_t}\|^2_{V_t^{-1}}$ at $2r\log(1+L^2T/(r\lambda))$ and therefore
the number of rounds at which any width can exceed a fixed level; and positive
coherence lower-bounds $\rho_{\lambda I}$ at the start of learning only. None
of the three is a condition under which a warm start cannot help by more than
a constant factor, and we do not claim one here; what they control is how
many rounds of cold exploration are needed before $\Phi_t$ falls below the
level the warm start could have delivered immediately.


\subsection{Regret of the Handoff Hybrid}
\label{sec:handoff}

\begin{proposition}[Handoff regret in exact form]
\label{thm:handoff}
Fix $\tau$ and an expert $\Exp$. On the event $\mathcal E_\delta$ (probability
at least $1-\delta$), the handoff hybrid satisfies
\begin{equation}
\label{eq:handoff-bound}
R_T(\Hyb_\tau)\;\le\;
\underbrace{\sum_{t\le\tau}\Delta_{t,a_t^{\Exp}}}_{\text{acquisition regret}}
\;+\;
\min\Biggl\{\sum_{t=\tau+1}^{T}\Phi_t(V_t),\;\;
2\beta_T\max\Bigl\{1,\tfrac{L}{\sqrt\lambda}\Bigr\}\sqrt{2(T-\tau)\,\log\frac{\det V_{T+1}}{\det V_{\tau+1}}}\Biggr\},
\end{equation}
where $V_{\tau+1}=\lambda I+M_\tau$ and the widths inside $\Phi_t$ are those
of the hybrid's own trajectory. The determinant ratio satisfies
$\log\det V_{T+1}/\det V_{\tau+1}\le\sum_{j=1}^{d}\log\bigl(1+L^2(T-\tau)/\sigma_j(V_{\tau+1})\bigr)$
with $\sigma_j$ the eigenvalues of $V_{\tau+1}$, and cold LinUCB satisfies the
same bound with $\tau=0$ and $V_1=\lambda I$.
\end{proposition}

We state the bound with the exact determinant ratio and make no asymptotic
simplification of it. Two remarks explain why.

\emph{The determinant term cannot certify a multiplier.} Suppose
$M_\tau\succeq\nu\tau P$ on an $r$-dimensional subspace containing all actions,
the situation analyzed in earlier work. Then the second term in
\eqref{eq:handoff-bound} is at most
$2\beta_T\sqrt{2r(T-\tau)\log(1+L^2(T-\tau)/(\lambda+\nu\tau))}$, and its ratio
to the cold bound $2\beta_T\sqrt{2rT\log(1+L^2T/\lambda)}$ is
\[
\Psi(\tau,\nu,T)=\sqrt{\frac{T-\tau}{T}}\cdot\sqrt{\frac{\log\bigl(1+L^2(T-\tau)/(\lambda+\nu\tau)\bigr)}{\log\bigl(1+L^2T/\lambda\bigr)}}.
\]
For $\tau=o(T)$ both logarithms are $\Theta(\log T)$ and $\Psi\to1$; for
$\tau=T^{\alpha}$ with $\nu\tau\gg\lambda$ it tends to $\sqrt{1-\alpha}$, bounded
away from zero. It is tempting to read $\Psi$ as
$\sqrt{\lambda/(\lambda+\nu\tau)}$ by comparing the arguments of the two
logarithms; that reading is wrong. The correct conclusion is the opposite
one: a warm start that improves only the $\sqrt{T}$-rate term buys at most a
constant factor, and any multiplier this bound can certify must come from the first term in the
minimum, that is, from competitors whose widths the warm start pushes below
their gaps so that $\Phi_t$ drops from the order of the gap to the order of
the (now small) width or to zero. That is what happens in
Proposition~\ref{prop:three-way}, where $\sum_t\Phi_t$ falls from $T$ to $O(1)$
while the determinant bound barely moves.

\emph{What the bound says about multiplication.} Combining
\eqref{eq:handoff-bound} with Proposition~\ref{prop:decomp}, the hybrid beats
cold LinUCB by more than a constant factor in this bound only if
\[
\sum_{t>\tau}\bigl(\Delta_{t,A^{\Cold}_t}-\Phi_t(V^{\Hyb}_t)\bigr)
\;\gg\;
\sum_{t\le\tau}\Delta_{t,a_t^{\Exp}}-R_{[1,\tau]}(\Cold),
\]
i.e., the width reduction that the routed data delivers on the competitors
cold LinUCB would have played must exceed the net regret paid to acquire it;
and it beats the expert only if the expert's tail regret exceeds the hybrid's.
Since $\Phi_t$ only upper-bounds the hybrid's regret, these are conditions for
the bound to certify a multiplier, not for a multiplier to occur.
Neither condition is implied by low expert accuracy being ``near-optimal
early'' or by an abstract coverage parameter. Both are properties of
\emph{which} actions the expert proposes, relative to the directions cold
exploration leaves unresolved, and of what those actions cost.

\emph{Adaptive $\tau$ and interleaving.} The router of
Appendix~\ref{app:meta-algorithm} decides at each round whether to follow the
expert or LinUCB by comparing confidence widths, so the expert rounds need not
be contiguous. Proposition~\ref{thm:handoff} applies verbatim with
$\sum_{t\le\tau}$ replaced by the sum over expert-routed rounds and with the
LinUCB-phase terms summed over LinUCB-routed rounds, since
Lemma~\ref{lem:width-limited} is a round-level statement and the elliptical
potential argument only uses that $V_t$ is non-decreasing.


\subsection{Width-Limited Regret with the Realized Estimation Error}
\label{app:two-objects}

Lemma~\ref{lem:width-limited} uses only the radius. With the realized
estimation error in its place it reads as follows, which shows how the two
objects of a warm start act on it.

\begin{proposition}[Width-limited regret with estimation error]
\label{prop:two-channel}
At a round $t$ at which LinUCB acts, let
$\rho_t:=\|\widehat\theta_t-\theta^\star\|_{V_t}$ be the realized estimation
error, so $\rho_t\le\beta_t$ on $\mathcal E_\delta$. Then on $\mathcal E_\delta$
\begin{equation}
\label{eq:phi-two}
\Delta_{t,A_t}\ \le\ \Phi_t(V_t,\rho_t)\ :=\
\max_{a\in\mathcal A_t}\min\bigl\{\Delta_{t,a},\ (\beta_t+\rho_t)\|x_{t,a}\|_{V_t^{-1}}\bigr\},
\end{equation}
and $\Phi_t(V,\rho)$ is non-increasing in $V$ and non-decreasing in $\rho$. A warm start acts
on the first argument through $M$ and on the second through $c$. Since
$0\le\rho_t\le\beta_t$ we have $\Phi_t(V_t,\rho_t)\ge\tfrac12\Phi_t(V_t)$: the
response sum can tighten the width-limited bound by at most a factor of two,
whereas the design matrix can reduce it by the price ratio, which is unbounded.
\end{proposition}

\subsection{What the Transfer Can Buy}
\label{app:transfer-regret}

Under the model of Section~\ref{sec:two-channels}, the transfer's effect on a
single round is bounded as follows.

\begin{proposition}[What the transfer can buy]
\label{prop:transfer-regret}
Under the misspecified model, at a round $t$ at which LinUCB acts and for any
reference $\widetilde\theta\in\R^d$, write
$\widetilde\rho_t:=\|\widehat\theta_t-\widetilde\theta\|_{V_t}$ and
$\widetilde\varepsilon(x):=\bar r(x)-\langle\widetilde\theta,x\rangle$. Then
\[
\Delta_{t,A_t}\ \le\ (\beta_t+\widetilde\rho_t)\|x_{t,A_t}\|_{V_t^{-1}}
\ +\ (\widetilde\rho_t-\beta_t)_+\|x_{t,a^\star_t}\|_{V_t^{-1}}
\ +\ \bigl[\widetilde\varepsilon(x_{t,a^\star_t})-\widetilde\varepsilon(x_{t,A_t})\bigr],
\]
simultaneously for every reference. The design matrix fixes the widths; the
response sum fixes $\widetilde\rho_t$ for every reference at once. The
bound improves when the transfer moves $\widehat\theta$ toward a reference whose
residual is smaller on the optimal arm than on the arms LinUCB would otherwise
play, and the bound is width-limited whenever some reference has
$\widetilde\rho_t\le\beta_t$ and a nonpositive residual difference on the
played arm.
\end{proposition}

Proposition~\ref{prop:transfer-regret} is a statement about what the transfer
\emph{can} deliver, not a prediction of what a given expert will deliver: $m$
depends on the residuals at the actions the expert chooses, which are not
known before it is run. Theorem~\ref{thm:transfer-separation} shows how large
the difference between two experts can be on a constructed instance, and
Section~\ref{sec:amazon} measures $m$ on real data, finds that it accounts for
essentially the whole estimate at handoff, and shows that removing it
collapses the differences between experts of similar accuracy.


\subsection{Two Relaxations of the Price Bound}
\label{app:relaxations}

\begin{remark}[Two relaxations]
\label{rem:relaxations}
(i) The factor $2$ in the definition of $U_{\mathrm{cert}}$ is the crude
one. On the confidence event with the true self-normalized radius
$\beta^{\mathrm{true}}_t\le\beta$, the proof of Lemma~\ref{lem:width-limited}
gives $\Delta_{t,A_t}\le(\beta_t+\beta^{\mathrm{true}}_t)\|x_{t,A_t}\|_{V_t^{-1}}$
for the played arm, so the theorem holds with $C^\star_x$ replaced by
$(\beta+\beta^{\mathrm{true}})^2\|x\|_2^2/\Delta(x)^2$ throughout and
$U_{\mathrm{cert}}$ enlarged to the arms with
$(\beta+\beta^{\mathrm{true}})\|x\|<\Delta(x)$; $P_U$ is unchanged. (ii) If the
optimal arm is not offered at every round, so that the gap of $x$ varies
with $t$, write $\Delta^-(x)$ and $\Delta^+(x)$ for its smallest and largest
values over rounds at which $x$ is offered. Lemma~\ref{lem:cold-price} holds
with $C^\star_x$ computed from $\Delta^-(x)$ and each pull charged $\Delta^+(x)$,
and Lemma~\ref{lem:cert-cost} with $\Delta^-(x)$; the price ratio then carries a factor $\Delta^+(x)/\Delta^-(x)$,
which is $1$ under the standing assumption.
\end{remark}

\subsection{When Do Cheap Probes Exist?}
\label{app:cheap-probes}

The positive rows need an action of small gap whose feature vector carries a
non-negligible component along a direction that cold LinUCB can otherwise
cover only by pulling an arm of large gap. Two observations make this less
special than it looks. First, the probe need not be a single designated arm:
any near-optimal action whose embedding overlaps the bad direction serves,
and in the embedding spaces used in practice (Section~\ref{sec:experiments})
near-optimal actions overlap many others. Second, the expert does not need to
know that the probe is informative; it needs only to prefer it, which any
policy with prior knowledge of the reward does when the probe is nearly as
good as the best arm. What the expert supplies, and cold optimism cannot, is
the \emph{order} in which coverage is bought: optimism ranks a fresh bad arm
above a cheap probe because the bad arm's width is larger, and pays the bad
arm's price first (Lemma~\ref{lem:cat-trap} in the appendix); a policy that
already believes the bad arm is bad takes the probe at once. The price ratio
$\gamma^2/(\eps L^2)$ of rows~5 and~$5'$ is the exchange rate between those two
orders, and Theorem~\ref{thm:general-price} says it is the only exchange rate
available: an expert whose actions are not cheaper per unit of certified
coverage than cold exploration's cannot multiply, however accurate it is.

\section{The Price Ratio Attained, and the Catalog}
\label{app:catalog-statements}

This appendix states the rows of Table~\ref{tab:catalog} that the main text
only summarizes. The canonical instance (rows 1--2) is defined in
Section~\ref{sec:catalog} and the template in Appendix~\ref{app:template-thm}
below; proofs are in Appendix~\ref{app:catalog}.

The catalog has eleven rows rather than the $2^3=8$ of a binary table
because several terms admit qualitatively distinct failure modes---for
instance, $I_\tau$ can vanish because the expert buys the same directions cold
exploration buys at the same price (row~4), because it buys the wrong
directions (row~7), or because it buys the right directions too dearly
(row~8)---because the offline/online distinction (rows~1--2) and the two
ways of killing the tail term (rows~3 and~6) each warrant separate entries,
and because the positive case appears twice, with the saving collected in
the routed window (row~5) and on the tail (row~$5'$).

\subsection{Standing Conventions}
\label{sec:template}

Noise is i.i.d.\ $\sigma$-sub-Gaussian and independent of the actions; we use
the reward-table model in which the $m$-th pull of a given action vector
receives the $m$-th entry of an i.i.d.\ sequence attached to that vector
\citep[\S4.6]{lattimore2020bandit}, which has the same law as the sequential
model and lets us compare two policies on the same noise. We fix
\begin{equation}
\label{eq:conventions}
L\ge8,\quad \gamma\in(0,1],\quad 0<\eps<\gamma/L,\quad
\lambda\ \ge\ \max\bigl\{\sigma^2d\,\ell_0,\ 8\sigma^2\ell_1,\ 60\bigr\},\quad
\beta\equiv2\sqrt\lambda,
\end{equation}
with $\ell_0:=\log(1+TL^2)+2\log(1/\delta)$ and $\ell_1:=\log(2dT/\delta)$.
The confidence radius is constant. The lower bound on $\lambda$ makes
$2\sqrt\lambda$ at least the radius of \citet{abbasi2011improved} at every
round, so the confidence event of Section~\ref{sec:width} holds for each
policy with probability $\ge1-\delta$, and it makes the noise on any $\lambda/30$
pulls of one arm smaller than a quarter of that arm's confidence bonus. The
constant $m_0:=\lfloor\lambda/30\rfloor\ge2$ is the number of pulls of a
fresh bad arm that optimism forces before the good arm or the probe can be
preferred to it. Conditions on $T$ appear in each statement; they only say
that the horizon is long enough for the claimed rounds to take place.


\subsection{The Template, and the Price Ratio Attained}
\label{app:template-thm}

All instances below are built from the canonical one by three modifications:
directions may recur, a cheap informative action may be available, and that
action may disappear. Fix $K\ge1$ and set $d=K+1$ (or $d=2K+1$ when nuisance
directions are present). The parameter is $\theta^\star=e_1$. Directions
cycle: $j(t):=2+((t-1)\bmod K)$. Three actions may be offered at round $t$:
\[
g:=e_1\ (\mu=1),\qquad
b_t:=Le_{j(t)}\ (\mu=0),\qquad
p_t:=(1-\eps)e_1+\gamma\,e_{j(t)}\ (\mu=1-\eps),
\]
the \emph{good} arm, the \emph{bad} arm, and the \emph{probe}. The instance
$\mathcal I(K,L,\gamma,\eps,\tau_0)$ has $\mathcal A_t=\{g,b_t,p_t\}$ for $t\le\tau_0$
and $\mathcal A_t=\{g,b_t\}$ for $t>\tau_0$. The price of $e_j$, $j\ge2$, is
$\min\{1/L^2,\eps/\gamma^2\}$ while probes are offered and $1/L^2$ afterwards. The
expert is the frozen linear policy $\Exp_w$ that plays
$\arg\max_{a\in\mathcal A_t}\langle\widehat\theta_{\Exp},x_{t,a}\rangle$ with
$\widehat\theta_{\Exp}=e_1+w\sum_{j\ge2}e_j$; ties do not occur for the values
of $w$ we use. Throughout, $\tau_0=Kn$ for an integer $n$, so each direction is
probed exactly $n$ times by a policy that plays every probe.

We work under the standing conventions \eqref{eq:conventions} of
Appendix~\ref{sec:template}: $L\ge8$, $\gamma\in(0,1]$, $0<\eps<\gamma/L$,
$\lambda\ge\max\{\sigma^2d\,\ell_0,8\sigma^2\ell_1,60\}$ with $\ell_0,\ell_1$
logarithmic in $T/\delta$, and the constant radius $\beta=2\sqrt\lambda$, in
the reward-table noise model of Remark~\ref{rem:decomp}; we write
$m_0:=\lfloor\lambda/30\rfloor$ for the number of pulls of a fresh bad arm
that optimism forces.

On $\mathcal I(K,L,\gamma,\eps,\tau_0)$ the expert $\Exp_w$ with
$w\in(1/L,(1-\eps)/(L-\gamma))$ plays every probe while probes are offered and
every bad arm afterwards (Lemma~\ref{lem:expert-routes}). Its routed actions
cost $\eps$ each, cover $e_{j(t)}$ at rate $\gamma^2$ per round, and recur; after
$\tau_0$ it is the worst possible policy. The hybrid that hands off exactly
when the probes disappear therefore isolates the price ratio.

\begin{theorem}[Online multiplication, and the price ratio as the multiplier]
\label{thm:price}
Let $\mathcal I=\mathcal I(K,L,\gamma,\eps,Kn)$ under \eqref{eq:conventions}, with
$\Exp=\Exp_w$ for some $w\in(1/L,(1-\eps)/(L-\gamma))$, and $\tau=\tau_0=Kn$. With
probability at least $1-3\delta$ the following hold simultaneously.
\begin{enumerate}[itemsep=1pt,topsep=1pt,leftmargin=*]
\item[(i)] \emph{(Expert.)} $R_T(\Exp)=\eps Kn+(T-Kn)$.
\item[(ii)] \emph{(Cold, lower bound.)} If $n\ge m_0$ then
$R_{[1,Kn]}(\Cold)\ge Km_0$; in particular $A_{\tau_0}\ge Km_0-\eps Kn$.
\item[(iii)] \emph{(Cold, upper bound.)} For every $n\ge1$, $R_T(\Cold)\le K(8\beta^2+1)+\eps Kn$.
\item[(iv)] \emph{(Hybrid.)} If $n\ge36\lambda L^2/\gamma^2$ and $K\ge36L^2$, then
$\Hyb_{\tau_0}$ plays $g$ at every round $t>\tau_0$, so
$R_T(\Hyb_{\tau_0})=\eps Kn$, $I_{\tau_0}=R_{[Kn+1,T]}(\Cold)\ge0$, and
$B_{\tau_0}=T-Kn$.
\end{enumerate}
Consequently, for $n\ge36\lambda L^2/\gamma^2$ and $K\ge36L^2$,
\begin{equation}
\label{eq:price-ratio}
\frac{m_0}{\eps n}\ \le\ \frac{R_T(\Cold)}{R_T(\Hyb_{\tau_0})}\ \le\ 1+\frac{8\beta^2+1}{\eps n},
\qquad
\frac{R_T(\Exp)}{R_T(\Hyb_{\tau_0})}=1+\frac{T-Kn}{\eps Kn},
\end{equation}
and at the smallest admissible $n=\lceil36\lambda L^2/\gamma^2\rceil$ the first
ratio is $\Theta(\gamma^2/(\eps L^2))$, with explicit constants
$\gamma^2/(2220\,\eps L^2)\le R_T(\Cold)/R_T(\Hyb_{\tau_0})\le1+\gamma^2/(\eps L^2)$
(the constants of this section are not optimized; see
Appendix~\ref{app:catalog}, ``Remarks on the constants'').
\end{theorem}

The bracket \eqref{eq:price-ratio} is the price ratio: cold exploration buys
coverage $L^2$ of a bad direction per unit of regret (price $1/L^2$), the
probe buys $\gamma^2$ per $\eps$ (price $\eps/\gamma^2$), and the multiplier over cold
LinUCB is their ratio $\gamma^2/(\eps L^2)$ up to constants. The conditions on
$n$ and $K$ in (iv) are not about prices alone. Enough coverage of $e_j$ is
necessary for the handoff matrix to certify $b_j$ but not sufficient: every
probe is orthogonal to $v=\gamma\,e_1-(1-\eps)\sum_{j\ge2}e_j$, so the width of
$e_j$ stays at least $(1-\eps)/\sqrt{\lambda(\gamma^2+K(1-\eps)^2)}$ however many
probes are played, and certification needs $K$ of order $L^2$; the price
ratio does not see this geometric condition. The instances are large by
design: orthogonal bad directions stand in for an action space too large for
cold exploration to exhaust (AmazonCat has $13{,}330$ categories), and at the
smallest admissible sizes the window alone exceeds $3\times10^8$ rounds. The multiplier over
the expert, $(T-\tau_0)/(\eps\tau_0)$, diverges as the tail lengthens. Here
the saving is collected during the routed rounds, as $A_{\tau_0}$; keeping
cold LinUCB busy elsewhere during the window moves it after them, into
$I_{\tau_0}$, with the same price ratio (Proposition~\ref{prop:cat-tail-saving},
row~$5'$, below).

\paragraph{The catalog.}
Theorem~\ref{thm:price} is one row of a catalog of eleven instances, each
obtained from the canonical one by changing a single ingredient
(Table~\ref{tab:catalog}). Exactly two online rows multiply over both base policies,
rows~5 and~$5'$, and they realize the same price ratio, one during the routed
rounds and one after them.

\begin{table}[t]
\centering
\scriptsize
\setlength{\tabcolsep}{2.5pt}
\resizebox{\linewidth}{!}{%
\begin{tabular}{@{}clccccl@{}}
\toprule
\# & Instance & $A_\tau$ & $I_\tau$ & $B_\tau$ & Ratio to $\Cold$ & Ingredient isolated\\
\midrule
1 & canonical, offline (Prop.~\ref{prop:three-way}) & not charged & $I^{\mathrm{off}}\asymp T$ & $B^{\mathrm{off}}\asymp T$ & $\asymp T$ & data is free\\
2 & canonical, online (Prop.~\ref{prop:one-pass}) & $\le\tau$ & $0$ & $\le0$ & $\le T/(T-\tau)$ & directions never recur\\
3 & canonical, oracle (Cor.~\ref{cor:cat-oracle}) & $\tau$ & $0$ & $-(T-\tau)$ & additive & expert too good\\
4 & recurring, bad-arm expert (Prop.~\ref{prop:cat-mimic}) & $0$ & $0$ & $\ge0$ & $\le961$ & same price as cold\\
5 & recurring, cheap probes (Thm.~\ref{thm:price}) & $\ge Km_0-\eps\tau_0$ & $\ge0$ & $T-\tau_0$ & $\asymp \gamma^2/(\eps L^2)$ & price ratio, in the window\\
$5'$ & as 5, cold kept busy in the window (Prop.~\ref{prop:cat-tail-saving}) & $-\eps\tau_0$ & $\ge Km_0$ & $T-\tau_0$ & $\asymp \gamma^2/(\eps L^2)$ & price ratio, on the tail\\
6 & as 5, expert optimal on tail (Prop.~\ref{prop:cat-tail}) & as 5 & as 5 & $0$ & as 5 & expert optimal late\\
7 & nuisance probes (Prop.~\ref{prop:cat-wrong}) & & $<R_{[\tau_0+1,T]}(\Cold)$ & & $\le65$ & wrong directions\\
8 & as 5 with $\eps\ge \gamma^2/L^2$ (Cor.~\ref{cor:cat-dear}) & & & & $\le2$ & probe too dear\\
9 & cold already cheap (Prop.~\ref{prop:cat-cheap}) & & & & $o(1)$ avg.\ reward & nothing to buy\\
10 & bad arms with gap $\eta$ (Prop.~\ref{prop:cat-gap}) & & & & $\le1+\gamma^2/(\eps L^2)$ & gap rescales prices\\
\bottomrule
\end{tabular}}
\caption{The catalog: what the three terms of Proposition~\ref{prop:decomp}
can be (statements in Appendix~\ref{app:catalog-statements}, proofs in
Appendix~\ref{app:catalog}). Row 1 is the offline identity, in which $A$ is
not charged; rows 2--10 are online handoffs. Rows 3 and 6 kill $B_\tau$ from
the two sides. Rows 4, 7 and 8 kill the multiplier over cold LinUCB: the
expert buys what cold buys at the same price, buys directions nobody needs,
or buys the right directions too dearly. Row 9 is the case in which cold
LinUCB's bill is already negligible, and row 10 shows that the gap rescales
both prices. Rows 5 and $5'$ realize the same price ratio, collected as
$A_\tau$ and as $I_\tau$. Blank cells are not claimed; constants are not
optimized.}
\label{tab:catalog}
\end{table}

Across these linear instances, apart from the direct advantage of an expert
that is simply better than cold LinUCB while it plays (row~3), a gain over
cold LinUCB appears only when the expert buys the directions cold LinUCB would
otherwise pay for, more cheaply and early enough for them to recur, and a
gain over the expert only when, in addition, the expert is worse than the
warm-started learner on the tail.

\subsection{The Other Rows}
\label{app:catalog-rows}

\paragraph{Row 3: an oracle expert on the canonical instance.}
\begin{corollary}[Oracle expert on the canonical instance]
\label{cor:cat-oracle}
On the instance of Proposition~\ref{prop:one-pass}, under its hypothesis, let
$\Exp$ play $g$ at every round. With probability at least $1-\delta$, for every
$\tau$: $R_T(\Exp)=0$, $R_T(\Hyb_\tau)=T-\tau$, $A_\tau=\tau$, $I_\tau=0$,
$B_\tau=-(T-\tau)$. The hybrid improves on cold LinUCB by exactly the rounds
it routes and on the expert by nothing.
\end{corollary}

\paragraph{Row 4: an expert that mimics cold exploration.}
\begin{proposition}[Expert mimics cold exploration: same price, no gain]
\label{prop:cat-mimic}
Let $\mathcal I(K,L,\cdot,\cdot,0)$ be the instance with no probes
($\mathcal A_t=\{g,b_t\}$ for all $t$, $d=K+1$), and let $\Exp$ be any policy
that plays $b_t$ at every round (for instance $\Exp_w$ with $w>1/L$). With
probability at least $1-2\delta$: (i) for every $\tau\le Km_0$,
$\Hyb_\tau$ and $\Cold$ play the same action at every round, so
$A_\tau=I_\tau=0$ and $R_T(\Hyb_\tau)=R_T(\Cold)$; (ii) for every
$\tau>Km_0$, $R_T(\Hyb_\tau)\ge\tau>Km_0$ while $R_T(\Cold)\le K(4\beta^2+1)$,
so $R_T(\Cold)/R_T(\Hyb_\tau)\le(16\lambda+1)/m_0\le961$.
\end{proposition}

\paragraph{Row 5: discussion of Theorem~\ref{thm:price}.}
The bracket \eqref{eq:price-ratio} is the price ratio. Cold exploration buys
coverage $L^2$ of a bad direction per unit of regret, so its price is $1/L^2$;
the probe buys $\gamma^2$ per $\eps$, price $\eps/\gamma^2$; the ratio of the two prices
is $\gamma^2/(\eps L^2)$, and the multiplier over cold LinUCB is that ratio up to
constants, and is $O(1)$ when the ratio is. The confidence radius does not
appear: both policies must push the same widths below the same level, and
$\beta$ scales the required coverage identically for both. The multiplier
over the expert is $(T-\tau_0)/(\eps\tau_0)$ and diverges as the tail
lengthens, because the expert keeps paying for bad arms that the warm-started
learner has learned to avoid. The lower bound in (ii) is attained inside the
routed window, so on this instance $A_{\tau_0}>0$ whenever $\eps n<m_0$: the
expert's probes are cheaper than the exploration cold LinUCB is forced into
at the same rounds, and the data they leave is what prevents that exploration
from resuming after $\tau_0$. On this instance the saving is carried by
$A_{\tau_0}$, not by $I_{\tau_0}$, and for a structural reason: the
certificate in (iv) needs coverage of order $\beta^2L^2$ of each bad
direction, which the probe supplies at $\gamma^2$ per round while cold LinUCB's
forced pulls supply it at $L^2$ per round. The window long enough for the
expert to certify ($n\ge36\lambda L^2/\gamma^2$ rounds per direction) is therefore
longer by a factor of order $L^2/\gamma^2$ than the exploration cold LinUCB is
forced into ($m_0\le$ pulls $\le 8\beta^2+1$ per direction), and cold LinUCB
completes that exploration inside the window whenever the bad arm is offered
there (exact-arithmetic simulations of the template, at the reduced sizes of
Appendix~\ref{app:synthetic}, give $R_{[Kn+1,T]}(\Cold)=0$). Theorem~\ref{thm:general-price}
does not distinguish the two: the certified saving is bounded by the
same price ratio wherever it is collected. Row~$5'$ shows that the same ratio
can be collected after the routed rounds when the routed window keeps cold
LinUCB away from the directions it will be trapped on later.

\begin{remark}[What certification costs on the template]
On the template, the certificate used in Theorem~\ref{thm:price}(iv), namely
that the confidence bound alone excludes the bad arm in direction $e_j$,
$2\beta L\|e_j\|_{V_{\tau+1}^{-1}}<1$, requires by
Lemma~\ref{lem:no-free-coverage} that
$R_{[1,\tau]}(\Exp)>\kappa(e_j)\,(4\beta^2L^2-\lambda)$ for that direction, hence
a total acquisition regret exceeding
$K\min\{1/L^2,\eps/\gamma^2\}(4\beta^2L^2-\lambda)$ for all $K$ directions. This is
Lemma~\ref{lem:cert-cost} written out for the template.
\end{remark}

\paragraph{Row $5'$: the same price ratio, collected on the tail.}
For the saving to appear as $I_\tau$, the routed window must contain rounds
at which cold LinUCB is \emph{not} offered the bad arms it will later be
trapped on, and is busy with something else; otherwise, as just explained, it
completes its forced exploration inside the window. The following instance
keeps cold LinUCB busy with \emph{distractors}: arms of gap zero whose
feature vectors carry a large component in a fresh direction, so that
optimism prefers them to everything else while the expert probes. Fix
$K,K',n\in\N$, put $d=1+K+K'$, $\theta^\star=e_1$, $L':=12$, and
$\tau_0:=Kn$. Directions $J:=\{2,\dots,K+1\}$ are the bad directions and
$J':=\{K+2,\dots,K+K'+1\}$ the distractor directions, cycled by
$j(t):=2+((t-1)\bmod K)$ and $j'(t):=K+2+((t-1)\bmod K')$. The arms are
\[
g=e_1\ (\mu=1),\quad p_j=(1-\eps)e_1+\gamma\,e_j\ (\mu=1-\eps),\quad
b_j=Le_j\ (\mu=0),\quad b'_{j'}=e_1+L'e_{j'}\ (\mu=1),
\]
and the instance $\mathcal J(K,K',L,\gamma,\eps,n)$ offers
$\mathcal A_t=\{g,p_{j(t)},b'_{j'(t)}\}$ for $t\le\tau_0$ and
$\mathcal A_t=\{g,b_{j(t)}\}$ for $t>\tau_0$. The expert plays the probe while
probes are offered and the bad arm afterwards; the linear policy
$\widehat\theta_{\Exp}=e_1+w\sum_{j\in J}e_j-\sum_{j'\in J'}e_{j'}$ with
$w\in(1/L,(1-\eps)/(L-\gamma))$ does this, by the computation of
Lemma~\ref{lem:expert-routes} together with $\langle\widehat\theta_{\Exp},b'_{j'}\rangle=1-L'<1$.
The prices of $e_j$, $j\in J$, are $\eps/\gamma^2$ in the window and $1/L^2$ on
the tail, as in Theorem~\ref{thm:price}; the distractor directions have
price $0$ and no bad arm uses them.

\begin{proposition}[The price ratio collected as $I_\tau$]
\label{prop:cat-tail-saving}
Let $\mathcal J=\mathcal J(K,K',L,\gamma,\eps,n)$ under \eqref{eq:conventions}
strengthened to $L\ge13$ and $\lambda\ge16\sigma^2d\,\ell_0$, with
$n\ge36\lambda L^2/\gamma^2$, $K\ge36L^2$, $K'\ge L'^2Kn/\lambda$, and
$T-\tau_0\ge Km_0$. With probability at least $1-3\delta$ the following hold
simultaneously.
\begin{enumerate}[itemsep=1pt,topsep=1pt,leftmargin=*]
\item[(i)] \emph{(Expert.)} $R_T(\Exp)=\eps Kn+(T-Kn)$.
\item[(ii)] \emph{(Cold.)} Cold LinUCB plays the distractor at every round
$t\le\tau_0$ and the bad arm at each of the first $m_0$ rounds at which every
tail direction is offered, so $R_{[1,\tau_0]}(\Cold)=0$ and
$Km_0\le R_{[\tau_0+1,T]}(\Cold)\le K(4\beta^2+1)$.
\item[(iii)] \emph{(Hybrid.)} $\Hyb_{\tau_0}$ plays $g$ at every round
$t>\tau_0$, so $R_T(\Hyb_{\tau_0})=\eps Kn$.
\end{enumerate}
Consequently $A_{\tau_0}=-\eps Kn<0$, $Km_0\le I_{\tau_0}\le K(16\lambda+1)$,
$B_{\tau_0}=T-Kn$, and
\[
\frac{m_0}{\eps n}\ \le\ \frac{R_T(\Cold)}{R_T(\Hyb_{\tau_0})}\ \le\ \frac{16\lambda+1}{\eps n},
\]
which at the smallest admissible $n$ lies between $\gamma^2/(2220\,\eps L^2)$ and
$\gamma^2/(2\eps L^2)$: the bracket of Theorem~\ref{thm:price} up to its constant,
now carried entirely by $I_{\tau_0}$.
\end{proposition}

The proof (Appendix~\ref{app:proof-tail-saving}) has one new ingredient: on
the confidence event with the true radius, which the strengthened condition
on $\lambda$ makes at most $\tfrac58\beta$, every index satisfies
$U_t(x)\ge\mu(x)+\tfrac38\beta\|x\|_{V_t^{-1}}$, and the distractor's width is
at least $L'/\sqrt{2\lambda}$ as long as it has been pulled at most
$\lambda/L'^2$ times, which the condition on $K'$ guarantees. Its index then
exceeds $7$ while the indices of $g$ and of the probe are below $4.3$ and
$5.6$. Cold LinUCB never probes, reaches the tail with every bad direction
untouched, and is trapped there exactly as in Theorem~\ref{thm:price}(ii);
the hybrid's window is the same as in Theorem~\ref{thm:price}(iv). In the
language of Theorem~\ref{thm:general-price}, rows~5 and~$5'$ have the same
certified set $U=\{e_j:j\in J\}$, the same prices, and the same $P_U=\gamma^2/(\eps L^2)$;
they differ only in which rounds cold LinUCB's counterfactual exploration
falls on. On the real streams of Section~\ref{sec:amazon}, too, the saving
is collected after routing stops, as in row~$5'$, although there its source
is mostly the transfer of Section~\ref{sec:two-channels}.


\paragraph{Row 6: the same expert, optimal on the tail.}
\begin{proposition}[Expert optimal on the tail: large $I_\tau$, no gain over the expert]
\label{prop:cat-tail}
Let $\mathcal I=\mathcal I(K,L,\gamma,\eps,Kn)$ with $n\ge36\lambda L^2/\gamma^2$ and
$K\ge36L^2$, and let $\Exp=\Exp_w$ with $w\in(\eps/\gamma,1/L)$, which plays every
probe and then $g$ at every round. With probability at least $1-3\delta$,
$R_T(\Exp)=R_T(\Hyb_{\tau_0})=\eps Kn$ and $B_{\tau_0}=0$, while $A_{\tau_0}$
and $I_{\tau_0}$ are exactly as in Theorem~\ref{thm:price}.
\end{proposition}

\paragraph{Row 7: probes that cover the wrong directions.}
\begin{proposition}[Probes cover the wrong directions]
\label{prop:cat-wrong}
Let $d=2K+1$ and replace the probe by
$p'_t:=(1-\eps)e_1+\gamma\,e_{K+j(t)}$, whose second coordinate is a nuisance
direction that no bad arm uses, with $\mathcal A_t=\{g,b_t,p'_t\}$ for
$t\le\tau_0=Kn$ and $\{g,b_t\}$ afterwards. Let $\Exp$ play $p'_t$ for
$t\le\tau_0$ and $b_t$ afterwards (for instance the linear policy with
$\widehat\theta_{\Exp}=e_1+w\sum_{j\ge2}e_j$, $w\in(1/L,(1-\eps)/(L-\gamma))$).
Assume $n\ge36\lambda L^2/\gamma^2$, $K\ge36L^2$, and $T-\tau_0\ge K\lfloor\lambda/2\rfloor$.
With probability at least $1-3\delta$,
\[
R_T(\Hyb_{\tau_0})\ \ge\ \eps Kn+K\lfloor\lambda/2\rfloor,
\qquad
R_T(\Cold)\ \le\ K(4\beta^2+1)+\eps Kn,
\qquad
\frac{R_T(\Cold)}{R_T(\Hyb_{\tau_0})}\le65 .
\]
The principal angle between the expert's Gram matrix and the span of the bad
directions is $90^\circ$ and $I_{\tau_0}\le R_{[\tau_0+1,T]}(\Cold)-K\lfloor\lambda/2\rfloor$.
\end{proposition}

\paragraph{Row 8: probes that are too dear.}
\begin{corollary}[Row 8: probes no cheaper than cold exploration]
\label{cor:cat-dear}
In the setting of Theorem~\ref{thm:price}, if $\eps\ge \gamma^2/L^2$ (the probe's
price $\eps/\gamma^2$ is at least cold exploration's price $1/L^2$) then
$R_T(\Cold)/R_T(\Hyb_{\tau_0})\le2$ for every $n\ge36\lambda L^2/\gamma^2$.
\end{corollary}

\paragraph{Row 9: cold LinUCB already cheap.}
\begin{proposition}[Cold LinUCB already cheap: no multiplier in average reward]
\label{prop:cat-cheap}
On the no-probe instance $\mathcal I(K,L,\cdot,\cdot,0)$, with probability
at least $1-\delta$, $R_T(\Cold)\le K(4\beta^2+1)=K(16\lambda+1)$. Hence for
\emph{every} policy $\pi$, $(R_T(\Cold)-R_T(\pi))/T\le K(16\lambda+1)/T$, which with
$\lambda=\sigma^2d\ell_0$ is $O(\sigma^2K^2\ell_0/T)$.
\end{proposition}

\paragraph{Row 10: the gap.}
\begin{proposition}[The gap enters only through the price]
\label{prop:cat-gap}
Fix $\eta\in(0,1]$ and replace $g$ by $g_\eta:=\eta e_1$ ($\mu=\eta$) and
$p_t$ by $p_{t,\eta}:=(1-\eps)\eta e_1+\gamma\,e_{j(t)}$ ($\mu=(1-\eps)\eta$), so that
the gaps of $b_t$ and $p_t$ are $\eta$ and $\eps\eta$ and the prices of
$e_j$ are $\eta/L^2$ and $\eps\eta/\gamma^2$. Let $\Exp=\Exp_w$ with
$w\in(\eta/L,(1-\eps)\eta/(L-\gamma))$. Under \eqref{eq:conventions}, with
$n\ge36\lambda L^2/(\eta^2\gamma^2)$, $K\ge36L^2/\eta^2$, and $\tau=\tau_0=Kn$, with
probability at least $1-3\delta$,
\[
\frac{\eta Km_0}{\eps\eta Kn}\ \le\
\frac{R_T(\Cold)}{R_T(\Hyb_{\tau_0})}\ \le\ 1+\frac{8\beta^2/\eta^2+1}{\eps n}\ \le\ 1+\frac{\gamma^2}{\eps L^2},
\]
and at the smallest admissible $n$ the left-hand side is at least
$\eta^2\gamma^2/(2220\,\eps L^2)$.
The upper bound is the price ratio $(\eta/L^2)/(\eps\eta/\gamma^2)=\gamma^2/(\eps L^2)$,
in which $\eta$ cancels; the lower bound loses a factor $\eta^2$ because our
lower bound on cold LinUCB counts only the $m_0$ pulls per direction that
optimism forces regardless of the gap.
\end{proposition}

\paragraph{Where the experiments sit.}
The experiments of Section~\ref{sec:experiments} relate to the rows as
follows. AmazonCat with Flan-T5 has the timing of row~$5'$: multipliers
$1.71\times$, $2.72\times$ and $3.29\times$, routed rounds concentrated in the
first few thousand, and $96$--$100\%$ of the saving accruing after the last
routed round (Figure~\ref{fig:amazon-dynamics}). Its source, however, is
mostly the transfer of Section~\ref{sec:two-channels}, which lies outside the
linear model of the catalog: the three experts buy the same coverage at the
same price and still differ (Section~\ref{sec:amazon}). AmazonCat with BGE
features shows the same timing with a smaller gain (up to $1.23\times$ on
held-out seeds), since cold LinUCB already reaches about a third of the
correct answers. WikiLinks with Flan-T5 does not multiply (at most
$1.05\times$): the routed data does reach useful directions ($I_\tau>0$ on
every seed), but no entity is frequent enough for the expert to steer the
learner toward (Section~\ref{sec:wiki}). WikiLinks with Qwen2.5 is row~3: the
expert is better than the learner on its own, and at the selected $c_{\mathrm{LLM}}$
the Router simply follows it.

\section{Proofs}
\label{app:proofs}

Section~\ref{app:proof-decomp} proves the decomposition identity.
Section~\ref{app:proof-three-way} proves
Proposition~\ref{prop:three-way}.
Section~\ref{app:proof-one-pass} proves Proposition~\ref{prop:one-pass}.
Sections~\ref{app:proof-width}--\ref{app:proof-transfer} prove
Lemma~\ref{lem:width-limited} and the two width identities of
Appendix~\ref{app:widths}, Section~\ref{app:proof-handoff} proves
Proposition~\ref{thm:handoff}, and the remaining sections prove the results of
Sections~\ref{sec:price} and~\ref{sec:two-channels}, including
Theorem~\ref{thm:general-price}.

\subsection{Proof of Proposition~\ref{prop:decomp}}
\label{app:proof-decomp}

\begin{proof}
Split $[T]$ into $[1,\tau]$ and $[\tau+1,T]$. On $[1,\tau]$ the hybrid plays
$a_t^{\Exp}$, so $R_{[1,\tau]}(\Hyb_\tau)=R_{[1,\tau]}(\Exp)$. Hence
\[
R_T(\Cold)-R_T(\Hyb_\tau)
=\bigl(R_{[1,\tau]}(\Cold)-R_{[1,\tau]}(\Exp)\bigr)
+\bigl(R_{[\tau+1,T]}(\Cold)-R_{[\tau+1,T]}(\Hyb_\tau)\bigr)
=A_\tau+I_\tau ,
\]
and
\[
R_T(\Exp)-R_T(\Hyb_\tau)
=\bigl(R_{[1,\tau]}(\Exp)-R_{[1,\tau]}(\Exp)\bigr)
+\bigl(R_{[\tau+1,T]}(\Exp)-R_{[\tau+1,T]}(\Hyb_\tau)\bigr)
=B_\tau .
\]
Both are identities between sums of gaps, so they hold for every noise
realization and therefore in expectation.
\end{proof}

\subsection{Proof of Proposition~\ref{prop:three-way}}
\label{app:proof-three-way}

We first record the standard subgaussian concentration lemma used in the
high-probability part of the argument, and then construct the linear bandit
instance and verify the three claims of Proposition~\ref{prop:three-way}.

\begin{lemma}[Subgaussian concentration; {\citealp[Corollary~5.5]{lattimore2020bandit}}]
\label{lem:subgaussian}
Assume that $X_1-\mu,\dots,X_n-\mu$ are independent, $\sigma$-subgaussian random variables. Then for any $\varepsilon \ge 0$,
\[
\mathbb{P}\!\left(\hat{\mu} \ge \mu+\varepsilon\right)
\le
\exp\!\left(-\frac{n\varepsilon^2}{2\sigma^2}\right),
\qquad
\mathbb{P}\!\left(\hat{\mu} \le \mu-\varepsilon\right)
\le
\exp\!\left(-\frac{n\varepsilon^2}{2\sigma^2}\right),
\]
where
\[
\hat{\mu}=\frac{1}{n}\sum_{t=1}^n X_t.
\]
\end{lemma}

\subparagraph{Problem setup.}
Let the dimension satisfy $d \ge n+1$, where $n$ is the horizon. The unknown parameter is
\[
\theta^\star = e_1 \in \mathbb{R}^d,
\]
where $\{e_i\}_{i=1}^d$ denotes the standard basis. Rewards follow the stochastic linear model
\[
X_t = \langle \theta^\star, A_t \rangle + \eta_t,
\]
where $\eta_t$ are independent mean-zero $\sigma$-subgaussian random variables.

At each round $t=1,\dots,n$, the learner is offered two actions,
\[
g:=e_1,
\qquad
b_t:=L e_{t+1},
\]
where $L>1$ is a fixed constant. Their expected rewards are
\[
\langle \theta^\star,g\rangle = 1,
\qquad
\langle \theta^\star,b_t\rangle = 0.
\]
Both LinUCB variants below use the rule
\[
\mathrm{UCB}_t(a)
=
a^\top \hat{\theta}_{t-1}
+
\beta \|a\|_{V_{t-1}^{-1}},
\]
with
\[
V_{t-1}=V_0+\sum_{s<t} A_sA_s^\top,
\qquad
\hat{\theta}_{t-1}
=
V_{t-1}^{-1}\sum_{s<t} A_s X_s.
\]

\subparagraph{Algorithm 1: a constant policy.}
Consider a policy $\pi_{\mathrm{const}}$ that selects an action by maximizing $\langle \hat{\theta}_{\mathrm{const}}, a\rangle$ for some fixed parameter estimate $\hat{\theta}_{\mathrm{const}}\in\mathbb{R}^d$ that is never updated. For concreteness, take
\[
\hat{\theta}_{\mathrm{const}}
=
\frac{1}{L}\sum_{j=2}^{n+1} e_j,
\]
so that
\[
\langle \hat{\theta}_{\mathrm{const}},g\rangle = 0,
\qquad
\langle \hat{\theta}_{\mathrm{const}},b_t\rangle = 1
\quad\text{for every } t\le n.
\]
Then $\pi_{\mathrm{const}}$ selects $b_t$ at every round $t$. Since each pull of $b_t$ yields expected reward $0$ while $g$ has expected reward $1$, the cumulative pseudo-regret of $\pi_{\mathrm{const}}$ on the horizon $n$ is exactly $n$. This proves part (1) of Proposition~\ref{prop:three-way}.

The failure mode is structural rather than statistical: $\pi_{\mathrm{const}}$ never queries direction $e_1$, so no amount of data can refine $\hat{\theta}_{\mathrm{const}}$ in the reward-relevant direction. The policy has no learning mechanism. We will see, however, that the trajectory it produces is precisely what suppresses the misplaced optimism of LinUCB in the bad directions.

\subparagraph{Algorithm 2: LinUCB with isotropic initialization.}
The first LinUCB variant uses
\[
V_0=\lambda I_d,
\qquad
\hat{\theta}_0=0.
\]
At $t=1$,
\[
\|g\|_{V_0^{-1}}^2=\frac{1}{\lambda},
\qquad
\|b_1\|_{V_0^{-1}}^2=\frac{L^2}{\lambda}.
\]
Since $\hat{\theta}_0=0$, we have
\[
\mathrm{UCB}_1(b_1)>\mathrm{UCB}_1(g),
\]
so the algorithm selects $b_1$.

Now suppose inductively that the algorithm has selected $b_1,\dots,b_{t-1}$ in the first $t-1$ rounds. Then
\[
V_{t-1}
=
\lambda I_d+\sum_{s=1}^{t-1} b_s b_s^\top
=
\lambda I_d + L^2\sum_{j=2}^{t} e_j e_j^\top.
\]
Hence $V_{t-1}$ is diagonal in the standard basis. In particular, for the fresh direction $e_{t+1}$,
\[
e_{t+1}^\top V_{t-1}^{-1} e_{t+1}=\frac{1}{\lambda},
\qquad
\|b_t\|_{V_{t-1}^{-1}}^2=\frac{L^2}{\lambda},
\qquad
\|g\|_{V_{t-1}^{-1}}^2=\frac{1}{\lambda}.
\]

The key point is that although the observed rewards on $b_1,\dots,b_{t-1}$ contain noise, the estimate
\[
\hat{\theta}_{t-1}
=
V_{t-1}^{-1}\sum_{s=1}^{t-1} b_s X_s
\]
lies entirely in $\mathrm{span}\{e_2,\dots,e_t\}$. Therefore,
\[
g^\top \hat{\theta}_{t-1}=0,
\qquad
b_t^\top \hat{\theta}_{t-1}=0,
\]
because both $g=e_1$ and the fresh arm $b_t=L e_{t+1}$ are orthogonal to $\mathrm{span}\{e_2,\dots,e_t\}$.

Thus, even in the stochastic-reward setting, the comparison between $g$ and $b_t$ is determined purely by the confidence terms:
\[
\mathrm{UCB}_t(g)=\frac{\beta}{\sqrt{\lambda}},
\qquad
\mathrm{UCB}_t(b_t)=\frac{\beta L}{\sqrt{\lambda}}.
\]
Since $L>1$, we have $\mathrm{UCB}_t(b_t)>\mathrm{UCB}_t(g)$ for all $t\le n$. Hence the isotropically initialized algorithm selects $b_t$ at every round and, like $\pi_{\mathrm{const}}$, incurs cumulative pseudo-regret $n$. This proves part (2) of Proposition~\ref{prop:three-way}.

\subparagraph{Algorithm 3: LinUCB with an anisotropic warm start.}
We now define an initialization that is realized by exactly the trajectory $\pi_{\mathrm{const}}$ would produce. Fix an integer $m\ge 1$, and let
\[
V_0'
=
\lambda I_d + mL^2 \sum_{j=2}^{n+1} e_j e_j^\top,
\qquad
\hat{\theta}_0'=0.
\]
This corresponds to having observed each bad direction $L e_j$, $j=2,\dots,n+1$, exactly $m$ times before the online interaction begins --- precisely the kind of data $\pi_{\mathrm{const}}$ generates if rolled out for $mn$ rounds on a sequence of action sets containing each bad direction. Thus $V_0'$ is a genuine design matrix obtainable from a collection of outer products produced by a constant logging policy.

At $t=1$,
\[
\|g\|_{(V_0')^{-1}}^2=\frac{1}{\lambda},
\qquad
\|b_1\|_{(V_0')^{-1}}^2
=
L^2\cdot \frac{1}{\lambda+mL^2}
=
\frac{L^2}{\lambda+mL^2}.
\]
Hence, if
\[
m > \frac{\lambda(L^2-1)}{L^2},
\]
then $\frac{1}{\lambda}>\frac{L^2}{\lambda+mL^2}$, and therefore
\[
\mathrm{UCB}_1'(g)>\mathrm{UCB}_1'(b_1).
\]
So the warm-started algorithm selects $g$ at round $1$.

Suppose now that the algorithm continues to select $g$ for the first $s$ rounds. Then
\[
V_s'
=
\lambda I_d
+
mL^2 \sum_{j=2}^{n+1} e_j e_j^\top
+
s\, e_1 e_1^\top.
\]
Since each reward observed on $g$ has mean $1$,
\[
\sum_{i=1}^s X_i = s + \sum_{i=1}^s \eta_i,
\]
and therefore
\[
g^\top \hat{\theta}_s'
=
e_1^\top (V_s')^{-1}\sum_{i=1}^s g X_i
=
\frac{1}{\lambda+s}\sum_{i=1}^s X_i.
\]
Applying Lemma~\ref{lem:subgaussian} with $\mu=1$, we obtain that with probability at least $1-\delta$, simultaneously for all $1\le s\le n$,
\[
\frac{1}{s}\sum_{i=1}^s X_i
\ge
1-\sqrt{\frac{2\sigma^2\log(2n/\delta)}{s}}.
\]
On this event,
\[
g^\top \hat{\theta}_s'
\ge
\frac{s}{\lambda+s}
\left(
1-\sqrt{\frac{2\sigma^2\log(2n/\delta)}{s}}
\right).
\]

For every fresh bad arm $b_{s+1}=Le_{s+2}$,
\[
b_{s+1}^\top \hat{\theta}_s' = 0,
\]
since the online trajectory has only sampled the direction $e_1$, and the warm-start contributes only to the design matrix, not to the response vector. Moreover,
\[
\|b_{s+1}\|_{(V_s')^{-1}}^2
=
\frac{L^2}{\lambda+mL^2},
\qquad
\|g\|_{(V_s')^{-1}}^2
=
\frac{1}{\lambda+s}.
\]

On the same high-probability event,
\[
\mathrm{UCB}_{s+1}'(g)
\ge
\frac{s}{\lambda+s}
\left(
1-\sqrt{\frac{2\sigma^2\log(2n/\delta)}{s}}
\right)
+
\frac{\beta}{\sqrt{\lambda+s}},
\]
while
\[
\mathrm{UCB}_{s+1}'(b_{s+1})
=
\frac{\beta L}{\sqrt{\lambda+mL^2}}.
\]
Consequently, if $m$ is chosen large enough so that for every $1\le s\le n$,
\begin{equation}
\label{eq:m-condition}
\frac{s}{\lambda+s}
\left(
1-\sqrt{\frac{2\sigma^2\log(2n/\delta)}{s}}
\right)
+
\frac{\beta}{\sqrt{\lambda+s}}
>
\frac{\beta L}{\sqrt{\lambda+mL^2}},
\end{equation}
then, with probability at least $1-\delta$, the warm-started algorithm selects $g$ at every round, and the cumulative pseudo-regret over the $n$ online rounds is $0$ on this event. This proves part (3) of Proposition~\ref{prop:three-way}; condition~\eqref{eq:m-condition} is the explicit ``$m$ sufficiently large'' requirement referenced in the proposition statement. The condition is satisfiable by a finite $m$ whenever the left-hand side is positive for every $s$, which holds for any $\beta\ge\sigma\sqrt{2\log(2n/\delta)}$ (in particular for the radius of \citealp{abbasi2011improved}), since then $\beta/\sqrt{\lambda+s}\ge\sigma\sqrt{2\log(2n/\delta)}\,\sqrt s/(\lambda+s)$. \qed


\begin{remark}
The warm-started algorithm in the proof uses $\widehat\theta_0'=0$, i.e., it
carries the offline Gram matrix but not the offline response sum. Because
every offline observation is of an arm with expected reward $0$, the offline
response sum has mean zero in each bad direction and carrying it would only
add a mean-zero perturbation to $\widehat\theta$ in directions orthogonal to
$g$; the comparison between $g$ and $b_{s+1}$ is unchanged up to the noise
term already absorbed in the choice of $m$, so we omit it for simplicity.
\end{remark}

\subsection{Proof of Proposition~\ref{prop:one-pass}}
\label{app:proof-one-pass}

\begin{proof}
Work on the instance of Section~\ref{app:proof-three-way} with horizon $T$:
$\theta^\star=e_1$, $\mathcal A_t=\{g,b_t\}$ with $g=e_1$ and
$b_t=Le_{t+1}$. Fix an arbitrary expert $\Exp$ and a handoff time $\tau$.
Whatever the expert plays during rounds $1,\dots,\tau$, and whatever LinUCB
plays afterwards, the set of arms played before round $t$ is contained in
$\{g\}\cup\{b_s:s<t\}$, all of which are orthogonal to $e_{t+1}$. Hence for
every $t$, the design matrix $V_t=\lambda I+\sum_{s<t}x_{s,A_s}x_{s,A_s}^\top$
is block diagonal with the block on $e_{t+1}$ equal to $\lambda$, and
$\widehat\theta_t=V_t^{-1}c_t$ lies in $\mathrm{span}\{e_1,\dots,e_t\}$. Therefore
\[
\|b_t\|_{V_t^{-1}}^2=\frac{L^2}{\lambda},\qquad \langle\widehat\theta_t,b_t\rangle=0,
\qquad\text{for every }t\text{ and every trajectory.}
\]
So $\mathrm{UCB}_t(b_t)=\beta L/\sqrt\lambda$ exactly, for all $t$.

Now bound $\mathrm{UCB}_t(g)$. Let $s_t$ be the number of pulls of $g$ before
round $t$ and $X_1,\dots,X_{s_t}$ the rewards observed on those pulls; since
$V_t$ is diagonal with entry $\lambda+s_t$ on $e_1$ and the rewards on bad
arms enter $c_t$ only in directions orthogonal to $e_1$,
\[
\langle\widehat\theta_t,g\rangle=\frac{1}{\lambda+s_t}\sum_{i=1}^{s_t}X_i,
\qquad
\|g\|_{V_t^{-1}}^2=\frac{1}{\lambda+s_t}\le\frac1\lambda .
\]
By Lemma~\ref{lem:subgaussian} and a union bound over $s\in\{1,\dots,T\}$,
with probability at least $1-\delta$ we have
$\sum_{i\le s}X_i\le s+\sqrt{2\sigma^2s\log(2T/\delta)}$ for all $s\le T$
simultaneously. On that event,
\[
\langle\widehat\theta_t,g\rangle
\le\frac{s_t}{\lambda+s_t}+\frac{\sqrt{2\sigma^2s_t\log(2T/\delta)}}{\lambda+s_t}
\le 1+\frac{\sigma\sqrt{2\log(2T/\delta)}}{\sqrt{\lambda+s_t}}
\le 1+\zeta ,
\]
using $\sqrt{s_t}/(\lambda+s_t)\le1/\sqrt{\lambda+s_t}\le1/\sqrt\lambda$.
Hence $\mathrm{UCB}_t(g)\le1+\zeta+\beta/\sqrt\lambda$. Under the hypothesis
$\beta(L-1)/\sqrt\lambda>1+\zeta$ this is strictly less than
$\beta L/\sqrt\lambda=\mathrm{UCB}_t(b_t)$, so LinUCB, whenever it acts, plays
$b_t$. The argument did not use anything about rounds $1,\dots,\tau$ beyond
the fact that they cannot produce coverage of $e_{t+1}$ for $t>\tau$, so it
holds for every expert and every $\tau$. Since $\Delta_{t,b_t}=1$,
$R_{[\tau+1,T]}(\Hyb_\tau)=T-\tau$. Cold LinUCB plays $b_t$ on every round by
part~(2) of Proposition~\ref{prop:three-way} (no condition on $\beta$ is
needed there because $\widehat\theta_t\perp g$ on its own trajectory), so
$R_{[\tau+1,T]}(\Cold)=T-\tau$ and $I_\tau=0$.
\end{proof}

\subsection{Proof of Lemma~\ref{lem:width-limited}}
\label{app:proof-width}

\begin{proof}
Fix a round $t$ at which LinUCB acts, and write $x^\star=x_{t,a_t^\star}$,
$x=x_{t,A_t}$, $\wid(u)=\|u\|_{V_t^{-1}}$. On $\mathcal E_\delta$, Cauchy--Schwarz
gives $|\langle\widehat\theta_t-\theta^\star,u\rangle|\le\beta_t\wid(u)$ for
every $u$. By the selection rule,
$\langle\widehat\theta_t,x\rangle+\beta_t\wid(x)\ge\langle\widehat\theta_t,x^\star\rangle+\beta_t\wid(x^\star)\ge\langle\theta^\star,x^\star\rangle$,
and $\langle\widehat\theta_t,x\rangle\le\langle\theta^\star,x\rangle+\beta_t\wid(x)$.
Combining, $\Delta_{t,A_t}=\langle\theta^\star,x^\star-x\rangle\le2\beta_t\wid(x)$.
Thus $\Delta_{t,A_t}=\min\{\Delta_{t,A_t},2\beta_t\wid(x_{t,A_t})\}\le\max_a\min\{\Delta_{t,a},2\beta_t\wid(x_{t,a})\}=\Phi_t(V_t)$.
Monotonicity: $V\preceq V'$ implies $V'^{-1}\preceq V^{-1}$, so every width is
non-increasing and so is each term in the maximum.
\end{proof}

\subsection{Proof of Lemma~\ref{lem:loewner}}
\label{app:proof-loewner}

\begin{proof}
$\lambda I+M\succeq\lambda I$ gives $(\lambda I+M)^{-1}\preceq(\lambda I)^{-1}$,
which is the first claim. For the second, $\lambda I+M\succeq\lambda I+mP$, so
$x^\top(\lambda I+M)^{-1}x\le x^\top(\lambda I+mP)^{-1}x$, and for
$x\in\mathrm{Im}(P)$ the right-hand side is $\|x\|_2^2/(\lambda+m)$ because
$(\lambda I+mP)^{-1}$ acts as $(\lambda+m)^{-1}$ on $\mathrm{Im}(P)$. For the
third, if $x\perp\mathrm{Im}(M)$ then $x$ is an eigenvector of $\lambda I+M$
with eigenvalue $\lambda$, so $x^\top(\lambda I+M)^{-1}x=\|x\|_2^2/\lambda$.
\end{proof}

\subsection{Proof of Lemma~\ref{lem:transfer}}
\label{app:proof-transfer}

\begin{proof}
By the Sherman--Morrison formula,
$(V+xx^\top)^{-1}=V^{-1}-\dfrac{V^{-1}xx^\top V^{-1}}{1+x^\top V^{-1}x}$. Multiplying
on both sides by $z$,
\[
z^\top(V+xx^\top)^{-1}z
= z^\top V^{-1}z-\frac{(z^\top V^{-1}x)^2}{1+\|x\|^2_{V^{-1}}}
=\|z\|^2_{V^{-1}}\Bigl(1-\frac{(z^\top V^{-1}x)^2}{(1+\|x\|^2_{V^{-1}})\|z\|^2_{V^{-1}}}\Bigr).
\]
The bracket is at least $1-\dfrac{\|x\|^2_{V^{-1}}\|z\|^2_{V^{-1}}}{(1+\|x\|^2_{V^{-1}})\|z\|^2_{V^{-1}}}>0$
by Cauchy--Schwarz in the $V^{-1}$ inner product, so $\rho_V\in[0,1)$. At
$V=\lambda I$, $z^\top V^{-1}x=\langle z,x\rangle/\lambda$,
$\|x\|^2_{V^{-1}}=\|x\|_2^2/\lambda$, $\|z\|^2_{V^{-1}}=\|z\|_2^2/\lambda$, and
substituting gives
$\rho_{\lambda I}(x,z)=\dfrac{\langle x,z\rangle^2}{\|x\|_2^2\|z\|_2^2}\cdot\dfrac{\|x\|_2^2}{\lambda+\|x\|_2^2}$.
\end{proof}

\subsection{Proof of Proposition~\ref{thm:handoff}}
\label{app:proof-handoff}

\begin{proof}
The acquisition term is the definition of $R_{[1,\tau]}(\Hyb_\tau)=R_{[1,\tau]}(\Exp)$.
For the LinUCB phase we give two bounds and take the minimum.

\emph{First bound.} By Lemma~\ref{lem:width-limited} applied at each
$t>\tau$, $\sum_{t>\tau}\Delta_{t,A_t}\le\sum_{t>\tau}\Phi_t(V_t)$.

\emph{Second bound.} From the proof of Lemma~\ref{lem:width-limited},
$\Delta_{t,A_t}\le2\beta_t\|x_{t,A_t}\|_{V_t^{-1}}\le2\beta_T\|x_{t,A_t}\|_{V_t^{-1}}$
for $t>\tau$ (the radius $\beta_t$ is non-decreasing in $t$). Cauchy--Schwarz
over $t=\tau+1,\dots,T$ gives
\[
\sum_{t>\tau}\Delta_{t,A_t}\le2\beta_T\sqrt{(T-\tau)\sum_{t>\tau}\|x_{t,A_t}\|^2_{V_t^{-1}}} .
\]
By the elliptical potential lemma \citep[Lemma~11]{abbasi2011improved},
\[
\sum_{t>\tau}\min\{1,\|x_{t,A_t}\|^2_{V_t^{-1}}\}\le2\log\frac{\det V_{T+1}}{\det V_{\tau+1}};
\]
since $\|x_{t,A_t}\|^2_{V_t^{-1}}\le L^2/\lambda$, the truncation at $1$ costs at
most a factor $\max\{1,L^2/\lambda\}$ on the sum of squares, which is the
factor $\max\{1,L/\sqrt\lambda\}$ in \eqref{eq:handoff-bound} after the square
root; with $L^2\le\lambda$ nothing is lost. This gives
the second term of \eqref{eq:handoff-bound}.

\emph{Determinant ratio.} Let $V_{\tau+1}=\sum_j\sigma_ju_ju_j^\top$ be an
eigendecomposition and $G=\sum_{t>\tau}x_{t,A_t}x_{t,A_t}^\top$. Then
$\det(V_{\tau+1}+G)/\det V_{\tau+1}=\det(I+V_{\tau+1}^{-1/2}GV_{\tau+1}^{-1/2})\le\prod_j(1+u_j^\top Gu_j/\sigma_j)$
by Hadamard's inequality for positive definite matrices, and
$u_j^\top Gu_j=\sum_{t>\tau}\langle u_j,x_{t,A_t}\rangle^2\le L^2(T-\tau)$.
Taking logarithms gives the stated bound. If all played actions lie in an
$r$-dimensional subspace $\mathrm{Im}(P)$ then $G$ is supported there and only
the $r$ eigenvalues of $V_{\tau+1}$ restricted to $\mathrm{Im}(P)$ appear, which
is how the factor $r$ in the remark after the theorem arises; if in addition
$V_{\tau+1}\succeq(\lambda+\nu\tau)P$ on that subspace, each $\sigma_j\ge\lambda+\nu\tau$.

The cold bound is the case $\tau=0$, $V_1=\lambda I$, $\sigma_j=\lambda$. The
ratio $\Psi$ in the text is the quotient of the two isotropic bounds and is
stated only to show that it does not tend to zero.
\end{proof}

\begin{remark}[Validity of the confidence event with an online warm start]
The hybrid's LinUCB phase uses $(V_{\tau+1},c_{\tau+1})$ built from the
expert-routed rounds together with the observed rewards on those rounds. The
self-normalized bound of \citet{abbasi2011improved} applies to any sequence of
(feature, reward) pairs in which the feature at round $s$ is measurable with
respect to the history before $s$, whoever chose it, so the event
$\mathcal E_\delta$ covers the hybrid's trajectory. If instead the warm start
carried $M_\tau$ but dropped $c_\tau$ (``passive'' updates), the estimate
$\widehat\theta_t$ would be biased toward zero in the directions of $M_\tau$
and the event would not be guaranteed; the Router in our experiments updates
both, and Algorithm~\ref{alg:confidence-gated} reflects this.
\end{remark}

\subsection{Proofs for Section~\ref{sec:two-channels}}
\label{app:proof-two-channels}

\begin{proofof}{Proposition~\ref{prop:two-channel}}
Write $x^\star=x_{t,a_t^\star}$, $x=x_{t,A_t}$, $\wid(u)=\|u\|_{V_t^{-1}}$.
Cauchy--Schwarz in the $V_t$ inner product gives
$|\langle\widehat\theta_t-\theta^\star,u\rangle|\le\rho_t\wid(u)$ for every $u$;
this is the inequality in the proof of Lemma~\ref{lem:width-limited} with the
realized error in place of its bound. The selection rule gives
\[
\langle\widehat\theta_t,x\rangle+\beta_t\wid(x)
\ \ge\ \langle\widehat\theta_t,x^\star\rangle+\beta_t\wid(x^\star)
\ \ge\ \langle\theta^\star,x^\star\rangle+(\beta_t-\rho_t)\wid(x^\star)
\ \ge\ \langle\theta^\star,x^\star\rangle,
\]
using $\rho_t\le\beta_t$, and
$\langle\widehat\theta_t,x\rangle\le\langle\theta^\star,x\rangle+\rho_t\wid(x)$.
Subtracting, $\Delta_{t,A_t}\le(\beta_t+\rho_t)\wid(x)$, and taking the maximum
over $a$ of the minimum with the gap gives \eqref{eq:phi-two}. Monotonicity in
$V$ is Lemma~\ref{lem:width-limited}; in $\rho$ it is immediate. Finally
$\min\{\Delta,(\beta+\rho)w\}\ge\min\{\Delta,\beta w\}\ge\tfrac12\min\{\Delta,2\beta w\}$.
\end{proofof}

\begin{proofof}{Proposition~\ref{prop:response-sum}}
$c=\sum_sr_sx_s=\sum_s(\langle\theta^\star,x_s\rangle+\eta_s)x_s
=\bigl(\sum_sx_sx_s^\top\bigr)\theta^\star+\sum_s\eta_sx_s=M\theta^\star+\xi$.
Conditionally on the features, $\xi$ is Gaussian with mean zero and covariance
$\sigma^2\sum_sx_sx_s^\top=\sigma^2M$. Then
$\widehat\theta_1-\theta^\star=V_1^{-1}(M\theta^\star+\xi)-\theta^\star
=V_1^{-1}(M-V_1)\theta^\star+V_1^{-1}\xi=-\lambda V_1^{-1}\theta^\star+V_1^{-1}\xi$.
Both $V_1=\lambda I+M$ and the law of $\xi$ are functions of $M$ alone, so the
conditional law of $(V_1,c_1)$ given $M$ is the same whatever policy produced
the actions. The LinUCB phase is a deterministic function of $(V_1,c_1)$ and of
the future contexts and noise, which are independent of the warm start; hence
the LinUCB phases agree in law, and so does the expected tail regret of the
hybrid. $B_\tau$ is not covered, since $R_{[\tau+1,T]}(\Exp)$ depends on the
expert's own actions.
\end{proofof}

\begin{proofof}{Proposition~\ref{prop:transfer-regret}}
Fix $\widetilde\theta$ and write $\wid(u)=\|u\|_{V_t^{-1}}$,
$x^\star=x_{t,a^\star_t}$, $x=x_{t,A_t}$. Cauchy--Schwarz gives
$|\langle\widehat\theta_t-\widetilde\theta,u\rangle|\le\widetilde\rho_t\wid(u)$.
By the definition of $\widetilde\varepsilon$,
\[
\bar r(x^\star)=\langle\widetilde\theta,x^\star\rangle+\widetilde\varepsilon(x^\star)
\le\langle\widehat\theta_t,x^\star\rangle+\widetilde\rho_t\wid(x^\star)+\widetilde\varepsilon(x^\star).
\]
The selection rule gives
$\langle\widehat\theta_t,x^\star\rangle\le\langle\widehat\theta_t,x\rangle+\beta_t\wid(x)-\beta_t\wid(x^\star)$,
and
$\langle\widehat\theta_t,x\rangle\le\langle\widetilde\theta,x\rangle+\widetilde\rho_t\wid(x)
=\bar r(x)-\widetilde\varepsilon(x)+\widetilde\rho_t\wid(x)$.
Chaining the three displays,
\[
\bar r(x^\star)-\bar r(x)\le(\beta_t+\widetilde\rho_t)\wid(x)
+(\widetilde\rho_t-\beta_t)\wid(x^\star)
+\widetilde\varepsilon(x^\star)-\widetilde\varepsilon(x),
\]
and $(\widetilde\rho_t-\beta_t)\wid(x^\star)\le(\widetilde\rho_t-\beta_t)_+\wid(x^\star)$.
\end{proofof}

\subsection{Proof of Theorem~\ref{thm:general-price}}
\label{app:proof-general-price}

We work in the setting of Section~\ref{sec:price}, on the confidence event
$\mathcal E_\delta$ of cold LinUCB, and prove the two steps of the proof
sketch as lemmas.

\begin{lemma}[Cold LinUCB pays its own price, once per arm]
\label{lem:cold-price}
$R_T(\Cold)\le\sum_{x\in\mathcal S}\Delta(x)\min\{n^\star_x,N_x\}$.
\end{lemma}

\begin{proof}
Fix $x\in \mathcal S$ and suppose cold LinUCB plays $x$ at round $t$ after
$i$ earlier pulls of $x$. On $\mathcal E_\delta$ the index of $x$ is at most
$\mu(x)+2\beta_t\|x\|_{V_t^{-1}}$ and the index of $x^\star\in\mathcal A_t$
is at least $\mu(x^\star)$, so playing $x$ requires
$\Delta(x)\le2\beta\|x\|_{V_t^{-1}}$ (this is Lemma~\ref{lem:width-limited}).
The $i$ earlier pulls are among the summands of $V_t$, so
$V_t\succeq\lambda I+i\,xx^\top$; inversion reverses the Loewner order, and
the Sherman--Morrison identity gives
$\|x\|^2_{V_t^{-1}}\le\|x\|_2^2/(\lambda+i\|x\|_2^2)$. Hence
$\Delta(x)^2(\lambda+i\|x\|_2^2)\le4\beta^2\|x\|_2^2$, that is,
$\lambda+i\|x\|_2^2\le C^\star_x$. If $C^\star_x<\lambda$ this fails already at $i=0$, and
$x$ is never played. Otherwise every pull has
$i\le(C^\star_x-\lambda)/\|x\|_2^2$, so $x$ is pulled at most $n^\star_x$ times, and at
most $N_x$ times because it is offered on $N_x$ rounds. Each pull costs
$\Delta(x)$ because $x^\star\in\mathcal A_t$, and optimal pulls cost nothing.
Summing over $x\in \mathcal S$ gives the lemma.
\end{proof}

\begin{lemma}[Certification costs at least the price]
\label{lem:cert-cost}
For every policy $\pi$, every cutoff $t_0\in\{0,\dots,T\}$ and every set of
arms $U\subseteq U_{\mathrm{cert}}(t_0)$,
\[
\sum_{x\in U}(C^\star_x-\lambda)_+\ \le\ \sum_{x\in U}\sum_{t\le t_0}\langle x^\pi_t,u_x\rangle^2
\ \le\ \frac{R_{[1,t_0]}(\pi)}{\kappa^\pi_U(t_0)}+C^{\mathrm{free}}_U(t_0),
\]
with $1/\kappa^\pi_U(t_0):=0$ when $\kappa^\pi_U(t_0)=+\infty$.
\end{lemma}

\begin{proof}
Fix $x\in U$, write $V:=V^\pi_{t_0+1}$, and let
$C(u_x):=\sum_{t\le t_0}\langle x^\pi_t,u_x\rangle^2$ be the coverage of $u_x$
after round $t_0$. Then $x^\top Vx=\|x\|_2^2(\lambda+C(u_x))$, and
Cauchy--Schwarz in the $V$ inner product,
$\|x\|_2^4=\langle x,x\rangle^2\le\|x\|^2_{V^{-1}}\,x^\top Vx$, gives
$\|x\|^2_{V^{-1}}\ge\|x\|_2^2/(\lambda+C(u_x))$. Certification of $x$,
$4\beta^2\|x\|^2_{V^{-1}}<\Delta(x)^2$, therefore forces
$\lambda+C(u_x)>C^\star_x$; since also $C(u_x)\ge0$, $C(u_x)\ge(C^\star_x-\lambda)_+$,
which summed over $U$ is the first inequality. For the second, exchange the
order of summation: $\sum_{x\in U}C(u_x)=\sum_{t\le t_0}\mathrm{cov}_U(x^\pi_t)$.
The rounds with $\Delta(x^\pi_t)=0$ contribute $C^{\mathrm{free}}_U(t_0)$ by
definition. On a round with $\Delta(x^\pi_t)>0$, the definition of
$\kappa^\pi_U(t_0)$ as a minimum gives
$\mathrm{cov}_U(x^\pi_t)\le\Delta(x^\pi_t)/\kappa^\pi_U(t_0)$, trivially when
$\mathrm{cov}_U(x^\pi_t)=0$ (including the case $\kappa^\pi_U(t_0)=+\infty$).
Summing, these rounds contribute at most $R_{[1,t_0]}(\pi)/\kappa^\pi_U(t_0)$.
\end{proof}

\begin{proofof}{Theorem~\ref{thm:general-price}}
Split the sum of Lemma~\ref{lem:cold-price} into $x\in U$ and $x\in \mathcal S\setminus U$. For
$x\in U$, $\min\{n^\star_x,N_x\}\le n^\star_x\le(C^\star_x-\lambda)_+/\|x\|_2^2+1$, so
$\Delta(x)\min\{n^\star_x,N_x\}\le\kappa_x(C^\star_x-\lambda)_++\Delta(x)\le\kappa^{\Cold}_U(C^\star_x-\lambda)_++\Delta(x)$.
Summing over $U$ and applying Lemma~\ref{lem:cert-cost} bounds these terms by
$\kappa^{\Cold}_U\bigl(R_{[1,t_0]}(\pi)/\kappa^\pi_U(t_0)+C^{\mathrm{free}}_U(t_0)\bigr)+\sum_{x\in U}\Delta(x)$,
and $\kappa^{\Cold}_U/\kappa^\pi_U(t_0)=P_U$; the terms with $x\in \mathcal S\setminus U$
are carried unchanged. Finally $R_{[1,t_0]}(\pi)\le R_T(\pi)$ because
pseudo-regret is nonnegative on every round, so dividing by $R_T(\pi)$ gives
\eqref{eq:general-price}. For $\pi=\Hyb_\tau$
and $t_0=\tau$, $R_{[1,\tau]}(\Hyb_\tau)=R_{[1,\tau]}(\Exp)$ because the
hybrid plays the expert's actions on $[1,\tau]$.
\end{proofof}

\subsection{The Price Bound on the Template}
\label{app:price-template}

Here we check the bound of Theorem~\ref{thm:general-price} on the template
instance.

On the instance of Theorem~\ref{thm:price} take for $U$ the $K$ bad arms
$b_j=Le_j$, all of which are certified at $\tau_0$ by (iv). Then
$\kappa_{b_j}=1/L^2$, so $\kappa^{\Cold}_U=1/L^2$;
the expert plays only probes in the window, $\mathrm{cov}_U(p_j)=\gamma^2$ and
$\Delta(p_j)=\eps$, so $\kappa^\pi_U(\tau_0)=\eps/\gamma^2$ and $P_U=\gamma^2/(\eps L^2)$;
$C^{\mathrm{free}}_U=0$ because $g\perp U$; $\sum_U\Delta=K$; and the
uncertified arms are the probes, for which
$\Delta(p_j)\min\{n^\star_{p_j},N_{p_j}\}\le\eps n$. The theorem gives
$R_T(\Cold)/R_T(\Hyb_{\tau_0})\le \gamma^2/(\eps L^2)+(K+\eps Kn)/(\eps Kn)=\gamma^2/(\eps L^2)+1+1/(\eps n)$,
consistent with the upper bound $1+(8\beta^2+1)/(\eps n)\le1+\gamma^2/(\eps L^2)$
proved directly. The same computation applies verbatim to
Proposition~\ref{prop:cat-tail-saving}, where the distractors have gap $0$
and contribute nothing.

\subsection{Proof of Corollary~\ref{cor:diverge}}
\label{app:proof-diverge}

\begin{proof}
From \eqref{eq:general-price}, $P_U\ge R_T(\Cold)/R_T(\pi)-C$ for the
fixed multiple $C$, so $P_U\to\infty$; since $\kappa^{\Cold}_U\le\max_{x\in\mathcal S}\kappa_x$ is a
property of the instance, this is $\kappa^\pi_U(t_0)\to0$ relative to $\kappa^{\Cold}_U$.
On the canonical instance the arms offered after round $\tau$ are
$b_t=Le_{t+1}$, $t>\tau$, and no arm played at or before round $\tau$ has a
component along $e_{t+1}$, so $\|b_t\|_{V_{\tau+1}^{-1}}=L/\sqrt\lambda$ and
$2\beta L/\sqrt\lambda>1=\Delta(b_t)$ under the hypothesis of
Proposition~\ref{prop:one-pass}; hence no arm offered after round $\tau$ is certified.
\end{proof}

\section{Proofs for the Catalog of Section~\ref{sec:catalog}}
\label{app:catalog}

Throughout this appendix the conventions \eqref{eq:conventions} are in force,
$\alpha:=1-\eps$, $m_0:=\lfloor\lambda/30\rfloor$, and ``LinUCB acts at round
$t$'' means that the action at $t$ is chosen by the LinUCB rule (every round
for $\Cold$; rounds $t>\tau$ for $\Hyb_\tau$). Ties in the LinUCB rule are
broken arbitrarily; every comparison we use is strict. For a policy $\pi$ and
a round $t$ we write $s(t)$, $m_j(t)$, $n_j(t)$ for the numbers of pulls of
$g$, $b$ in direction $e_j$, and $p$ in direction $e_j$ before round $t$
(all of them counted from the start, including expert-routed rounds), and we
drop the argument $t$ when no confusion arises. In the variant of
Proposition~\ref{prop:cat-gap} the roles of $g$ and $p$ are played by
$g_\eta$ and $p_{t,\eta}$ and $\alpha:=(1-\eps)\eta$; we prove everything for
general $\eta\in(0,1]$ and set $\eta=1$ outside that proposition.

\subsection{Events and their probabilities}

\begin{lemma}[Good events]
\label{lem:cat-events}
Fix a policy $\pi$ on any instance of the template. Let
$\mathcal E_{\mathrm{conf}}(\pi)$ be the event that
$\|\widehat\theta_t-\theta^\star\|_{V_t}\le\beta$ for every $t\in[T]$, where
$(V_t,c_t)$ are formed from all of $\pi$'s data before $t$. Let
$\mathcal E_{\mathrm{noise}}(\pi)$ be the event that for every direction
$j\in\{2,\dots,d\}$ and every $m\in[T]$, the sum $S_j^{(m)}$ of the noises on
the first $m$ pulls of $b$ in direction $e_j$ by $\pi$ satisfies
$|S_j^{(m)}|\le\sigma\sqrt{2m\ell_1}$. Then
$\Prb(\mathcal E_{\mathrm{conf}}(\pi))\ge1-\delta$ and
$\Prb(\mathcal E_{\mathrm{noise}}(\pi))\ge1-\delta$. Moreover, under
\eqref{eq:conventions}, $\beta\ge4\sigma\sqrt{2\ell_1}$.
\end{lemma}

\begin{proof}
All arm vectors of the template have Euclidean norm at most $L$: $\|g_\eta\|=\eta\le1$,
$\|b_t\|=L$, $\|p_{t,\eta}\|^2=\alpha^2+\gamma^2\le2\le L^2$. Also $\|\theta^\star\|=1$.
By Theorem~2 of \citet{abbasi2011improved}, with probability at least
$1-\delta$, simultaneously for all $t$,
$\|\widehat\theta_t-\theta^\star\|_{V_t}\le\sigma\sqrt{2\log\bigl(\det(V_t)^{1/2}\det(\lambda I)^{-1/2}/\delta\bigr)}+\sqrt\lambda$.
Since $V_t\preceq(\lambda+tL^2)I$, $\det(V_t)^{1/2}\det(\lambda I)^{-1/2}\le(1+tL^2/\lambda)^{d/2}$,
so the radius is at most
$\sigma\sqrt{d\log(1+TL^2/\lambda)+2\log(1/\delta)}+\sqrt\lambda\le\sigma\sqrt{d\,\ell_0}+\sqrt\lambda\le2\sqrt\lambda=\beta$,
using $\lambda\ge1$ and $\lambda\ge\sigma^2d\ell_0$. This proves the first claim.
The theorem of \citet{abbasi2011improved} only requires the noise to be
conditionally sub-Gaussian given the past and the current action, which holds
in the reward-table model, and it applies to the data of any policy, not only
to LinUCB's own.

For the second claim, in the reward-table model the noises on successive pulls
of $b$ in direction $e_j$ form an i.i.d.\ $\sigma$-sub-Gaussian sequence
$(\eta_{j,i})_{i\ge1}$ regardless of the policy, and $S_j^{(m)}=\sum_{i\le m}\eta_{j,i}$.
By Hoeffding's inequality for sub-Gaussian sums,
$\Prb(|S_j^{(m)}|>\sigma\sqrt{2m\ell_1})\le2e^{-\ell_1}=\delta/(dT)$. A union
bound over $d-1<d$ directions and $T$ values of $m$ gives probability at most
$\delta$ for the complement. Finally $\beta=2\sqrt\lambda\ge2\sqrt{8\sigma^2\ell_1}=4\sigma\sqrt{2\ell_1}$.
\end{proof}

\begin{lemma}[Confidence consequences]
\label{lem:cat-conf}
On $\mathcal E_{\mathrm{conf}}(\pi)$, for every $t$ and every $x\in\R^d$,
$|\langle\widehat\theta_t-\theta^\star,x\rangle|\le\beta\|x\|_{V_t^{-1}}$.
Hence the index $U_t(x):=\langle\widehat\theta_t,x\rangle+\beta\|x\|_{V_t^{-1}}$
satisfies $\mu(x)\le U_t(x)\le\mu(x)+2\beta\|x\|_{V_t^{-1}}$, where
$\mu(x)=\langle\theta^\star,x\rangle$.
\end{lemma}

\begin{proof}
Cauchy--Schwarz in the $V_t$ inner product:
$|\langle\widehat\theta_t-\theta^\star,x\rangle|\le\|\widehat\theta_t-\theta^\star\|_{V_t}\|x\|_{V_t^{-1}}\le\beta\|x\|_{V_t^{-1}}$.
\end{proof}

\subsection{The state of LinUCB on the template}

Every action vector of the template is $g_\eta$, some $b$ in direction $e_j$,
or some $p_\eta$ in direction $e_j$, and the probe is the only one that couples
two coordinates. The design matrix therefore has an arrowhead structure whose
inverse is explicit.

\begin{lemma}[Arrowhead inverse]
\label{lem:arrowhead}
Let $V=\lambda I+Se_1e_1^\top+\sum_{j=2}^{d}\bigl(M_je_je_j^\top+n_jq_jq_j^\top\bigr)$
with $S,M_j,n_j\ge0$ and $q_j=\alpha e_1+\gamma e_j$. Put
\[
D_j:=\lambda+M_j+n_j\gamma^2,\qquad u_j:=\alpha \gamma n_j,
\]
\[
Q:=\lambda+S+\alpha^2\sum_{j\ge2}n_j-\sum_{j\ge2}\frac{u_j^2}{D_j}
=\lambda+S+\alpha^2\sum_{j\ge2}\frac{n_j(\lambda+M_j)}{D_j}.
\]
Then $Q\ge\lambda$ and
\begin{align*}
(V^{-1})_{11}&=\frac1Q,&
(V^{-1})_{1j}&=-\frac{u_j}{D_jQ},\\
(V^{-1})_{jj}&=\frac1{D_j}+\frac{u_j^2}{D_j^2Q},&
(V^{-1})_{jj'}&=\frac{u_ju_{j'}}{D_jD_{j'}Q}\ \ (j\ne j'),
\end{align*}
for $j,j'\ge2$. In particular, if $n_j=0$ then
$(V^{-1})_{jj}=1/D_j$ and $(V^{-1})_{1j}=(V^{-1})_{jj'}=0$ for all $j'\ne j$; and
if $n_j\ge1$ and $M_j\ge1$ then $(V^{-1})_{jj}\le 1/D_j+1/M_j$.
\end{lemma}

\begin{proof}
Write $V=\begin{pmatrix}a&u^\top\\ u&D\end{pmatrix}$ with
$a=\lambda+S+\alpha^2\sum_jn_j$, $u=(u_2,\dots,u_d)^\top$, and
$D=\mathrm{diag}(D_2,\dots,D_d)$; this is the matrix in the statement because
$q_jq_j^\top=\alpha^2e_1e_1^\top+\alpha \gamma(e_1e_j^\top+e_je_1^\top)+\gamma^2e_je_j^\top$.
The Schur complement of $D$ is $a-u^\top D^{-1}u=Q$, and the block inverse
formula gives
$(V^{-1})_{11}=Q^{-1}$, $(V^{-1})_{1,\cdot}=-Q^{-1}u^\top D^{-1}$, and
$(V^{-1})_{\cdot,\cdot}=D^{-1}+D^{-1}uu^\top D^{-1}Q^{-1}$ on the lower block,
which are the displayed entries. The second expression for $Q$ follows from
$\alpha^2n_j-u_j^2/D_j=\alpha^2n_j(D_j-\gamma^2n_j)/D_j=\alpha^2n_j(\lambda+M_j)/D_j$,
and it shows $Q\ge\lambda$. If $n_j=0$ then $u_j=0$ and the special case
follows. For the last claim, $Q\ge\alpha^2n_j(\lambda+M_j)/D_j\ge\alpha^2n_jM_j/D_j$,
so $u_j^2/(D_j^2Q)\le\alpha^2\gamma^2n_j^2D_j/(D_j^2\alpha^2n_jM_j)=\gamma^2n_j/(D_jM_j)\le1/M_j$,
using $\gamma^2n_j\le D_j$.
\end{proof}

\begin{lemma}[State of a policy on the template]
\label{lem:cat-state}
Let $\pi$ be any policy on an instance of the template and $t$ a round. Then
$V_t$ is the matrix of Lemma~\ref{lem:arrowhead} with $S=\eta^2s$, $M_j=m_jL^2$,
$n_j$ the probe counts, and $q_j=p_\eta$'s direction vector. For a direction
$j\ge2$ with $n_j=0$:
\begin{enumerate}[itemsep=1pt,topsep=1pt,leftmargin=*]
\item[(a)] $\|e_j\|^2_{V_t^{-1}}=1/(\lambda+m_jL^2)$ and
$\widehat\theta_{t,j}=LS_j^{(m_j)}/(\lambda+m_jL^2)$, where $S_j^{(m_j)}$ is the
sum of the noises on the $m_j$ pulls of $b$ in direction $e_j$;
\item[(b)] $\|g_\eta\|^2_{V_t^{-1}}=\eta^2/Q\le1/\lambda$ and
$\|p_\eta\|^2_{V_t^{-1}}=\alpha^2/Q+\gamma^2/(\lambda+m_jL^2)\le2/\lambda$ for the probe
in direction $e_j$.
\end{enumerate}
For any direction with $m_j\ge1$, $\|e_j\|^2_{V_t^{-1}}\le2/(m_jL^2)$.
\end{lemma}

\begin{proof}
The form of $V_t$ is immediate from the definitions. For (a), Lemma~\ref{lem:arrowhead}
with $n_j=0$ gives $(V_t^{-1})_{jj}=1/D_j=1/(\lambda+m_jL^2)$ and that the $j$-th
row of $V_t^{-1}$ vanishes off the diagonal. The $j$-th coordinate of the
response vector $c_t$ is $L\sum r$ over the pulls of $b$ in direction $e_j$
plus $\gamma\sum r$ over the pulls of $p$ in that direction; the second sum is
empty, and since $\mu(b)=0$ the first is $LS_j^{(m_j)}$. Hence
$\widehat\theta_{t,j}=(V_t^{-1})_{jj}LS_j^{(m_j)}$. For (b),
$\|g_\eta\|^2_{V_t^{-1}}=\eta^2(V_t^{-1})_{11}=\eta^2/Q$ and
$\|p_\eta\|^2_{V_t^{-1}}=\alpha^2(V_t^{-1})_{11}+2\alpha \gamma(V_t^{-1})_{1j}+\gamma^2(V_t^{-1})_{jj}
=\alpha^2/Q+\gamma^2/D_j$, using $(V_t^{-1})_{1j}=0$; both bounds use $Q\ge\lambda$,
$D_j\ge\lambda$, $\alpha\le1$, $\gamma\le1$, $\eta\le1$. The last claim is the last
claim of Lemma~\ref{lem:arrowhead} with $M_j=m_jL^2\ge L^2\ge1$ together with
$1/D_j\le1/(m_jL^2)$, and it holds trivially when $n_j=0$.
\end{proof}

\subsection{Three reusable facts}

\begin{lemma}[Optimism trap]
\label{lem:cat-trap}
Let $\pi$ be a policy on an instance of the template and suppose LinUCB acts
for $\pi$ at round $t$. On $\mathcal E_{\mathrm{conf}}(\pi)\cap\mathcal E_{\mathrm{noise}}(\pi)$,
if the direction $j=j(t)$ satisfies $n_j(t)=0$ and $m_j(t)<m_0$, then
\[
U_t(b_t)\ \ge\ \tfrac34\,\frac{\beta L}{\sqrt{\lambda+m_j(t)L^2}}\ >\ \max\{U_t(g_\eta),U_t(p_{t,\eta})\},
\]
where the probe term is omitted if the probe is not offered at $t$. In
particular LinUCB plays $b_t$.
\end{lemma}

\begin{proof}
Write $m=m_j(t)$ and $D=\lambda+mL^2$. By Lemma~\ref{lem:cat-state}(a),
$U_t(b_t)=L\widehat\theta_{t,j}+\beta L\|e_j\|_{V_t^{-1}}=L^2S_j^{(m)}/D+\beta L/\sqrt D$.
On $\mathcal E_{\mathrm{noise}}$, $|L^2S_j^{(m)}/D|\le L^2\sigma\sqrt{2m\ell_1}/D\le L\sigma\sqrt{2\ell_1}/\sqrt D$,
since $L\sqrt m\le\sqrt D$; by Lemma~\ref{lem:cat-events} this is at most
$(\beta/4)L/\sqrt D$. This gives the first inequality. For the second, by
Lemma~\ref{lem:cat-conf} and Lemma~\ref{lem:cat-state}(b),
$U_t(g_\eta)\le\eta+2\beta\eta/\sqrt Q\le1+2\beta/\sqrt\lambda$ and
$U_t(p_{t,\eta})\le(1-\eps)\eta+2\beta\sqrt{2/\lambda}\le1+2\sqrt2\,\beta/\sqrt\lambda$.
Since $\beta=2\sqrt\lambda$, $1\le\beta/(2\sqrt\lambda)$, so both indices are at
most $(\tfrac12+2\sqrt2)\beta/\sqrt\lambda$. It remains to check
$\tfrac34L/\sqrt D>(\tfrac12+2\sqrt2)/\sqrt\lambda$, i.e.\
$D<\lambda L^2\,(3/4)^2/(\tfrac12+2\sqrt2)^2$. Numerically
$(3/4)^2/(\tfrac12+2\sqrt2)^2=0.5625/11.0784>0.0507$. Since $m\le m_0-1<\lambda/30$,
$D<\lambda+\lambda L^2/30=\lambda L^2(1/L^2+1/30)\le\lambda L^2(1/64+1/30)<0.0491\,\lambda L^2$,
using $L\ge8$. The inequality holds with room to spare.
\end{proof}

\begin{lemma}[Bad-pull cap]
\label{lem:cat-cap}
Let $\pi$ be a policy on an instance of the template. On
$\mathcal E_{\mathrm{conf}}(\pi)$, at any round $t$ at which LinUCB acts for
$\pi$ and plays $b_t$ with direction $j$, $m_j(t)\le8\beta^2/\eta^2$; if
moreover $n_j(t)=0$, then $m_j(t)\le4\beta^2/\eta^2$. Consequently the number
of rounds at which LinUCB acting for $\pi$ plays $b$ in direction $e_j$ is at
most $8\beta^2/\eta^2+1$, and at most $4\beta^2/\eta^2+1$ if no probe in
direction $e_j$ is ever played by $\pi$.
\end{lemma}

\begin{proof}
By Lemma~\ref{lem:width-limited}, a bad arm with gap $\eta$ is played only if
$\eta\le2\beta L\|e_j\|_{V_t^{-1}}$, i.e.\ $\|e_j\|^2_{V_t^{-1}}\ge\eta^2/(4\beta^2L^2)$.
If $m_j(t)\ge1$, Lemma~\ref{lem:cat-state} gives $\|e_j\|^2_{V_t^{-1}}\le2/(m_j(t)L^2)$,
so $m_j(t)\le8\beta^2/\eta^2$; with $n_j(t)=0$ the sharper bound
$\|e_j\|^2_{V_t^{-1}}=1/(\lambda+m_j(t)L^2)\le1/(m_j(t)L^2)$ gives $m_j(t)\le4\beta^2/\eta^2$.
Since $m_j(\cdot)$ increases by one at each such round, the number of rounds
at which $b$ in direction $e_j$ is played with $m_j(t)\le M$ is at most $M+1$.
\end{proof}

\begin{lemma}[Expert routing]
\label{lem:expert-routes}
Let $\eps L<\gamma$ and let $\Exp_w$ be the frozen linear policy with
$\widehat\theta_{\Exp}=e_1+w\sum_{j\ge2}e_j$ on the instance with $g_\eta,b_t,p_{t,\eta}$.
\begin{enumerate}[itemsep=1pt,topsep=1pt,leftmargin=*]
\item[(a)] If $w\in(\eta/L,(1-\eps)\eta/(L-\gamma))$, a nonempty interval, then
$\Exp_w$ plays $p_{t,\eta}$ whenever the probe is offered and $b_t$ otherwise.
\item[(b)] If $w\in(\eps\eta/\gamma,\eta/L)$, a nonempty interval, then $\Exp_w$ plays
$p_{t,\eta}$ whenever the probe is offered and $g_\eta$ otherwise.
\end{enumerate}
All three comparisons are strict.
\end{lemma}

\begin{proof}
The scores are $\langle\widehat\theta_{\Exp},g_\eta\rangle=\eta$,
$\langle\widehat\theta_{\Exp},b_t\rangle=wL$, and
$\langle\widehat\theta_{\Exp},p_{t,\eta}\rangle=(1-\eps)\eta+w\gamma$. The probe beats
$g_\eta$ iff $w\gamma>\eps\eta$ and beats $b_t$ iff $w<(1-\eps)\eta/(L-\gamma)$; $b_t$
beats $g_\eta$ iff $w>\eta/L$. The interval in (a) is nonempty iff
$(1-\eps)\eta L>\eta(L-\gamma)$, i.e.\ $\gamma>\eps L$, and every $w$ in it satisfies
$w>\eta/L>\eps\eta/\gamma$ because $\gamma>\eps L$. The interval in (b) is nonempty iff
$\eps\eta/\gamma<\eta/L$, again $\gamma>\eps L$; for $w$ in it, $w<\eta/L<(1-\eps)\eta/(L-\gamma)$.
\end{proof}

\begin{lemma}[Widths at handoff after $n$ probes in each direction]
\label{lem:cat-handoff}
Let $V=\lambda I+n\sum_{j=2}^{K+1}q_jq_j^\top$ with $q_j=\alpha e_1+\gamma e_j$, $\alpha>0$,
$n\ge1$. Then for every $j\in\{2,\dots,K+1\}$,
\[
\|e_j\|^2_{V^{-1}}\ \le\ \frac1{n\gamma^2}+\frac1{K\lambda},
\qquad\text{and}\qquad
Q=\lambda+\frac{\alpha^2Kn\lambda}{\lambda+n\gamma^2}\ \ge\ \lambda\Bigl(1+\frac{\alpha^2K}{2\gamma^2}\Bigr)\ \text{ if } n\gamma^2\ge\lambda .
\]
\end{lemma}

\begin{proof}
Apply Lemma~\ref{lem:arrowhead} with $S=0$, $M_j=0$, $n_j=n$: $D_j=D:=\lambda+n\gamma^2$,
$u_j=\alpha \gamma n$, and $Q=\lambda+\alpha^2Kn\lambda/D$. Then
$(V^{-1})_{jj}=1/D+(\alpha \gamma n)^2/(D^2Q)$. The first term is at most $1/(n\gamma^2)$.
For the second, $Q\ge\alpha^2Kn\lambda/D$, so
$(\alpha \gamma n)^2/(D^2Q)\le(\alpha \gamma n)^2D/(D^2\alpha^2Kn\lambda)=\gamma^2n/(DK\lambda)\le1/(K\lambda)$,
using $\gamma^2n\le D$. The bound on $Q$ follows from $D\le2n\gamma^2$ when $n\gamma^2\ge\lambda$.
\end{proof}

\subsection{Proof of Lemma~\ref{lem:no-free-coverage}}

Fix $u$. At a round $t\in I$ with $\langle x_{t,\pi_t},u\rangle=0$ the
contribution to the coverage is zero. At a round with
$\langle x_{t,\pi_t},u\rangle\ne0$, the definition \eqref{eq:price} gives
$\Delta_{t,\pi_t}/\langle x_{t,\pi_t},u\rangle^2\ge\kappa_t(u)\ge\kappa(u)$, i.e.\
$\langle x_{t,\pi_t},u\rangle^2\le\Delta_{t,\pi_t}/\kappa(u)$ (if $\kappa(u)=0$ the bound
is vacuous). Summing over $I$ gives the first claim. For the second, let
$u$ be a unit vector and $V=V_{\tau+1}=\lambda I+M_\tau$. By Cauchy--Schwarz in
the $V$ inner product,
$1=(u^\top u)^2=(u^\top V^{1/2}V^{-1/2}u)^2\le(u^\top Vu)(u^\top V^{-1}u)$, so
$\|u\|^2_{V^{-1}}\ge1/(\lambda+u^\top M_\tau u)$, and
$u^\top M_\tau u=\sum_{t\le\tau}\langle x_{t,A_t},u\rangle^2\le R_{[1,\tau]}(\Exp)/\kappa(u)$
by the first claim applied to the expert on $[1,\tau]$. \qed

\subsection{Proof of Theorem~\ref{thm:price}}

We work on $\mathcal E_{\mathrm{conf}}(\Cold)\cap\mathcal E_{\mathrm{noise}}(\Cold)\cap\mathcal E_{\mathrm{conf}}(\Hyb_{\tau_0})$,
which has probability at least $1-3\delta$ by Lemma~\ref{lem:cat-events}. Set
$\eta=1$.

\paragraph{(i) The expert.}
By Lemma~\ref{lem:expert-routes}(a), $\Exp_w$ plays $p_t$ at every $t\le\tau_0$
(gap $\eps$) and $b_t$ at every $t>\tau_0$ (gap $1$), so $R_T(\Exp)=\eps Kn+(T-Kn)$.

\paragraph{(ii) Cold LinUCB, lower bound.}
We claim that at every round $t\in[T]$ with $m_{j(t)}(t)<m_0$, cold LinUCB
plays $b_t$. We argue by induction on $t$. Suppose the claim holds for all
rounds before $t$ and let $j=j(t)$ with $m_j(t)<m_0$. If cold LinUCB had
played the probe in direction $e_j$ at some earlier round $t'$ (so $j(t')=j$),
then $m_j(t')\le m_j(t)<m_0$ and the claim at $t'$ says it played $b_{t'}$,
a contradiction; hence $n_j(t)=0$. Lemma~\ref{lem:cat-trap} now gives that
LinUCB plays $b_t$. (The base case $t=1$ is the same argument with no earlier
rounds.) Direction $e_j$ is offered at exactly $n$ rounds in $[1,Kn]$, and at
each of the first $\min\{n,m_0\}$ of them $m_j<m_0$, so cold LinUCB plays the
bad arm, paying $1$ each time. With $n\ge m_0$ this gives $R_{[1,Kn]}(\Cold)\ge Km_0$,
and $A_{\tau_0}=R_{[1,Kn]}(\Cold)-R_{[1,Kn]}(\Exp)\ge Km_0-\eps Kn$ by (i).

\paragraph{(iii) Cold LinUCB, upper bound.}
Cold LinUCB acts at every round. By Lemma~\ref{lem:cat-cap} it plays $b$ in
each of the $K$ directions at most $8\beta^2+1$ times over the horizon, each
costing $1$. Probes exist only at rounds $t\le Kn$ and cost $\eps$ each, so
probe pulls cost at most $\eps Kn$. Pulls of $g$ cost nothing. Hence
$R_T(\Cold)\le K(8\beta^2+1)+\eps Kn$.

\paragraph{(iv) The hybrid.}
During $t\le\tau_0$ the hybrid plays the expert's actions, which by (i) are
the probes, each direction exactly $n$ times. Hence
$V_{\tau_0+1}=\lambda I+n\sum_{j=2}^{K+1}p^{(j)}p^{(j)\top}$ with $p^{(j)}=\alpha e_1+\gamma e_j$,
and Lemma~\ref{lem:cat-handoff} gives, for every bad direction,
$\|e_j\|^2_{V_{\tau_0+1}^{-1}}\le1/(n\gamma^2)+1/(K\lambda)$. With $n\gamma^2\ge36\lambda L^2=9\beta^2L^2$
and $K\lambda\ge36\lambda L^2=9\beta^2L^2$ this is at most $2/(9\beta^2L^2)$, so
$2\beta L\|e_j\|_{V_{\tau_0+1}^{-1}}\le2\sqrt2/3<1$. Now fix $t>\tau_0$. The
hybrid's design matrix satisfies $V_t\succeq V_{\tau_0+1}$, so by
Lemma~\ref{lem:loewner} $\|e_{j(t)}\|_{V_t^{-1}}\le\|e_{j(t)}\|_{V_{\tau_0+1}^{-1}}$,
and by Lemma~\ref{lem:cat-conf}, $U_t(b_t)\le0+2\beta L\|e_{j(t)}\|_{V_t^{-1}}<1\le U_t(g)$.
Since $\mathcal A_t=\{g,b_t\}$, LinUCB plays $g$, with gap $0$. Therefore
$R_{[\tau_0+1,T]}(\Hyb_{\tau_0})=0$, $R_T(\Hyb_{\tau_0})=R_{[1,\tau_0]}(\Exp)=\eps Kn$,
$I_{\tau_0}=R_{[Kn+1,T]}(\Cold)-0\ge0$, and $B_{\tau_0}=(T-Kn)-0$.

\paragraph{The ratios.}
From (ii), (iii), (iv): $R_T(\Cold)/R_T(\Hyb_{\tau_0})\ge Km_0/(\eps Kn)=m_0/(\eps n)$
and $\le(K(8\beta^2+1)+\eps Kn)/(\eps Kn)=1+(8\beta^2+1)/(\eps n)$; the
expert ratio is (i) over (iv). At $n=\lceil36\lambda L^2/\gamma^2\rceil$, since
$36\lambda L^2/\gamma^2\ge36\cdot60\cdot64>1$ we have $n\le37\lambda L^2/\gamma^2$, and
$m_0\ge\lambda/30-1\ge\lambda/60$ because $\lambda\ge60$; hence
$m_0/(\eps n)\ge(\lambda/60)\gamma^2/(37\eps\lambda L^2)=\gamma^2/(2220\,\eps L^2)$. For the
upper bound, $8\beta^2+1=32\lambda+1$ and $n\ge36\lambda L^2/\gamma^2$ give
$(8\beta^2+1)/(\eps n)\le(32\lambda+1)\gamma^2/(36\eps\lambda L^2)\le \gamma^2/(\eps L^2)$
since $\lambda\ge1$. \qed

\subsection{Proof of Corollary~\ref{cor:cat-dear}}

Let $\eps\ge \gamma^2/L^2$ and $n\ge36\lambda L^2/\gamma^2$, so that
$\eps n\ge(\gamma^2/L^2)(36\lambda L^2/\gamma^2)=36\lambda$. With $\beta=2\sqrt\lambda$
from \eqref{eq:conventions} we have $8\beta^2+1=32\lambda+1$, hence
\[
\frac{8\beta^2+1}{\eps n}\ \le\ \frac{32\lambda+1}{36\lambda}
\ =\ \frac89+\frac{1}{36\lambda}\ \le\ \frac89+\frac{1}{2160}\ <\ 1 ,
\]
using $\lambda\ge60$. The upper bound of \eqref{eq:price-ratio} then gives
$R_T(\Cold)/R_T(\Hyb_{\tau_0})\le1+(8\beta^2+1)/(\eps n)<2$. \qed

\subsection{Proof of Proposition~\ref{prop:cat-tail-saving}}
\label{app:proof-tail-saving}

\paragraph{Events.}
All arm vectors of $\mathcal J$ have norm at most $L$: $\|b'_{j'}\|_2^2=1+L'^2=145\le L^2$
since $L\ge13$, and the others are as in Lemma~\ref{lem:cat-events}. Let
$\mathcal E'_{\mathrm{conf}}(\Cold)$ be the event that
$\|\widehat\theta_t-\theta^\star\|_{V_t}\le\tfrac54\sqrt\lambda$ for every $t$
on cold LinUCB's trajectory. By Theorem~2 of \citet{abbasi2011improved} and
the computation in the proof of Lemma~\ref{lem:cat-events}, the
self-normalized radius is at most $\sigma\sqrt{d\ell_0}+\sqrt\lambda\le\tfrac14\sqrt\lambda+\sqrt\lambda$
under $\lambda\ge16\sigma^2d\ell_0$, so
$\Prb(\mathcal E'_{\mathrm{conf}}(\Cold))\ge1-\delta$, and
$\mathcal E'_{\mathrm{conf}}(\Cold)\subseteq\mathcal E_{\mathrm{conf}}(\Cold)$
since $\tfrac54\sqrt\lambda\le\beta=2\sqrt\lambda$. We work on
$\mathcal E'_{\mathrm{conf}}(\Cold)\cap\mathcal E_{\mathrm{noise}}(\Cold)\cap\mathcal E_{\mathrm{conf}}(\Hyb_{\tau_0})$,
of probability at least $1-3\delta$ (Lemma~\ref{lem:cat-events}; the noise
event concerns the pulls of $b$ in the directions $e_j$, $j\in J$, and its
proof is unchanged). On $\mathcal E'_{\mathrm{conf}}(\Cold)$, Cauchy--Schwarz
gives $|\langle\widehat\theta_t-\theta^\star,x\rangle|\le\tfrac54\sqrt\lambda\|x\|_{V_t^{-1}}=\tfrac58\beta\|x\|_{V_t^{-1}}$
for every $x$, hence for the index $U_t(x)=\langle\widehat\theta_t,x\rangle+\beta\|x\|_{V_t^{-1}}$,
\begin{equation}
\label{eq:sharp-index}
\mu(x)+\tfrac38\beta\|x\|_{V_t^{-1}}\ \le\ U_t(x)\ \le\ \mu(x)+\tfrac{13}8\beta\|x\|_{V_t^{-1}} .
\end{equation}

\paragraph{(i) The expert.}
The expert plays $p_{j(t)}$ (gap $\eps$) at every $t\le\tau_0$ and $b_{j(t)}$
(gap $1$) at every $t>\tau_0$, so $R_T(\Exp)=\eps Kn+(T-Kn)$.

\paragraph{(ii) Cold LinUCB in the window.}
We show by induction on $t\le\tau_0$ that cold LinUCB plays $b'_{j'(t)}$ at
round $t$. Suppose this holds before $t$; then no probe and no bad arm has
been played, so for every $j\in J$ no played vector has a component along
$e_j$, $V_te_j=\lambda e_j$, and $V_t^{-1}e_j=e_j/\lambda$. Three width bounds
follow. First, $\|g\|^2_{V_t^{-1}}=(V_t^{-1})_{11}\le1/\lambda$ because
$V_t\succeq\lambda I$. Second, with $j=j(t)$,
$\|p_j\|^2_{V_t^{-1}}=(1-\eps)^2(V_t^{-1})_{11}+2(1-\eps)\gamma\,e_1^\top V_t^{-1}e_j+\gamma^2(V_t^{-1})_{jj}
=(1-\eps)^2(V_t^{-1})_{11}+\gamma^2/\lambda\le(1+\gamma^2)/\lambda\le2/\lambda$, using
$e_1^\top V_t^{-1}e_j=e_1^\top e_j/\lambda=0$. Third, with $j'=j'(t)$ and $m$
the number of earlier pulls of $b'_{j'}$: the only played vector with a
component along $e_{j'}$ is $b'_{j'}$ itself, so
$e_{j'}^\top V_te_{j'}=\lambda+mL'^2$, and Cauchy--Schwarz in the $V_t$ inner
product gives
$\|b'_{j'}\|^2_{V_t^{-1}}\ge\langle b'_{j'},e_{j'}\rangle^2/(e_{j'}^\top V_te_{j'})=L'^2/(\lambda+mL'^2)$.
Direction $j'$ is offered at $\lceil\tau_0/K'\rceil$ rounds in the window, so
$m\le\tau_0/K'\le\lambda/L'^2$ by the condition on $K'$, and
$\|b'_{j'}\|^2_{V_t^{-1}}\ge L'^2/(2\lambda)$. Now apply
\eqref{eq:sharp-index} with $\beta=2\sqrt\lambda$:
\begin{gather*}
U_t(b'_{j'})\ \ge\ 1+\tfrac38\cdot2\sqrt\lambda\cdot\frac{L'}{\sqrt{2\lambda}}=1+\frac{3L'}{4\sqrt2}>7.3,
\\
U_t(g)\ \le\ 1+\tfrac{13}8\cdot2\sqrt\lambda\cdot\frac1{\sqrt\lambda}=4.25,
\qquad
U_t(p_j)\ \le\ 1-\eps+\tfrac{13}8\cdot2\sqrt\lambda\cdot\frac{\sqrt2}{\sqrt\lambda}<5.6,
\end{gather*}
using $L'=12$. So cold LinUCB plays $b'_{j'(t)}$, whose gap is $0$. The base
case $t=1$ is the same computation with $m=0$. Hence $R_{[1,\tau_0]}(\Cold)=0$,
and at round $\tau_0+1$ every direction $e_j$, $j\in J$, is untouched:
$m_j=n_j=0$.

\paragraph{(ii) Cold LinUCB on the tail.}
For $t>\tau_0$ the action set is $\{g,b_{j(t)}\}$ and the directions $e_j$,
$j\in J$, are touched only by the bad arms, so for each of them $e_j$ is an
eigenvector of $V_t$ with eigenvalue $\lambda+m_j(t)L^2$, the $j$-th
coordinate of the response vector is $LS_j^{(m_j(t))}$, and the conclusions of
Lemma~\ref{lem:cat-state}(a) hold verbatim; Lemma~\ref{lem:cat-state}(b)'s
bound $U_t(g)\le1+2\beta/\sqrt\lambda$ holds on $\mathcal E_{\mathrm{conf}}(\Cold)$
by Lemma~\ref{lem:cat-conf} and $\|g\|_{V_t^{-1}}\le1/\sqrt\lambda$. These are
the only properties of the template that the proofs of
Lemmas~\ref{lem:cat-trap} and~\ref{lem:cat-cap} use for a direction with
$n_j=0$, so both lemmas apply to cold LinUCB on $\mathcal J$ at every tail
round (with no probe term). The induction of Theorem~\ref{thm:price}(ii) now
shows that cold LinUCB plays $b_t$ at every tail round with $m_{j(t)}(t)<m_0$;
since $T-\tau_0\ge Km_0$ and the tail directions cycle, each direction is
offered at least $m_0$ times on the tail, giving $R_{[\tau_0+1,T]}(\Cold)\ge Km_0$.
Lemma~\ref{lem:cat-cap} with $n_j=0$ bounds the pulls of $b$ in each
direction by $4\beta^2+1$, and $g$ costs nothing, so
$R_{[\tau_0+1,T]}(\Cold)\le K(4\beta^2+1)=K(16\lambda+1)$.

\paragraph{(iii) The hybrid.}
During $t\le\tau_0$ the hybrid plays the expert's actions, which by (i) are
the probes, each direction $j\in J$ exactly $n$ times; no distractor and no
bad arm is played. Hence $V_{\tau_0+1}=\lambda I+n\sum_{j\in J}p_jp_j^\top$,
which is the matrix of Theorem~\ref{thm:price}(iv), and the argument there
(Lemma~\ref{lem:cat-handoff}, Lemma~\ref{lem:loewner}, Lemma~\ref{lem:cat-conf}
on $\mathcal E_{\mathrm{conf}}(\Hyb_{\tau_0})$) shows that
$U_t(b_t)<1\le U_t(g)$ at every $t>\tau_0$, so the hybrid plays $g$ and
$R_{[\tau_0+1,T]}(\Hyb_{\tau_0})=0$. Therefore $R_T(\Hyb_{\tau_0})=\eps Kn$.

\paragraph{The terms and the ratios.}
$A_{\tau_0}=R_{[1,\tau_0]}(\Cold)-R_{[1,\tau_0]}(\Exp)=0-\eps Kn$;
$I_{\tau_0}=R_{[\tau_0+1,T]}(\Cold)-0\in[Km_0,K(16\lambda+1)]$;
$B_{\tau_0}=(T-Kn)-0$. Since $R_T(\Cold)=R_{[\tau_0+1,T]}(\Cold)$ and
$R_T(\Hyb_{\tau_0})=\eps Kn$, the ratio lies between $m_0/(\eps n)$ and
$(16\lambda+1)/(\eps n)$. At $n=\lceil36\lambda L^2/\gamma^2\rceil$ the lower bound
is at least $\gamma^2/(2220\,\eps L^2)$ exactly as in the proof of
Theorem~\ref{thm:price}, and $(16\lambda+1)/(\eps n)\le(16\lambda+1)\gamma^2/(36\eps\lambda L^2)\le \gamma^2/(2\eps L^2)$
since $\lambda\ge1$. \qed

\subsection{Proof of Proposition~\ref{prop:cat-mimic}}

Work on $\mathcal E_{\mathrm{conf}}(\Cold)\cap\mathcal E_{\mathrm{noise}}(\Cold)$
(probability $\ge1-2\delta$). There are no probes, so $n_j\equiv0$ for every
policy.

\emph{(i)} Let $t\le Km_0$ and $j=j(t)$. Direction $e_j$ was offered at
$\lceil t/K\rceil-1\le m_0-1$ earlier rounds, so $m_j(t)\le m_0-1<m_0$ for cold
LinUCB, and Lemma~\ref{lem:cat-trap} (with the probe term absent) says cold
LinUCB plays $b_t$. Hence cold LinUCB plays $b_t$ at every $t\le Km_0$, and
so does $\Exp$. Fix $\tau\le Km_0$. The two policies $\Cold$ and $\Hyb_\tau$
play the same action vectors in the same order during $[1,\tau]$; in the
reward-table model they therefore observe the same rewards, so their states
$(V_{\tau+1},c_{\tau+1})$ coincide. From round $\tau+1$ on both run LinUCB
from that state on the same action sets and the same reward table, so they
play the same actions at every round. Thus $R_I(\Cold)=R_I(\Hyb_\tau)$ for every
interval $I$, $A_\tau=R_{[1,\tau]}(\Cold)-R_{[1,\tau]}(\Exp)=\tau-\tau=0$, and $I_\tau=0$.

\emph{(ii)} For any $\tau$, $R_T(\Hyb_\tau)\ge R_{[1,\tau]}(\Exp)=\tau$. By
Lemma~\ref{lem:cat-cap} with $n_j\equiv0$, cold LinUCB plays each bad
direction at most $4\beta^2+1$ times, so $R_T(\Cold)\le K(4\beta^2+1)=K(16\lambda+1)$.
For $\tau>Km_0$, $R_T(\Cold)/R_T(\Hyb_\tau)\le K(16\lambda+1)/(Km_0)\le(16\lambda+1)/(\lambda/60)=960+60/\lambda\le961$. \qed

\subsection{Proof of Proposition~\ref{prop:cat-tail}}

Work on the event of Theorem~\ref{thm:price}. By Lemma~\ref{lem:expert-routes}(b),
$\Exp_w$ plays every probe during $[1,\tau_0]$ and $g$ afterwards, so
$R_T(\Exp)=\eps Kn$. The hybrid's routed rounds are the same probes as in
Theorem~\ref{thm:price}(iv), so its state at handoff and its subsequent play
are identical to those established there: $R_T(\Hyb_{\tau_0})=\eps Kn$ and
$R_{[\tau_0+1,T]}(\Hyb_{\tau_0})=0$. Hence $B_{\tau_0}=0-0=0$. Cold LinUCB does not
depend on the expert, so $A_{\tau_0}$ and $I_{\tau_0}$ are as in
Theorem~\ref{thm:price}. \qed

\subsection{Proof of Proposition~\ref{prop:cat-wrong}}

Work on $\mathcal E_{\mathrm{conf}}(\Cold)\cap\mathcal E_{\mathrm{conf}}(\Hyb_{\tau_0})\cap\mathcal E_{\mathrm{noise}}(\Hyb_{\tau_0})$
(probability $\ge1-3\delta$). Here $d=2K+1$; coordinates $2,\dots,K+1$ are the
bad directions and $K+2,\dots,2K+1$ the nuisance directions. The probe $p'_t$
couples $e_1$ with a nuisance coordinate only, so Lemma~\ref{lem:arrowhead}
applies with the nuisance coordinates playing the role of the probed $j$'s
and the bad coordinates having $n_j=0$ for every policy at every round.

\emph{Cold LinUCB.} Since $n_j\equiv0$ on bad coordinates,
Lemma~\ref{lem:cat-cap} gives at most $4\beta^2+1$ bad pulls in each of the $K$
directions, costing $1$ each; probe pulls occur only at $t\le Kn$ and cost
$\eps$ each; so $R_T(\Cold)\le K(4\beta^2+1)+\eps Kn$.

\emph{The hybrid.} During $[1,\tau_0]$ it plays the $Kn$ probes, paying $\eps Kn$.
At handoff, $V_{\tau_0+1}=\lambda I+n\sum_{i}q_iq_i^\top$ over the $K$ nuisance
directions with $q_i=\alpha e_1+\gamma e_{K+1+i}$, and the bad coordinates are
untouched: by Lemma~\ref{lem:arrowhead} (with $M_j=n_j=0$ on them),
$\|e_j\|^2_{V_{\tau_0+1}^{-1}}=1/\lambda$ for every bad direction and the estimate
has zero $j$-th coordinate. By Lemma~\ref{lem:cat-handoff} applied to the
nuisance block, $Q\ge\lambda(1+\alpha^2K/(2\gamma^2))$ once $n\gamma^2\ge\lambda$, which holds
since $n\ge36\lambda L^2/\gamma^2$. As $\alpha=1-\eps\ge1-1/L\ge7/8$ and $\gamma\le1$,
$Q\ge\lambda(1+49K/128)\ge16\lambda=4\beta^2$ because $K\ge36L^2\ge2304$. Hence at
every $t>\tau_0$, since $Q$ only grows, $U_t(g)\le1+2\beta/\sqrt Q\le2$ by
Lemma~\ref{lem:cat-conf}. On the other hand, for $t>\tau_0$ and $j=j(t)$ with
$m_j(t)=m$, Lemma~\ref{lem:cat-state}(a) and $\mathcal E_{\mathrm{noise}}(\Hyb_{\tau_0})$
give, exactly as in the proof of Lemma~\ref{lem:cat-trap},
$U_t(b_t)\ge\tfrac34\beta L/\sqrt{\lambda+mL^2}$, which exceeds $2$ as soon as
$\lambda+mL^2<\tfrac9{64}\beta^2L^2=\tfrac9{16}\lambda L^2$, i.e.\ $m<\tfrac9{16}\lambda-\lambda/L^2$.
With $L\ge8$, $\tfrac9{16}\lambda-\lambda/L^2\ge\tfrac{35}{64}\lambda>\lambda/2$, so this holds
whenever $m\le\lfloor\lambda/2\rfloor-1$. Thus at every round $t>\tau_0$ at which
$m_{j(t)}(t)<\lfloor\lambda/2\rfloor$ the hybrid plays the bad arm. Each direction
is offered at least $\lfloor(T-\tau_0)/K\rfloor\ge\lfloor\lambda/2\rfloor$ times after
$\tau_0$, so each is pulled at least $\lfloor\lambda/2\rfloor$ times, and
$R_{[\tau_0+1,T]}(\Hyb_{\tau_0})\ge K\lfloor\lambda/2\rfloor$. Hence
$R_T(\Hyb_{\tau_0})\ge\eps Kn+K\lfloor\lambda/2\rfloor$.

\emph{The ratio and $I_{\tau_0}$.} Since $\lambda\ge60$, $\lfloor\lambda/2\rfloor\ge\lambda/4$,
and for $x\ge0$, $(a+x)/(b+x)\le\max\{a/b,1\}$; with $a=K(16\lambda+1)$,
$b=K\lambda/4$, $x=\eps Kn$ this gives $R_T(\Cold)/R_T(\Hyb_{\tau_0})\le64+4/\lambda\le65$.
Finally $I_{\tau_0}=R_{[\tau_0+1,T]}(\Cold)-R_{[\tau_0+1,T]}(\Hyb_{\tau_0})\le R_{[\tau_0+1,T]}(\Cold)-K\lfloor\lambda/2\rfloor$.
The expert's Gram matrix $n\sum_iq_iq_i^\top$ has range in
$\mathrm{span}\{e_1,e_{K+2},\dots,e_{2K+1}\}$, which is orthogonal to the span of
the bad directions, so the principal angle is $90^\circ$. \qed

\subsection{Proof of Proposition~\ref{prop:cat-cheap}}

On $\mathcal E_{\mathrm{conf}}(\Cold)$, Lemma~\ref{lem:cat-cap} with $n_j\equiv0$
gives $R_T(\Cold)\le K(4\beta^2+1)=K(16\lambda+1)$; since $R_T(\pi)\ge0$ for every
$\pi$, $(R_T(\Cold)-R_T(\pi))/T\le K(16\lambda+1)/T$, and with $\lambda=\sigma^2d\ell_0$
and $d=K+1$ this is $O(\sigma^2K^2\ell_0/T)$. \qed

\subsection{Proof of Proposition~\ref{prop:cat-gap}}

The proof of Theorem~\ref{thm:price} goes through with $\eta$ carried along;
we indicate the changes. Work on the same three events. (i)~By
Lemma~\ref{lem:expert-routes}(a), $R_T(\Exp)=\eps\eta Kn+\eta(T-Kn)$.
(ii)~Lemma~\ref{lem:cat-trap} was proved for general $\eta$, and the induction
is unchanged; each forced bad pull now costs $\eta$, so
$R_{[1,Kn]}(\Cold)\ge\eta Km_0$. (iii)~Lemma~\ref{lem:cat-cap} gives at most
$8\beta^2/\eta^2+1$ bad pulls in each direction at cost $\eta$, and probe pulls
cost $\eps\eta$, so $R_T(\Cold)\le\eta K(8\beta^2/\eta^2+1)+\eps\eta Kn$.
(iv)~At handoff, Lemma~\ref{lem:cat-handoff} with $\alpha=(1-\eps)\eta>0$ gives
$\|e_j\|^2_{V_{\tau_0+1}^{-1}}\le1/(n\gamma^2)+1/(K\lambda)\le2\eta^2/(9\beta^2L^2)$ under
$n\gamma^2\ge9\beta^2L^2/\eta^2$ and $K\lambda\ge9\beta^2L^2/\eta^2$, i.e.\ $n\ge36\lambda L^2/(\eta^2\gamma^2)$
and $K\ge36L^2/\eta^2$. Then $U_t(b_t)\le2\beta L\|e_j\|_{V_t^{-1}}\le(2\sqrt2/3)\eta<\eta\le U_t(g_\eta)$
for all $t>\tau_0$ by Lemmas~\ref{lem:loewner} and~\ref{lem:cat-conf}, so the
hybrid plays $g_\eta$ and $R_T(\Hyb_{\tau_0})=\eps\eta Kn$. The ratio bounds
follow: the lower bound is $\eta Km_0/(\eps\eta Kn)=m_0/(\eps n)$, and at
$n=\lceil36\lambda L^2/(\eta^2\gamma^2)\rceil\le37\lambda L^2/(\eta^2\gamma^2)$ this is at least
$(\lambda/60)\eta^2\gamma^2/(37\eps\lambda L^2)=\eta^2\gamma^2/(2220\,\eps L^2)$; the upper bound
is $1+(8\beta^2/\eta^2+1)/(\eps n)\le1+(32\lambda+\eta^2)\gamma^2/(36\eps\lambda L^2)\le1+\gamma^2/(\eps L^2)$. \qed

\subsection{Proof of Corollary~\ref{cor:cat-oracle}}

On the event of Proposition~\ref{prop:one-pass}, the hybrid plays $b_t$ at
every $t>\tau$, so $R_{[\tau+1,T]}(\Hyb_\tau)=T-\tau$, while the routed rounds
cost nothing because the expert plays the optimal arm; hence
$R_T(\Hyb_\tau)=T-\tau$ and $R_T(\Exp)=0$. Cold LinUCB also plays $b_t$ at
every round by Proposition~\ref{prop:one-pass} with $\tau=0$, so
$R_{[1,\tau]}(\Cold)=\tau$ and $R_{[\tau+1,T]}(\Cold)=T-\tau$. Therefore
$A_\tau=\tau-0$, $I_\tau=(T-\tau)-(T-\tau)=0$, and $B_\tau=0-(T-\tau)$. \qed

\subsection{Remarks on the constants}

The constants $36$, $30$, $65$, $961$, $2220$ come from three places: the
factor $\tfrac34$ by which noise can reduce a bad arm's index below its bonus,
the factor $\tfrac12+2\sqrt2$ by which the good arm's or the probe's index can
exceed its mean at a cold start, and the choice $\beta=2\sqrt\lambda$. None of
them is tight. What matters in Theorem~\ref{thm:price} is the structure of
\eqref{eq:price-ratio}: the lower bound is (forced cold pulls per direction)
divided by (probe cost per direction), the upper bound is (maximal cold pulls
per direction) divided by the same, and both numerators are independent of
$\eps$, $\gamma$, and $K$. The gap between $m_0\asymp\lambda$ and $8\beta^2\asymp\lambda$
is a constant factor, and that is why the multiplier is pinned to the price
ratio up to constants. The same structure with $\eta$ gives the bounds in
Proposition~\ref{prop:cat-gap}, where the looseness is a factor $\eta^2$ rather
than a constant because the forced pulls in Lemma~\ref{lem:cat-trap} are
counted before the good arm's index has come down to its mean; a sharper
count of cold LinUCB's pulls after the good arm is well estimated, as in the
proof of Proposition~\ref{prop:cat-wrong}, would remove it, but with probes
available to cold LinUCB the required control of the estimate in a coupled
coordinate is more delicate and we have not carried it out.


\section{Proof of Theorem~\ref{thm:transfer-separation} and Numerical Checks}
\label{app:proof-transfer-sep}

\begin{proofof}{Theorem~\ref{thm:transfer-separation}}
Every offered arm is a multiple of a coordinate vector, so for every policy
$V_t$ is diagonal and each coordinate of $\widehat\theta_t$ is computed
separately. If before round $t$ the action $D$ has been played $n_D$ times with
total reward $s_D$, and $R$ has been played $n_R$ times with total reward
$s_R$, then
\[
\mathrm{UCB}_t(D)=\frac{s_D}{\lambda+n_D}+\frac{\beta}{\sqrt{\lambda+n_D}},\qquad
\mathrm{UCB}_t(R)=\frac{s_R}{\lambda+n_R}+\frac{\beta}{\sqrt{\lambda+n_R}},\qquad
\mathrm{UCB}_t(b_t)=\omega,
\]
the last because coordinate $t+2$ has not been touched before round $t$.
Since $L>1$, $\omega=\beta L/\sqrt\lambda>\beta/\sqrt{\lambda+n}$ for every
$n\ge0$. Rewards observed on a fresh arm change only its own coordinate, which
never recurs, so they do not affect the indices of $D$, $R$ or later fresh
arms. The window is a fixed sequence of types; the linear model describes the
rounds after it, and pseudo-regret is measured against $\theta^\star$
throughout.

\emph{(i)} On the window, $\Exp_A$ plays $D$ on the $k$ rounds of type $D$
(right) and the $k$ of type $R$ (wrong), and $R$ on the $2k$ of type
$\varnothing$ (wrong); $\Exp_B$ plays $R$ on the $k$ rounds of type $D$ (wrong)
and the $k$ of type $R$ (right), and $D$ on the $2k$ of type $\varnothing$
(wrong). Each plays $D$ and $R$ exactly $2k$ times, so
$M^A_\tau=M^B_\tau=2k(e_1e_1^\top+e_2e_2^\top)$, each pays $2k(p-q)$ for its
plays of $R$, and each is right $k$ times. The response sums are $ke_1$ and
$ke_2$, which gives the two estimates.

\emph{(ii)} Cold LinUCB has $n_D=n_R=0$, so $\mathrm{UCB}_t(D)=\mathrm{UCB}_t(R)
=\beta/\sqrt\lambda<\omega$ and it plays $b_t$; by induction it does so at every
round, paying $p-f$ each time. If the routed rounds contribute no rewards, then
after $\tau$ we have $s_D=s_R=0$, both indices equal $\beta/\sqrt{\lambda+2k}<\omega$,
and the same induction applies.

\emph{(iii)} After the window $\Hyb^B_\tau$ has $n_D=2k$ and $s_D=0$, because
$\Exp_B$ played $D$ only on rounds of type $\varnothing$. Hence
$\mathrm{UCB}_t(D)=\beta/\sqrt{\lambda+2k}<\omega$, and as long as $D$ is not
played these statistics do not change; so $D$ is never played. Every round
after $\tau$ then costs at least $m:=\min\{p-q,p-f\}$, and the window cost
$2k(p-q)\ge2km$, so $R_T(\Hyb^B_\tau)\ge2km+(T-4k)m=(T-2k)m$.

\emph{(iv)} Let $X_i:=\ind\{y_{\tau+i}=D\}$, which are i.i.d.\ Bernoulli$(p)$
(rounds of types $R$, $F$ and $\varnothing$ all give $D$ reward $0$),
$S_n:=\sum_{i\le n}X_i$, and
$G:=\{k+S_n>\omega(\lambda+2k+n)\ \text{for all }n\ge0\}$. On $G$, $\Hyb^A_\tau$
plays $D$ at every round after $\tau$: if it has done so at rounds
$\tau+1,\dots,\tau+n$, then $n_D=2k+n$ and $s_D=k+S_n$, so
$\mathrm{UCB}(D)>s_D/(\lambda+n_D)>\omega$; and $n_R=2k$, $s_R=0$, so
$\mathrm{UCB}(R)=\beta/\sqrt{\lambda+2k}<\omega$. It therefore plays $D$ at round
$\tau+n+1$. Every round after $\tau$ then has gap $0$, and
$R_T(\Hyb^A_\tau)=2k(p-q)$.

For the probability, $n=0$ is in $G$ because $h>0$. For $n\ge1$, put
$Y_i:=\omega-X_i\in[\omega-1,\omega]$, so that $\E Y_i=\omega-p<0$ and $G$ fails exactly when
$\sum_{i\le n}Y_i\ge k-\omega(\lambda+2k)=h$ for some $n\ge1$. By Hoeffding's lemma,
$\E e^{\theta Y_i}\le\exp(\theta(\omega-p)+\theta^2/8)$, which equals $1$ at
$\theta=8(p-\omega)>0$. So $\exp(\theta\sum_{i\le n}Y_i)$ is a nonnegative
supermartingale starting at $1$, and Ville's maximal inequality gives
$\Prb(G^c)\le e^{-\theta h}=e^{-8(p-\omega)h}$. If also $\omega\le p/2$ then
$8(p-\omega)\ge4p$, and if $\omega\le1/4$ then $h\ge(k-2\omega\lambda)/2$, so
$8(p-\omega)h\ge2p(k-2\omega\lambda)$.

\emph{(v)} Off $G$, $\Hyb^A_\tau$ plays $D$ (gap $0$) until it first plays a
fresh arm; from then on the statistics of $D$ and $R$ no longer change, both
indices stay below $\omega$, and it plays fresh arms at cost $p-f$ per round.
So $R_T(\Hyb^A_\tau)\le2k(p-q)+(p-f)T$ off $G$, and
$\E R_T(\Hyb^A_\tau)\le2k(p-q)+(p-f)T\,\Prb(G^c)$. The stated $k$ gives
$h\ge\ln T/(8(p-\omega))$, hence $\Prb(G^c)\le1/T$ and
$\E R_T(\Hyb^A_\tau)\le2k(p-q)+p-f$, while $R_T(\Cold)=(p-f)T$ is
deterministic. The ratios follow from (ii)--(iv).
\end{proofof}

\begin{remark}[Settling, and radius schedules]
\label{rem:transfer-extensions}
(a) With a constant radius the theorem also shows the all-or-nothing
behaviour seen on AmazonCat (Appendix~\ref{app:dominant}): once
$\Hyb^A_\tau$ plays a fresh arm, the statistics of $D$ stop changing and every
later fresh arm has index $\omega$, so it never plays $D$ again. It either
settles on $D$ for good or never does.
(b) Parts (i)--(v) use the radius only through the comparisons above, so they
hold for any non-increasing schedule $\beta_t$ with
$\omega:=\sup_{t>\tau}\beta_tL/\sqrt\lambda$. The all-or-nothing behaviour of
(a) needs the radius to be constant once the learner has left $D$; the
schedule of our experiments is constant after $t\approx1{,}500$.
\end{remark}

\paragraph{An expert that sees only the features.}
Theorem~\ref{thm:transfer-separation} compares two experts that both see the
type. A third expert plays $D$ and $R$ equally often, $2k$ times each, but
chooses which window rounds get $D$ without seeing the type, so that its right
answers come at the natural rate. It has the same coverage and price as
$\Exp_A$ and $\Exp_B$ and, on average, the same number of right answers. We
have no theorem for it. Numerically, with $f=0$, the parameters of the first row of
Table~\ref{tab:transfer-check} and $500$ seeds, its hybrid spends the whole
second half of the horizon on $D$ in $62\%$, $72\%$ and $84\%$ of runs at
$k=2,4,8$, and its mean regret is $576$, $412$ and $238$, against $28$, $2$
and $4$ for $\Exp_A$, $1{,}498$, $1{,}495$ and $1{,}491$ for $\Exp_B$, and
$1{,}500$ for cold LinUCB. The two informed experts sit on either side of it:
selection toward the frequent action helps far more than unselected data, and
selection away from it removes the gain.

\paragraph{Numerical checks.}
The instance is simple enough to simulate exactly. We ran LinUCB with the full
$d\times d$ design matrix and, for long horizons, the scalar recursion used in
the proof; the two chose the same action at every round of $36$ runs with
$T\le300$. With $f=0$, $\lambda=0.1$, $L=2$ and $T=5{,}000$, across $k\in\{1,\dots,32\}$
and the three settings of $(p,q,\omega)$ in Table~\ref{tab:transfer-check}: cold
LinUCB's regret equalled $pT$ in each of the $600$ runs checked;
$\Hyb^B_\tau$ never played $D$ after round $\tau$ in $12{,}000$ runs; a hybrid
that left $D$ never returned to it; when either hybrid received the routed
covariance without the routed rewards, it never played $D$ after $\tau$
($300$ runs each); and with a radius large enough that
$\beta(L-1)/\sqrt\lambda\ge1$, so $\omega>p$, neither hybrid played $D$
after $\tau$ ($200$ runs), as Remark~\ref{rem:transfer-regime} states. With
$f=0.25>p-q$ ($p=0.3$, $q=0.1$, $\omega=0.135$, $\lambda=1$, $L=2$), so that
$\Exp_B$ is right on $35\%$ of later rounds against $\Exp_A$'s $30\%$, cold
LinUCB's regret equalled $(p-f)T$ in every run, $\Hyb^B_\tau$ never played $D$
after $\tau$, its regret was never below $(T-2k)\min\{p-q,p-f\}$, and
$\Hyb^A_\tau$ settled on $D$ in $98\%$ of runs at $k=2$ and in every run at
$k\in\{4,8\}$ ($T\in\{2{,}000,20{,}000\}$, $100$--$300$ runs per cell). With $k$
chosen as in (v), the mean regret of $\Hyb^A_\tau$ stayed below
$2k(p-q)+p-f$ at $T=2{,}000$, $20{,}000$ and $200{,}000$, and the ratio of
cold LinUCB's regret to it grew from $28$ to $227$ to $1{,}923$.
Table~\ref{tab:transfer-check} compares the frequency with which
$\Hyb^A_\tau$ settles on $D$ with the bound of
Theorem~\ref{thm:transfer-separation}(iv).

\begin{table}[h]
\centering
\caption{Frequency with which $\Hyb^A_\tau$ settles on $D$ after the routed
window ($2{,}000$ seeds per cell, $\lambda=0.1$, $L=2$, $T=5{,}000$) against the
lower bound $1-e^{-8(p-\omega)h}$ of Theorem~\ref{thm:transfer-separation}(iv),
with $f=0$ and $h=k(1-2\omega)-\omega\lambda$. The bound holds in every cell.}
\label{tab:transfer-check}
\small
\begin{tabular}{@{}llccccc@{}}
\toprule
$(p,q,\omega)$ & & $k=1$ & $k=2$ & $k=3$ & $k=4$ & $k=6$ \\
\midrule
\multirow{2}{*}{$(0.3,\,0.05,\,0.135)$} & observed & $0.862$ & $0.979$ & $0.996$ & $1.000$ & $1.000$ \\
 & bound & $0.612$ & $0.852$ & $0.943$ & $0.978$ & $0.997$ \\
\midrule
\multirow{2}{*}{$(0.5,\,0.1,\,0.225)$} & observed & $0.865$ & $0.983$ & $0.996$ & $0.999$ & $1.000$ \\
 & bound & $0.687$ & $0.907$ & $0.972$ & $0.992$ & $0.999$ \\
\midrule
\multirow{2}{*}{$(0.2,\,0,\,0.09)$} & observed & $0.870$ & $0.979$ & $0.997$ & $1.000$ & $1.000$ \\
 & bound & $0.510$ & $0.762$ & $0.884$ & $0.944$ & $0.987$ \\
\bottomrule
\end{tabular}
\end{table}

\section{Numerical Validation of the Catalog}
\label{app:synthetic}

Every claim of the catalog (Section~\ref{sec:catalog} and
Appendix~\ref{app:catalog-statements}) is a statement about an explicitly
constructed instance, so each can be checked by simulating that instance. We
do so for the arrowhead identity and for each catalog row whose proof asserts
a numerical bound. The simulations use the instances of the
corresponding results, $10$ seeds, and no tuning; they are reported in
Table~\ref{tab:synthetic}. The admissible sizes are out of reach (under
\eqref{eq:conventions} $L\ge8$, so Theorem~\ref{thm:price} needs $K\ge2{,}304$
and $n\ge138{,}240$), so several blocks use smaller parameters than the proofs
require: $L=3$, $K=32$ and $n=19{,}440$ for rows~4, 5 and~7, and $K=4$,
$n=300$ for row~$5'$. Those blocks check the mechanism and the direction of
each inequality at a feasible size, not the constants. Their purpose is to
catch an error in an inequality direction, not to estimate anything.

\begin{table}[h]
\centering
\caption{Simulated instances against their proved claims. ``Claim'' is what
the cited result asserts; ``simulated'' is what the instance produces. Row
numbers refer to Table~\ref{tab:catalog}. Each block uses the instance of the
corresponding result at the sizes stated in the text, which for rows~4--7 are
below the admissible sizes, so the absolute magnitudes are
not comparable across blocks; the first four row-5 entries are the instance at
$\eps=0.002$ and the last two the $\eps$-sweep instance.}
\label{tab:synthetic}
\small
\setlength{\tabcolsep}{4pt}
\begin{tabular}{@{}llll@{}}
\toprule
Result & Quantity & Claim & Simulated \\
\midrule
Lemma~\ref{lem:arrowhead} & max entrywise error, $200$ trials & $0$ & $<10^{-10}$ \\
\midrule
Prop.~\ref{prop:three-way} & cold regret & $T=2{,}000$ & $2{,}000$ \\
 & warm regret at $m=20$ & $0$ & $0$ (all seeds) \\
Prop.~\ref{prop:one-pass} & tail regret, good expert & $T-\tau$ & $T-\tau$ exactly \\
 & tail regret, bad expert & $T-\tau$ & $T-\tau$ exactly \\
\midrule
Prop.~\ref{prop:cat-mimic} (row 4) & hybrid vs.\ cold, $\tau\le Km_0$ & identical & identical \\
 & multiplier & $\le4\beta^2+1=961$ & $0.343$ \\
\midrule
Thm.~\ref{thm:price} (row 5) & expert regret, part (i) & $\eps Kn+(T{-}Kn)=201{,}244$ & $201{,}244$ \\
 & cold regret, part (iii) & $\le K(8\beta^2{+}1)+\eps Kn=62{,}716$ & $3{,}820$ \\
 & hybrid regret, part (iv) & $\eps Kn=1{,}244$ & $1{,}244$ \\
 & non-good plays after handoff & $0$ & $0$ \\
 & handoff certificate $2\beta L\|e_j\|_{V^{-1}}$ & $<1$ & $0.941$ \\
 & $(\text{ratio}-1)\cdot\eps$, $64\times$ sweep in $\eps$ & constant & $0.0044$--$0.0039$ \\
\midrule
Prop.~\ref{prop:cat-tail-saving} (row $5'$) & regret inside window & $0$ & $0.000$ \\
 & probes played & $0$ & $0$ \\
 & tail bad pulls per direction & $[m_0,\,4\beta^2{+}1]=[2,961]$ & $93$--$94$ \\
 & tail regret & $\ge Km_0=8$ & $374$ \\
\midrule
Prop.~\ref{prop:cat-wrong} (row 7) & multiplier & $\le65$ & $0.99$ \\
 & bad pulls per direction & $\Theta(\beta^2)$ & $214$ \\
\midrule
Cor.~\ref{cor:cat-dear} (row 8) & multiplier at $\eps\ge \gamma^2/L^2$ & $\le2$ & $0.63,0.39,0.06$ \\
\bottomrule
\end{tabular}
\end{table}

Two entries carry more weight than the rest. The certificate row checks the
geometric condition of Theorem~\ref{thm:price}(iv) directly: with $K=36L^2$
and $n=36\lambda L^2/\gamma^2$ as in (iv), but at $L=3$ rather than $L\ge8$, the
handoff certificate is $0.941<1$. At the $K=32$ of the other row-5 entries it
is $2.19$, so those runs are not certified at handoff, although the hybrid
still plays $g$ on every tail round there. The $\eps$-sweep row tests the
mechanism rather than the conclusion. Theorem~\ref{thm:price} predicts that
the multiplier exceeds one by an amount of order $1/\eps$; over a $64$-fold
range of $\eps$ the product $(\text{ratio}-1)\cdot\eps$ moves only between
$0.0044$ and $0.0039$, so the predicted scaling is what the instance exhibits,
not just the predicted sign.

The gaps between claim and simulation are all in the direction the proofs
allow, and are the constants discussed in
Appendix~\ref{app:catalog}, ``Remarks on the constants'': cold regret on the
row-5 template comes in at $3{,}820$ against a bound of $62{,}716$, and the
row-4 and row-7 multipliers at $0.343$ and $0.99$ against caps of $961$ and
$65$. In every case the instance does what the row says it does---rows 4, 7
and 8 kill the multiplier, row $5'$ collects its saving only on the tail with
an empty window---while the numerical bound is loose by the stated constant
factors.

\section{Additional Experimental Results}
\label{app:more-exp}

This appendix gives the full results summarized in Section~\ref{sec:experiments}:
the decomposition and learning dynamics on AmazonCat/BERT, the complete
coverage-versus-transfer analysis, the BGE confidence sweep, the WikiLinks
results, a summary across configurations, and the comparison with other
routers.

\subsection{AmazonCat-13K with BERT Features}
\label{app:amazon-bert}

\paragraph{Choice of $c_{\mathrm{LLM}}$.} On AmazonCat the tuning seeds $11$--$20$ select
$c_{\mathrm{LLM}}=100$ for BERT features (mean Router accuracy $17.2$, $19.5$, $22.2$, $21.4$
and $5.4\%$ at $c_{\mathrm{LLM}}=1,10,10^2,10^3,10^4$) and $c_{\mathrm{LLM}}=0.1$ for BGE features
(Table~\ref{tab:bge-csweep}). On WikiLinks the Flan-T5 grid is
$c_{\mathrm{LLM}}\in\{1,10,10^2,10^3,10^4\}$ and selects $c_{\mathrm{LLM}}=10^4$
(Table~\ref{tab:wiki-csweep}); the Qwen grid is
$c_{\mathrm{LLM}}\in\{0.01,0.1,1,10,100\}$ and selects $c_{\mathrm{LLM}}=10$
(Table~\ref{tab:bge-csweep}; $c_{\mathrm{LLM}}=10$ and $100$ tie, and ties go to the
smaller value).

\begin{table}[h]
\centering
\caption{The decomposition of Proposition~\ref{prop:decomp} measured on the
BERT-768d runs, in correct answers, mean over ten seeds. $A_\tau$ is summed
over the routed rounds and $I_\tau,B_\tau$ over the rest, using the routed set
the Router actually selected (the identities hold for any routed set, by the
remark after Proposition~\ref{prop:decomp}). Every configuration collects
almost all of its advantage through $I_\tau$: the expert's own answers are
worth a few dozen correct rounds out of $1{,}578$ routed, while the data they
leave behind is worth thousands. $I_\tau$ is the only column that separates
the three experts, and the multiplying configuration is also the only one whose
$I_\tau$ is positive on every seed.}
\label{tab:amazon-decomp}
\small
\begin{tabular}{@{}lcccc@{}}
\toprule
LLM & $A_\tau$ & $I_\tau$ & $B_\tau$ & seeds with $I_\tau>0$ \\
\midrule
Flan-T5-small  & $47$ & $2{,}936\pm3{,}562$ & $4{,}582$ & $9/10$ \\
Flan-T5-base   & $27$ & $7{,}178\pm4{,}772$ & $9{,}599$ & $9/10$ \\
Flan-T5-large  & $53$ & $\mathbf{9{,}571\pm3{,}411}$ & $11{,}189$ & $\mathbf{10/10}$ \\
\bottomrule
\end{tabular}
\end{table}

\textbf{The multiplying configuration.} Flan-T5-large reaches
$27.64\pm5.56\%$ against cold LinUCB's $8.39\pm2.46\%$ with only $3.2\%$ of
rounds routed to the LLM overall; routing is concentrated in the first
$\sim5{,}000$ rounds. Table~\ref{tab:amazon-decomp} gives the shape of a
positive case directly. $A_{\tau}=53$ is negligible: it is positive, not
negative, because the routed rounds fall early, while cold LinUCB is still
near $2\%$, so even a $5\%$ expert is marginally the better player there; but
at $1{,}578$ routed rounds the whole term is worth $53$ correct answers
against a total advantage of $9{,}624$. $I_{\tau}=9{,}571$ carries $99.4\%$ of
it, and is positive on all ten seeds. $B_{\tau}=11{,}189$ is large because the
expert stays at $5\%$ while the warm-started learner passes $27\%$. This is
the row-$5'$ signature of Proposition~\ref{prop:cat-tail-saving}: the
acquisition is nearly free, and the saving is collected on the tail, after
routing has stopped. Figure~\ref{fig:amazon-dynamics} shows the same
accounting in time: every routed round falls inside the first $5{,}000$, and
$96\%$ of the $9{,}624$ correct answers the Router gains over cold LinUCB are
gained after its last routed round, while the expert itself stays frozen at
$5.3\%$. The cumulative-accuracy curves part earlier than that only because an
early lead weighs heavily in a short prefix.


\begin{figure}[h]
\centering
\includegraphics[width=0.85\linewidth]{multiplication.png}
\caption{Learning dynamics on AmazonCat-13K (BERT-768d, $c_{\mathrm{LLM}}=100$), mean over
ten seeds with shaded standard error. The shaded window contains every routed
round: all calls to the LLM fall inside the first $5{,}000$ rounds, after which
the expert is switched off. The three Flan-T5 experts (held-out accuracies
$5.20\%$, $3.61\%$ and $5.27\%$, the band at the bottom) are consulted on
equally few rounds, yet their Routers separate by nearly a factor of two; and
$96$--$100\%$ of each Router's advantage over cold LinUCB, counted in correct
answers, accrues after its last routed round. That is the row-$5'$ signature:
the saving is collected on the tail.}
\label{fig:amazon-dynamics}
\end{figure}

\subsection{Coverage and the Transfer}
\label{app:amazon-transfer}

Section~\ref{sec:two-channels} says where that difference can live. The three
experts cover the same directions to within $5\%$: summed over the routed
rounds, the coverage of the directions cold LinUCB is trapped on is
$1.40\times10^5$, $1.46\times10^5$ and $1.47\times10^5$, at prices
$9.5\times10^{-3}$, $1.03\times10^{-2}$ and $1.04\times10^{-2}$, so the price
ratio of Theorem~\ref{thm:general-price}, measured after the run, is $1.73$,
$1.59$ and $1.58$ for the three---close to one another, and ordered opposite
to the observed multipliers. We compute it from per-round logs, taking for $U$
the leading directions of the features cold LinUCB played on rounds it lost
after routing stopped, in place of certified arms. The streams are outside
the theorem's setting, so this is an explanatory statistic rather than a
bound the theorem certifies, and by construction it is blind to the
transfer. It bounds the regret ratio, and the observed regret ratios lie
below it for every expert; since the real stream is outside the theorem's
setting we read this as consistency, not confirmation. Since coverage does
not separate the experts, we test directly whether the transfer does, and
Figure~\ref{fig:two-channel} shows that it does (full numbers and tests in
Table~\ref{tab:two-channel}, Appendix~\ref{app:two-channel}).

The comparison is a replay in which the routed rounds are held fixed. For
each expert and seed, three versions of the learner run in lockstep on the
same stream: the full Router, which decides which rounds are routed, and two
ablated learners that route exactly the same rounds and play the expert's
answer there, whether or not their own confidence would have routed them. On
every other round each learner plays its own LinUCB choice and sees the true
reward. The three versions therefore have the same routed rounds, the same
routed actions and the same routed covariance, and differ only in what the
routed rounds add to the response sum; the full version reproduces the
Router of Table~\ref{tab:amazon-results} exactly. Removing $m$ collapses the
spread of the experts' mean per-seed multipliers from $1.74$ to $0.50$ and
turns every significant pairwise difference into a null result: large over
small falls from $t=5.18$ to $t=1.33$. Suppressing the response sum on the
routed rounds gives the same collapse, to a spread of $0.01$; and the two
transfer-free conditions do not differ from each other on any expert (paired
$t=-0.66$, $+0.64$, $+0.46$). Removing the transfer costs each expert in
proportion to what it gained: the per-seed multiplier falls by $2.04$ for
large ($t=4.57$), $1.27$ for base ($t=3.94$) and $0.70$ for small
($t=2.51$). Whether a run settles on the dominant category
(Appendix~\ref{app:dominant}) follows the same pattern: $2$, $8$ and $9$ of
ten runs with the transfer, and at most $2$ of ten for any expert without
it. A third route to the same conclusion
replaces the reward by $\langle\theta_{\mathrm{or}},\phi\rangle+\mathcal
N(0,\varsigma^2)$, making the model well specified so that
Proposition~\ref{prop:response-sum} applies up to the Router's adaptive
choice of routed rounds: at
$\varsigma\in\{0.001,0.01,0.499\}$ the three experts' regret ratios, measured
against the oracle's best action, span $0.019$, $0.010$ and $0.030$, against
$0.197$ on the real reward; the ratios themselves lie between $0.97$ and
$1.09$. Under linear rewards the experts are interchangeable, exactly as
predicted.

\textbf{Where the regret goes.} On the regret scale coverage and the transfer
have very different sizes. With the response sum suppressed on the routed
rounds, coverage alone gives regret ratios $1.035$, $1.029$ and $1.019$:
about $2$ to $3.5\%$ less regret than cold LinUCB, indistinguishable across
experts, and below the price ratio for each, consistent with
Theorem~\ref{thm:general-price} although the stream is outside its setting.
The full Router gives $1.070$, $1.187$ and $1.266$. For Flan-T5-base and
Flan-T5-large, therefore, almost all of the regret saved is saved by the
transfer, roughly $13$ of $16$ and $19$ of $21$ percentage points, and for
Flan-T5-small about half, $3$ of $6.5$. The separation between experts
survives the change of scale intact: on per-seed regret ratios, large beats
small with paired $t=+5.65$ and base beats small with $t=+3.32$ when the
transfer is present, and neither difference is significant with the transfer
removed ($t=+1.25$, $+1.24$) or the response sum suppressed ($t=-0.47$,
$-0.13$); Figure~\ref{fig:two-channel} shows the same
collapse on the accuracy scale. The price ratio is about the same for the
three experts, as is the coverage gain; it does not predict the size of that
gain, which it bounds loosely.

Finally, the transfer is not a small correction to the estimate at handoff; on
this data it is essentially the whole of it. The ridge oracle assigns a
correct action a score of about $2.5\times10^{-3}$ where the reward is $1$, so
the residuals dominate: $\|m\|/\|\widehat\theta_{\tau+1}\|=0.9998$, and the
identity of Definition~\ref{prop:transfer} holds on the logged trajectories
to $1.2\times10^{-14}$. Stripping $m$ makes the three estimates agree to four
decimals in held-out ranking quality ($0.2058$, $0.2058$, $0.2057$), while the
transfers themselves separate the experts with the same resolution the
multipliers have: base and large each beat small ($t=3.12$, $t=2.56$) and are
indistinguishable from one another ($t=0.36$), which is the pair the
multipliers also cannot separate.

\subsection{Coverage and the Transfer: Full Numbers}
\label{app:two-channel}

Table~\ref{tab:two-channel} gives the numbers behind
Figure~\ref{fig:two-channel}, together with the paired tests quoted in
Section~\ref{sec:amazon}.

\begin{table}[h]
\centering
\caption{Separating coverage from the transfer on AmazonCat/BERT-768d, ten seeds,
with the routed rounds held fixed: all three versions route exactly the rounds
the full Router chose and play the expert's answer there, so their routed
covariance is identical. In \emph{full} a routed round also adds its observed
reward to the response sum, as the Router does. In \emph{transfer removed} it
adds the expert-independent part: the learner is fed
$\langle\theta_{\mathrm{or}},\phi\rangle$ instead of the observed reward on
the routed rounds, which deletes $m$ of Definition~\ref{prop:transfer}. In
\emph{no response sum} it adds nothing to the response sum. On every other
round each version plays its own LinUCB choice. Multipliers are accuracy
multipliers, Router over cold LinUCB ($8.40\%$), with the regret ratio in
parentheses. The spread is the largest minus the smallest of the three
experts' mean per-seed multipliers; $t$-statistics are paired over seeds, on
each scale, $t^\star=2.26$. With the transfer present the experts separate;
without it, by either mechanism, no pairwise difference survives.}
\label{tab:two-channel}
\footnotesize
\setlength{\tabcolsep}{3pt}
\begin{tabular}{@{}lcccccc@{}}
\toprule
& \multicolumn{3}{c}{Router accuracy} & \multicolumn{3}{c}{Multiplier over cold} \\
\cmidrule(lr){2-4}\cmidrule(lr){5-7}
LLM & full & transfer removed & no response sum & full & transfer removed & no response sum \\
\midrule
Flan-T5-small & $14.37\pm6.12\%$ & $9.71\pm4.18\%$  & $11.53\pm6.61\%$ & $1.71\times$ ($1.07$) & $1.16\times$ ($1.01$) & $1.37\times$ ($1.04$) \\
Flan-T5-base  & $22.80\pm9.00\%$ & $12.56\pm5.66\%$ & $11.02\pm6.86\%$ & $2.72\times$ ($1.19$) & $1.50\times$ ($1.05$) & $1.31\times$ ($1.03$) \\
Flan-T5-large & $27.64\pm5.56\%$ & $12.49\pm6.35\%$ & $10.10\pm5.08\%$ & $3.29\times$ ($1.27$) & $1.49\times$ ($1.05$) & $1.20\times$ ($1.02$) \\
\midrule
\multicolumn{4}{@{}l}{runs settling on $1471$ (small / base / large)} & $2/8/9$ & $1/2/2$ & $1/2/1$ \\
\multicolumn{4}{@{}l}{spread of mean per-seed multipliers} & $1.74$ & $0.50$ & $0.01$ \\
\multicolumn{4}{@{}l}{large $-$ small, paired $t$} & $\mathbf{+5.18}$ & $+1.33$ & $+0.01$ \\
\multicolumn{4}{@{}l}{base $-$ small, paired $t$} & $\mathbf{+3.03}$ & $+1.36$ & $+0.03$ \\
\multicolumn{4}{@{}l}{large $-$ small, paired $t$, regret ratio} & $\mathbf{+5.65}$ & $+1.25$ & $-0.47$ \\
\multicolumn{4}{@{}l}{base $-$ small, paired $t$, regret ratio} & $\mathbf{+3.32}$ & $+1.24$ & $-0.13$ \\
\bottomrule
\end{tabular}
\end{table}

\subsection{BGE Features}
\label{app:bge}

\begin{table}[h]
\centering
\caption{Confidence sweep on the BGE-1024d configurations, tuning seeds
$11$--$20$, which select $c_{\mathrm{LLM}}$. At the
$c_{\mathrm{LLM}}=100$ used for BERT the expert width $u^{(1)}=\|\phi\|/\sqrt{c_{\mathrm{LLM}}}$ is
permanently below LinUCB's, the gate never engages, and the Router equals the
LLM exactly on every seed ($I_\tau=0$, zero variance). Lowering $c_{\mathrm{LLM}}$ raises the
expert's width and restores the gate. On AmazonCat, where cold LinUCB is
strong, the Router at $c_{\mathrm{LLM}}=0.1$ beats it while routing $0.05\%$ of rounds, and
the same $c_{\mathrm{LLM}}$ reproduces the gain on the evaluation seeds
(Table~\ref{tab:amazon-results}); at $c_{\mathrm{LLM}}\ge1$ it is dragged down toward a
weaker expert. On WikiLinks the expert is
already better than the learner, as in row~3 of Table~\ref{tab:catalog} (cold
LinUCB $18.88\%$ on these seeds, LLMs $18.72\%$, $29.47\%$ and $38.74\%$): the
selected $c_{\mathrm{LLM}}=10$ routes every round, so the Router is the expert, and at
$c_{\mathrm{LLM}}=1$, where the gate hands back, the Router falls well below the two
larger experts.}
\label{tab:bge-csweep}
\scriptsize
\setlength{\tabcolsep}{4pt}
\begin{tabular}{@{}llccccccc@{}}
\toprule
& & \multicolumn{3}{c}{routed \%} & \multicolumn{3}{c}{Router accuracy} & \\
\cmidrule(lr){3-5}\cmidrule(lr){6-8}
Dataset & $c_{\mathrm{LLM}}$ & 0.5B & 3B & 7B & 0.5B & 3B & 7B & best mult.\ \\
\midrule
\multirow{4}{*}{AmazonCat} & $0.01$ & $0.0$ & $0.0$ & $0.0$ & $34.77$ & $34.77$ & $34.77$ & $1.00\times$ \\
 & $\mathbf{0.1}$ & $0.04$ & $0.05$ & $0.05$ & $33.84$ & $39.68$ & $\mathbf{40.63}$ & $\mathbf{1.17\times}$ \\
 & $1$  & $75.2$ & $71.9$ & $73.5$ & $28.09$ & $29.26$ & $36.50$ & $1.05\times$ \\
 & $10$ & $100$ & $100$ & $100$ & $21.38$ & $22.39$ & $27.17$ & $0.78\times$ \\
\cmidrule(lr){1-9}
\multirow{5}{*}{WikiLinks} & $0.01$ & $0.0$ & $0.0$ & $0.0$ & $18.88$ & $18.88$ & $18.88$ & $1.00\times$ \\
 & $0.1$ & $0.002$ & $0.002$ & $0.002$ & $18.88$ & $18.87$ & $18.88$ & $1.00\times$ \\
 & $1$  & $13.6$ & $13.6$ & $13.7$ & $18.98$ & $20.55$ & $21.89$ & $1.01\times$ \\
 & $\mathbf{10}$ & $100$ & $100$ & $100$ & $18.72$ & $29.47$ & $\mathbf{38.74}$ & $1.00\times$ \\
 & $100$ & $100$ & $100$ & $100$ & $18.72$ & $29.47$ & $38.74$ & $1.00\times$ \\
\bottomrule
\end{tabular}
\end{table}

\textbf{BGE-1024d.} At $c_{\mathrm{LLM}}=100$ the expert width $u^{(1)}=\|\phi\|_2/\sqrt{c_{\mathrm{LLM}}}$
is permanently below LinUCB's $u^{(0)}$, so the Router defers on every round
and equals the LLM exactly: $I_\tau=0$ with zero variance across seeds. This
is a failure of the gate's hyperparameter, not evidence about $I_\tau$, and it
affected all six BGE configurations. Table~\ref{tab:bge-csweep} sweeps $c_{\mathrm{LLM}}$
down until the gate engages. The sweep is run on seeds $11$--$20$ and selects
$c_{\mathrm{LLM}}=0.1$; the same value is then evaluated on seeds $1$--$10$. There cold LinUCB
reaches $32.59\%$, while the Router reaches $40.20\%$ with Qwen-3B and
$38.71\%$ with Qwen-7B, accuracy multipliers of $1.23\times$ and $1.19\times$
(regret ratios $1.13$ and $1.10$), while routing $0.05\%$ of rounds. Both gains
are significant (paired $t=3.92$ and $2.59$ against cold, $t^\star=2.26$). With
Qwen-0.5B the gate engages just as often and the Router does not move at all
($1.00\times$, $t=0.01$), which is the expert-dependence the BERT rows also
show. The accounting has the BERT signature where there is a gain: $A_\tau$ is
$5$ and $6$ correct answers while $I_\tau$ is $3{,}801$ and $3{,}056$, positive
on $9$ and $8$ of the ten seeds, against $I_\tau=13$ for Qwen-0.5B. The
multipliers are smaller than on BERT, which is what a stronger learner implies:
when cold LinUCB already reaches a third of the correct answers, there is less
left to buy: cold LinUCB on BGE features settles on the dominant category
by itself in $7$ of $20$ seeds, so the expert has less to add
(Appendix~\ref{app:dominant}). Pushing $c_{\mathrm{LLM}}$ higher
drags the Router toward a weaker expert and below cold LinUCB. On WikiLinks
the Qwen experts are better than the learner ($29.47\%$ and $38.74\%$ against
$18.88\%$), as in row~3 of Table~\ref{tab:catalog}: with the expert above the
learner, the best the Router can do is follow it. The tuning seeds select
$c_{\mathrm{LLM}}=10$, at which the gate routes every round, so on the evaluation
seeds the Router equals the expert exactly ($I_\tau=B_\tau=0$); at
$c_{\mathrm{LLM}}=1$, the largest value at which the gate hands back, it routes
$14\%$ of rounds and reaches only $20.55\%$ and $21.89\%$ with Qwen-3B and
Qwen-7B.

\subsection{OneShotWikiLinks-311}
\label{app:wiki}

\textbf{Setup.} $622{,}000$ context--entity pairs over $311$ entities.
\emph{Flan-T5 (BERT-768d features)}: embeddings from
\texttt{bert-base-nli-mean-tokens}, joint feature $\phi(x,a)=x\odot a\in\mathbb R^{768}$,
expert Flan-T5 \citep{chung2024scaling} at 60M/250M/780M.
\emph{Qwen2.5 (BGE-1024d features)}: embeddings from
\texttt{BAAI/bge-large-en-v1.5} \citep{xiao2024c}, expert Qwen2.5-Instruct
\citep{hui2024qwen2} at 0.5B/3B/7B.

\begin{table}[h]
\centering
\caption{Final cumulative accuracy on OneShotWikiLinks-311, mean$\pm$std over
seeds $1$--$10$, at the $c_{\mathrm{LLM}}$ chosen on seeds $11$--$20$: $10^4$ for Flan-T5
(Table~\ref{tab:wiki-csweep}) and $10$ for Qwen2.5 (Table~\ref{tab:bge-csweep}).
LLM accuracy and cold LinUCB do not depend on $c_{\mathrm{LLM}}$. Mult.\ is the accuracy
multiplier Router/$\max$(LLM, Cold); the regret ratio is within $0.01$ of one
on every row. The lift column is the Router's gain over cold LinUCB in
percentage points. The Flan-T5 Routers consult the expert on $14.7\%$,
$13.4\%$ and $13.8\%$ of rounds; the Qwen2.5 Routers on every round, so they
equal the expert exactly.}
\label{tab:wiki-results}
\small
\begin{tabular}{@{}llccccc@{}}
\toprule
Family & Size & LLM & Cold LinUCB & Router & Lift (pp) & Mult.\ \\
\midrule
\multirow{3}{*}{Flan-T5} & small (60M) & $7.16\%$  & $10.71\pm 0.04\%$ & $10.71\pm 0.03\%$ & $0.00$ & $1.00\times$ \\
& base (250M)   & $10.56\%$ & $10.71\pm 0.04\%$ & $11.29\pm 0.02\%$ & $+0.58$ & $1.05\times$ \\
& large (780M)  & $12.43\%$ & $10.71\pm 0.04\%$ & $11.56\pm 0.04\%$ & $+0.85$ & $0.93\times$ \\
\midrule
\multirow{3}{*}{Qwen2.5} & 0.5B  & $18.72\%$ & $18.90\pm 0.02\%$ & $18.72\pm0.00\%$ & $-0.18$ & $0.99\times$ \\
& 3B  & $29.47\%$ & $18.90\pm 0.02\%$ & $29.47\pm0.00\%$ & $+10.57$ & $1.00\times$ \\
& 7B  & $38.74\%$ & $18.90\pm 0.02\%$ & $38.74\pm0.00\%$ & $+19.85$ & $1.00\times$ \\
\bottomrule
\end{tabular}
\end{table}

\textbf{Results.} Table~\ref{tab:wiki-results} reports final cumulative
accuracy at the selected $c_{\mathrm{LLM}}$. With Flan-T5 the Router routes $13$--$15\%$
of rounds to the LLM. Routing is driven by feature geometry, not by
reward: all three Flan-T5 sizes are routed on nearly the same share of
rounds at every $c_{\mathrm{LLM}}$ (Table~\ref{tab:wiki-csweep}) despite a large spread
in accuracy. The gain over cold LinUCB is significant for Flan-T5-base and
Flan-T5-large (paired $t=36.4$ and $52.2$, since the seeds vary little on
$622{,}000$ rounds) but small, and the Router stays below Flan-T5-large itself.

\textbf{Confidence sweep.} To rule out a suboptimal $c_{\mathrm{LLM}}$, the Flan-T5
configurations are swept over $c_{\mathrm{LLM}}\in\{1,10,10^2,10^3,10^4\}$ on the tuning seeds
$11$--$20$ (Table~\ref{tab:wiki-csweep}), and the Qwen configurations over
$c_{\mathrm{LLM}}\in\{0.01,0.1,1,10,100\}$ (Table~\ref{tab:bge-csweep}).

\begin{table}[h]
\centering
\caption{Confidence sweep on the WikiLinks Flan-T5 (BERT-768d) configurations,
tuning seeds $11$--$20$: share of rounds routed to the LLM and final Router
accuracy at each $c_{\mathrm{LLM}}$, mean over seeds. Cold LinUCB does not depend on $c_{\mathrm{LLM}}$
($10.70\%$ on these seeds); the LLMs score $7.16\%$, $10.56\%$ and $12.43\%$.
Best mult.\ is the largest Router accuracy over $\max\{$LLM, cold$\}$ among the
three experts. Mean Router accuracy over the experts is highest at
$c_{\mathrm{LLM}}=10^4$ ($11.18\%$, against $11.05\%$ at $10^3$ and $10.86\%$ at $10^2$),
which is therefore the value evaluated on seeds $1$--$10$. At that value the
Router consults the LLM on about $14\%$ of rounds and moves toward the LLM's
own accuracy rather than multiplying: it gains on the two experts that beat
cold LinUCB and stays below Flan-T5-large itself.}
\label{tab:wiki-csweep}
\scriptsize
\setlength{\tabcolsep}{4pt}
\begin{tabular}{@{}lccccccc@{}}
\toprule
& \multicolumn{3}{c}{routed \%} & \multicolumn{3}{c}{Router accuracy} & \\
\cmidrule(lr){2-4}\cmidrule(lr){5-7}
$c_{\mathrm{LLM}}$ & small & base & large & small & base & large & best mult.\ \\
\midrule
$1$ & $0.10$ & $0.09$ & $0.10$ & $10.77$ & $10.79$ & $10.82$ & $1.01\times$ \\
$10$ & $0.14$ & $0.14$ & $0.14$ & $10.78$ & $10.81$ & $10.82$ & $1.01\times$ \\
$10^2$ & $0.29$ & $0.27$ & $0.27$ & $10.83$ & $10.85$ & $10.90$ & $1.01\times$ \\
$10^3$ & $1.61$ & $1.47$ & $1.52$ & $10.97$ & $11.06$ & $11.13$ & $1.03\times$ \\
$\mathbf{10^4}$ & $14.7$ & $13.4$ & $13.8$ & $10.71$ & $11.26$ & $\mathbf{11.55}$ & $\mathbf{1.05\times}$ \\
\bottomrule
\end{tabular}
\end{table}

The decomposition explains the shape. For Flan-T5 the routed rounds do leave
usable data---$I_\tau$ is $2{,}392$, $2{,}859$ and $3{,}098$ correct answers and is
positive on all ten seeds---so the expert's data is not aimed at the wrong
directions, as in row~7 of Table~\ref{tab:catalog}. Unlike on AmazonCat,
$A_\tau$ is not negligible: it is $-2{,}395$, $732$ and $2{,}201$, and for
Flan-T5-small the expert loses on the routed rounds almost exactly what its
data gains afterwards, so the Router ends where cold LinUCB does. The gain stays small
because there is little for the expert to steer the learner toward: the most
frequent of the $311$ entities covers $0.36\%$ of contexts, whereas on
AmazonCat one category covers $29.8\%$ (Section~\ref{sec:amazon}). For Qwen the expert is far better than the
learner ($29$--$39\%$ against $19\%$): at the selected $c_{\mathrm{LLM}}$ every round
is routed, the whole difference from cold LinUCB is $A_\tau$ ($-1{,}122$,
$65{,}772$ and $123{,}441$ correct answers), and $I_\tau=B_\tau=0$.

\subsection{Summary Across Configurations}
\label{app:summary}

Twelve configurations span two datasets, two model families (60M to 7B), two
embedding spaces and four orders of magnitude in $c_{\mathrm{LLM}}$, with ten evaluation
seeds per row and every $c_{\mathrm{LLM}}$ chosen on a disjoint block of ten tuning seeds. One configuration family
multiplies clearly: AmazonCat with BERT-768d and Flan-T5 (accuracy multipliers
$1.71\times$, $2.72\times$, $3.29\times$; regret ratios $1.07$, $1.19$,
$1.27$). AmazonCat with BGE-1024d and Qwen, evaluated on seeds held out from
the choice of $c_{\mathrm{LLM}}$, gives a smaller but significant gain (up to $1.23\times$ and
$1.13$, paired $t=3.92$). The six WikiLinks configurations do not (at most
$1.05\times$); neither does BGE at the $c_{\mathrm{LLM}}$ that routes everything. In every configuration
that multiplies clearly, the accounting has the same shape: $A_\tau$ is at most a few
dozen correct answers, $I_\tau$ is in the thousands, and $B_\tau$ is large because the expert stays poor while the
warm-started learner does not. Where multiplication fails, the failing term of
Proposition~\ref{prop:decomp} can be named: on WikiLinks with Flan-T5
$I_\tau$ is positive but small, since no entity is frequent enough for the
expert to steer the learner toward; with Qwen the expert is above the learner's
ceiling, and the selected $c_{\mathrm{LLM}}$ routes every round to it;
at large $c_{\mathrm{LLM}}$ on BGE the gate never hands back and $I_\tau=0$ identically.

\subsection{Other Routers}
\label{app:routers}

Across the three Flan-T5 experts coverage is
expert-independent, as Corollary~\ref{cor:coverage-complete} requires, and the
transfer accounts for the whole of the observed separation and for most of the
regret saved (Figure~\ref{fig:two-channel}; Table~\ref{tab:two-channel} in
Appendix~\ref{app:two-channel}). The deterministic Router differs from the
randomized CORRAL master of \citet{agarwal2017corralling} and the CORRAL-style
routing of \citet{chen2024llmdecision} only in how it chooses the routed
rounds; both feed every observation to the bandit learner, so the
decomposition applies to both, and Proposition~\ref{prop:response-sum} adds
that two routers with the same routed Gram matrix leave warm starts of the
same law---any difference between routers must therefore appear in $A_\tau$ or
in which residuals their routed sets happen to collect.
Table~\ref{tab:baselines} runs both alternatives.

\begin{table}[h]
\centering
\caption{Alternative routing mechanisms on AmazonCat-13K (BERT-768d,
$c_{\mathrm{LLM}}=100$), ten seeds, paired on seed with the gate. CORRAL is the randomized
log-barrier master of \citet{agarwal2017corralling} with importance-weighted
updates; the offline arms spend $\tau$ expert pulls on contexts drawn from
\emph{outside} the online stream and then run plain LinUCB, so their
acquisition is not charged as regret and $A_\tau=0$ identically.
$\tau=1{,}600$ matches the gate's routed count. Cold LinUCB is bit-identical
across all four arms on every seed, so each comparison isolates the router.
The last column is the paired $t$ statistic for gate minus that arm over the
ten shared seeds ($t^\star=2.26$). Mult.\ is the accuracy multiplier, with
the regret ratio in parentheses.}
\label{tab:baselines}
\small
\setlength{\tabcolsep}{4pt}
\begin{tabular}{@{}llrcccrr@{}}
\toprule
LLM & Router & queries & Router acc.\ & Mult.\ & $A_\tau$ & $I_\tau$ & $t$ vs.\ gate \\
\midrule
\multirow{4}{*}{Flan-T5-small}
 & Gate (Alg.\ 1)    & $1{,}406$  & $14.37\pm6.12\%$ & $1.71\times$ ($1.07$) & $47$   & $2{,}936$ & --- \\
 & CORRAL            & $22{,}860$ & $8.87\pm1.37\%$  & $1.06\times$ ($1.01$) & $-508$ & $744$     & $\mathbf{2.64}$ \\
 & Offline $\tau{=}1{,}600$ & $1{,}600$ & $11.17\pm5.71\%$ & $1.33\times$ ($1.03$) & $0$ & $1{,}386$ & $1.16$ \\
 & Offline $\tau{=}5{,}000$ & $5{,}000$ & $15.95\pm7.11\%$ & $1.90\times$ ($1.09$) & $0$ & $3{,}778$ & $-0.59$ \\
\midrule
\multirow{4}{*}{Flan-T5-base}
 & Gate (Alg.\ 1)    & $1{,}570$ & $22.80\pm9.00\%$ & $2.72\times$ ($1.19$) & $27$  & $7{,}178$ & --- \\
 & CORRAL            & $6{,}939$ & $24.44\pm2.39\%$ & $2.91\times$ ($1.21$) & $-23$ & $8{,}045$ & $-0.63$ \\
 & Offline $\tau{=}1{,}600$ & $1{,}600$ & $20.97\pm8.89\%$ & $2.50\times$ ($1.16$) & $0$ & $6{,}286$ & $0.42$ \\
 & Offline $\tau{=}5{,}000$ & $5{,}000$ & $27.29\pm5.17\%$ & $3.25\times$ ($1.26$) & $0$ & $9{,}449$ & $-2.00$ \\
\midrule
\multirow{4}{*}{Flan-T5-large}
 & Gate (Alg.\ 1)    & $1{,}597$ & $27.64\pm5.56\%$ & $3.29\times$ ($1.27$) & $53$ & $9{,}571$ & --- \\
 & CORRAL            & $7{,}488$ & $25.42\pm2.31\%$ & $3.03\times$ ($1.23$) & $50$ & $8{,}460$ & $1.49$ \\
 & Offline $\tau{=}1{,}600$ & $1{,}600$ & $26.27\pm7.61\%$ & $3.13\times$ ($1.24$) & $0$ & $8{,}937$ & $0.42$ \\
 & Offline $\tau{=}5{,}000$ & $5{,}000$ & $25.76\pm6.53\%$ & $3.07\times$ ($1.23$) & $0$ & $8{,}681$ & $0.67$ \\
\bottomrule
\end{tabular}
\end{table}

Three things follow. First, no baseline beats the gate significantly on any
expert, and the gate beats CORRAL on Flan-T5-small ($t=2.64$) while issuing
$16\times$ fewer expert queries; on the other two experts CORRAL matches the
gate at four to five times the query budget. Second, the offline arm at a
matched budget of $1{,}600$ pulls is statistically indistinguishable from the
gate on all three experts even though its acquisition is free, so the online
routing decision costs essentially nothing---which is the content of
$A_\tau\approx0$, now visible from the other side. Tripling the free budget to
$5{,}000$ still does not separate the two. Where an alternative does come out
ahead of the gate in the mean, the margin is small, never significant, and
paid for with a much larger budget: CORRAL on Flan-T5-base gains $1.6$
points ($t=-0.63$) with $4.4\times$ the gate's expert queries, and the
offline arm at $\tau=5{,}000$ gains $1.6$ and $4.5$ points on Flan-T5-small
and Flan-T5-base ($t=-0.59$, $-2.00$) only by drawing $3$--$4\times$ as many
expert answers, all collected before the online run begins. None of these
trades buys a reliable improvement, so the gate is the cheapest router in
the comparison and no alternative is significantly better on any expert.

Third, and most directly, the decomposition survives the change of master. In
all twelve rows $I_\tau$ carries the value and $A_\tau$ is negligible beside
it, and $A_\tau$ is exactly where the choice of router shows up: it is $0$ by
construction offline, a few dozen for the gate, and \emph{negative} for CORRAL
on the two weaker experts ($-508$ and $-23$), which is the price of routing
$22{,}860$ rounds to a $5.2\%$ expert while cold LinUCB is improving. The
comparison also shows that coverage is not bought by volume: CORRAL routes
$16\times$ as many rounds as the gate on Flan-T5-small and ends with
$I_\tau=744$ against the gate's $2{,}936$. What the gate supplies is not more
expert pulls but pulls on the directions cold LinUCB is trapped on, which is
what the price $\kappa(u)$ of Section~\ref{sec:price} measures and what
Proposition~\ref{prop:response-sum} says is the whole of what coverage contributes.

\section{The Dominant Category on AmazonCat-13K}
\label{app:dominant}

This appendix gives the evidence behind the paragraph ``What the transfer
carries'' of Section~\ref{sec:amazon}. All numbers are on BERT-768d features
unless stated.

\paragraph{The label distribution.}
On the $50{,}000$ contexts of the online stream the four most frequent of the
$13{,}330$ categories are present on $29.8\%$, $17.0\%$, $11.3\%$ and $7.8\%$
of contexts (category identifiers $1471$, $7961$, $7892$, $7891$). A policy
that always predicts the first scores $29.77\%$; it uses the label
distribution, which an online learner is not given. Context-free bandits that
must discover it do not: over the same stream, $50{,}000$ rounds and seeds
$1$--$10$, UCB1 (one pull of every category first) scores $0.04\%$ and an
optimistic empirical-mean greedy rule $2.31\%$.

\paragraph{Where a run settles.}
We say a run \emph{settles} on a category if that category accounts for more
than $20\%$ of the learner's plays in the last fifth of the horizon. For the
dominant category the classification does not depend on the threshold: on the
runs of Table~\ref{tab:settle}, moving it anywhere between $10\%$ and $30\%$
changes the number of runs that settle on $1471$ by at most one on BERT
features and not at all on BGE features. The counts for the less frequent
categories are more sensitive on BERT, where runs that settle on $7892$
predict it on about a fifth of late rounds. Table~\ref{tab:settle}
reports where runs settle and how accurate the runs that settle on each
category are. On BERT features cold LinUCB settles on the dominant category
in $1$ of $20$ seeds; the two multiplying experts move it there in $17$ of
$20$. On BGE features cold LinUCB already settles there in $7$ of $20$ seeds,
which is why the expert has less to add; there the runs that settle on the
dominant category reach $41.2\%$, above the $29.8\%$ of always predicting
it, so the learner also uses the context. Knowing only where each run
settles predicts mean Router accuracy well: predicting a configuration's
accuracy as the frequency of each settled category times the mean accuracy
of runs settling there in the \emph{other} configurations gives $R^2=0.93$
over $25$ BERT configurations and $R^2=0.93$ over $7$ BGE configurations.
These are every configuration that routes fewer than half the rounds: on
BERT, the gate at $c_{\mathrm{LLM}}\in\{1,10,10^2,10^3\}$, CORRAL, the two offline arms of
Table~\ref{tab:baselines} and the transfer-removed ablation of
Table~\ref{tab:two-channel}, for each of the three experts, plus cold
LinUCB; on BGE, the gate at $c_{\mathrm{LLM}}\in\{0.01,0.1\}$ for each expert, plus cold
LinUCB.

\begin{table}[h]
\centering
\caption{Where runs settle on AmazonCat-13K (seeds $1$--$20$; BERT at $c_{\mathrm{LLM}}=100$,
BGE at $c_{\mathrm{LLM}}=0.1$). Each cell counts runs that settle on that category; the
percentage under each category is its frequency in the stream. The bottom
block gives the mean final accuracy of the runs, pooled over the four rows,
that settle on each category.}
\label{tab:settle}
\small
\setlength{\tabcolsep}{4pt}
\begin{tabular}{@{}llcccccc@{}}
\toprule
& & $1471$ & $7961$ & $7892$ & $7891$ & & \\
Features & Policy & ($29.8\%$) & ($17.0\%$) & ($11.3\%$) & ($7.8\%$) & none & Accuracy \\
\midrule
\multirow{4}{*}{BERT-768d} & Cold LinUCB & $1$ & $0$ & $9$ & $1$ & $9$ & $8.34\%$ \\
& Flan-T5-small & $3$ & $0$ & $14$ & $0$ & $3$ & $13.71\%$ \\
& Flan-T5-base & $17$ & $0$ & $3$ & $0$ & $0$ & $25.10\%$ \\
& Flan-T5-large & $17$ & $0$ & $3$ & $0$ & $0$ & $26.90\%$ \\
\midrule
\multirow{4}{*}{BGE-1024d} & Cold LinUCB & $7$ & $12$ & $0$ & $0$ & $1$ & $33.68\%$ \\
& Qwen-0.5B & $8$ & $9$ & $0$ & $2$ & $1$ & $33.23\%$ \\
& Qwen-3B & $17$ & $2$ & $0$ & $1$ & $0$ & $39.94\%$ \\
& Qwen-7B & $16$ & $3$ & $0$ & $1$ & $0$ & $39.67\%$ \\
\midrule
\multicolumn{2}{@{}l}{Accuracy of settled runs, BERT} & $28.08\%$ & --- & $11.61\%$ & $6.24\%$ & $5.93\%$ & \\
\multicolumn{2}{@{}l}{Accuracy of settled runs, BGE} & $41.21\%$ & $30.84\%$ & --- & $24.66\%$ & $25.96\%$ & \\
\bottomrule
\end{tabular}
\end{table}

\paragraph{A handful of routed answers decides it.}
Table~\ref{tab:dose} groups every Router run that routes fewer than half the
rounds by the number $k$ of routed rounds on which the expert proposed the
dominant category and was correct, and reports the fraction that settle on
it. On BERT features the fraction rises from $7\%$ at $k=0$ to $80\%$ at
$k=4$, and every one of the $54$ runs with $k\ge5$ settles there. On BGE
features the fraction at $k=0$ is already $43\%$, the rate at which the
learner finds the category by itself, and again every run with $k\ge5$
settles there. The pooled runs span experts, values of $c_{\mathrm{LLM}}$ and, on BERT, the
CORRAL master, so the curve is a property of the learner and the data rather
than of one expert. The three BERT experts differ in exactly this count: in
the $c_{\mathrm{LLM}}=100$ runs the mean $k$ is $0.20$ for Flan-T5-small, $3.70$ for
Flan-T5-base and $4.00$ for Flan-T5-large.

\begin{table}[h]
\centering
\caption{Runs that settle on the dominant category, by the number $k$ of
correct routed answers on it. BERT: $177$ runs ($c_{\mathrm{LLM}}\in\{1,10,10^2,10^3\}$ and
CORRAL). BGE: $99$ runs ($c_{\mathrm{LLM}}\in\{0.01,0.1\}$).}
\label{tab:dose}
\small
\begin{tabular}{@{}lcccccc@{}}
\toprule
$k$ & $0$ & $1$ & $2$ & $3$ & $4$ & $\ge5$ \\
\midrule
BERT-768d & $3/45$ ($7\%$) & $11/29$ ($38\%$) & $9/17$ ($53\%$) & $13/17$ ($76\%$) & $12/15$ ($80\%$) & $54/54$ ($100\%$) \\
BGE-1024d & $25/58$ ($43\%$) & $3/3$ & $5/9$ & $6/8$ & $8/9$ & $12/12$ ($100\%$) \\
\bottomrule
\end{tabular}
\end{table}

\paragraph{Controlled test.}
Table~\ref{tab:inject} changes the expert's answers on a fixed set of $15$
contexts and nothing else. The runs use an exact re-implementation of the
Router and of cold LinUCB, checked to make the same choice as the
experimental runner on every round of a $3{,}000$-round comparison, on
$10{,}000$ held-out contexts that do not appear in the online stream, with a
horizon of $10{,}000$ rounds and $20$ seeds; the seed fixes the order of the
contexts and is shared across conditions, so every comparison is paired.
Multipliers are relative to cold LinUCB on the same runs ($4.63\%$).
\emph{$+15$ on $1471$} replaces $15$ wrong answers of Flan-T5-small by the
dominant category on contexts that contain it. \emph{$+15$ random} gives the
same number of correct answers, on the same contexts where possible, on a
uniformly chosen true category other than $1471$. \emph{$+15$ on $7961$}
gives $15$ correct answers on the second most frequent category.
\emph{$-15$ on $1471$} replaces the $15$ correct dominant-category answers of
Flan-T5-large by wrong answers. The first three conditions raise held-out
accuracy by the same $0.15$ points; only the category differs.
\emph{$+15$ on $1471$, reward hidden} uses the same answers as \emph{$+15$ on
$1471$}, but when the Router routes one of the $15$ changed contexts the
learner adds the answer's feature vector to its design matrix and nothing to
its response sum: the change adds coverage along the dominant category and no
transfer. Only $2.15$ of the $15$ changed answers are routed per seed on
average (none in one seed), so the whole effect of the intervention passes
through two or three routed rounds.

\begin{table}[h]
\centering
\caption{Changing $15$ of the expert's $10{,}000$ held-out answers. Tests are
paired over $20$ seeds against the unmodified expert: a paired $t$-test on
Router accuracy and an exact McNemar test on whether the run settles on
$1471$.}
\label{tab:inject}
\small
\setlength{\tabcolsep}{4pt}
\resizebox{\linewidth}{!}{%
\begin{tabular}{@{}llccccc@{}}
\toprule
Expert & Change & Expert acc. & Router acc. & Mult. & Settle on $1471$ & $p$ (acc.\ / settle) \\
\midrule
\multirow{5}{*}{Flan-T5-small} & none & $5.34\%$ & $6.14\%$ & $1.33\times$ & $1/20$ & --- \\
& $+15$ random & $5.49\%$ & $7.63\%$ & $1.65\times$ & $2/20$ & $0.15$ / $1.00$ \\
& $+15$ on $7961$ & $5.49\%$ & $9.39\%$ & $2.03\times$ & $3/20$ & $0.010$ / $0.50$ \\
& $+15$ on $1471$ & $5.49\%$ & $\mathbf{12.10\%}$ & $\mathbf{2.62\times}$ & $\mathbf{8/20}$ & $\mathbf{0.002}$ / $\mathbf{0.016}$ \\
& $+15$ on $1471$, reward hidden & $5.49\%$ & $6.83\%$ & $1.48\times$ & $2/20$ & $0.47$ / $1.00$ \\
\midrule
\multirow{2}{*}{Flan-T5-large} & none & $5.32\%$ & $13.79\%$ & $2.98\times$ & $10/20$ & --- \\
& $-15$ on $1471$ & $5.17\%$ & $10.09\%$ & $2.18\times$ & $6/20$ & $0.067$ / $0.39$ \\
\bottomrule
\end{tabular}}
\end{table}

Hiding the reward of the added answers removes the gain: the Router falls
back to $1.48\times$, not significantly different from the unmodified expert,
and significantly below the same answers with their rewards (paired $t=3.46$,
$p=0.003$; settling $8/20$ against $2/20$, McNemar $p=0.031$). Since only
about two changed answers are routed per run, the comparison is between two
or three rewarded answers on the dominant category, which are enough to move
the learner there and double the multiplier, and the same two or three
answers with their rewards hidden, which leave the learner where it was. What
the learner uses is the reward, not the coverage along the category. At equal
accuracy, the gain grows with the frequency of the category the
added answers point to: none for random categories, a significant gain for
the second most frequent category (where $5$ of $20$ runs settle on that
category instead), and the largest gain for the dominant one. Removing the
dominant-category answers from Flan-T5-large lowers its multiplier but not
significantly; its other answers still lead $6$ of $20$ runs to the
dominant category. Individual seeds are sensitive to any change in the
routed data, which can tip a run toward or away from a category, so we
compare rates over seeds rather than reading individual pairs.

\paragraph{The same test on the evaluation stream.}
Table~\ref{tab:inject50k} repeats the test on the $50{,}000$-round stream and
seeds $1$--$10$ of Table~\ref{tab:amazon-results}, with the same simulator.
To keep the same share of the stream, $75$ of Flan-T5-small's wrong answers
are changed, again $+0.15$ points of expert accuracy ($5.20\%$ to $5.35\%$);
the unmodified row is the Router of Table~\ref{tab:amazon-results}, and cold
LinUCB is the same $8.40\%$. The conditions are those of
Table~\ref{tab:inject}: \emph{$+75$ on $1471$}, the same answers with their
rewards hidden, and $+75$ correct answers on random categories other than
$1471$. The Router routes $2.6$ of the $75$ changed answers per seed on
average (none in one seed), about the same number as on the held-out stream,
since it consults the expert almost only in the first few thousand rounds.
The result is the same: rewarded answers on the dominant category raise the
multiplier from $1.71\times$ to $2.56\times$ and settling from $2$ to $7$ of
ten runs, while the same answers with their rewards hidden give $1.84\times$,
not significantly different from the unmodified expert and significantly
below the same answers with their rewards (paired $t=2.36$, $p=0.042$;
settling $7/10$ against $3/10$, McNemar $p=0.13$). Answers on random
categories give no significant gain either. With ten seeds the settling
counts alone are not significant; the accuracy tests are.

\begin{table}[h]
\centering
\caption{Changing $75$ of the expert's answers on the $50{,}000$-round
evaluation stream, seeds $1$--$10$. Tests are paired over the ten seeds
against the unmodified expert: a paired $t$-test on Router accuracy and an
exact McNemar test on whether the run settles on $1471$. Multipliers and
regret ratios are relative to cold LinUCB on the same seeds ($8.40\%$).}
\label{tab:inject50k}
\small
\setlength{\tabcolsep}{4pt}
\resizebox{\linewidth}{!}{%
\begin{tabular}{@{}lcccccc@{}}
\toprule
Change & Expert acc. & Router acc. & Mult.\ (regret ratio) & Settle on $1471$ & Routed changed & $p$ (acc.\ / settle) \\
\midrule
none & $5.20\%$ & $14.37\pm6.12\%$ & $1.71\times$ ($1.07$) & $2/10$ & --- & --- \\
$+75$ random & $5.35\%$ & $15.42\pm9.41\%$ & $1.84\times$ ($1.08$) & $3/10$ & $2.4$ & $0.63$ / $1.00$ \\
$+75$ on $1471$ & $5.35\%$ & $\mathbf{21.50\pm6.80\%}$ & $\mathbf{2.56\times}$ ($\mathbf{1.17}$) & $\mathbf{7/10}$ & $2.6$ & $\mathbf{0.007}$ / $0.063$ \\
$+75$ on $1471$, reward hidden & $5.35\%$ & $15.45\pm5.71\%$ & $1.84\times$ ($1.08$) & $3/10$ & $2.6$ & $0.61$ / $1.00$ \\
\bottomrule
\end{tabular}}
\end{table}

\section{Confidence-Gated Meta-Algorithm: Full Description}
\label{app:meta-algorithm}

This appendix expands the brief description of
Section~\ref{sec:setup} into the full round-by-round walkthrough,
states the deterministic confidence-gated routing rule, and provides
the pseudocode used in our experiments. The Router is the adaptive-$\tau$ version of the handoff hybrid
$\Hyb_\tau$ of Section~\ref{sec:preliminaries}: it chooses, round by round,
whether the expert or LinUCB acts, and feeds every observation to LinUCB.
Proposition~\ref{thm:handoff} applies to it with the acquisition term summed over
expert-routed rounds.

Rather than encoding prior knowledge through a single improved
initialization, we assume access to \emph{multiple frozen policy
estimators} $\{(\widehat\theta_i,V^{(i)})\}_{i=1}^m$. Each pair should
be thought of as a constant policy obtained offline:
$\widehat\theta_i$ is a fixed linear predictor and $V^{(i)}$ is the
corresponding design matrix that captures the geometry of the data on
which that predictor was trained. Our goal is not to compete with the
single best expert in hindsight, an objective closer in spirit to
online model-selection methods such as CORRAL
\citep{agarwal2017corralling,chen2024llmdecision}, which provide
robustness guarantees of the form ``perform nearly as well as the best
base algorithm.'' Here we want to collect the data-value term $I_\tau$ of
Proposition~\ref{prop:decomp} when it is available while paying as little
acquisition regret as possible, and to fall back to LinUCB otherwise.

\paragraph{Setting.}
At each round $t=1,\dots,T$, the learner observes a finite action set
$\mathcal A_t$ with feature vectors
$\{x_{t,a}\}_{a\in\mathcal A_t}\subseteq\mathbb R^d$, selects an action
$A_t\in\mathcal A_t$, and observes reward
$r_t = \langle\theta^\star,x_{t,A_t}\rangle + \eta_t$,
with $\eta_t$ conditionally mean-zero sub-Gaussian.
We are given $m$ frozen estimators
$\{(\widehat\theta_i,V^{(i)})\}_{i=1}^m$ and an online LinUCB learner
with state $(V_t,c_t,\widehat\theta_t)$, where
$\widehat\theta_t = V_t^{-1}c_t$. The frozen estimators are never
updated; only LinUCB is updated online.

\paragraph{Expert proposals and geometry-aware confidence.}
Each frozen expert $i\in[m]$ proposes the greedy action
$a_t^{(i)}\in\arg\max_{a\in\mathcal A_t}\langle\widehat\theta_i,x_{t,a}\rangle$.
To measure how trustworthy expert $i$ is on the current round, we
evaluate the uncertainty of its proposed action under its own design
geometry,
\[
u_t^{(i)} \;:=\; \bigl\|x_{t,a_t^{(i)}}\bigr\|_{(V^{(i)})^{-1}},
\qquad
\|x\|_{M^{-1}} \;:=\; \sqrt{x^\top M^{-1}x}.
\]
The interpretation is simple. If the proposed action lies in directions
that are well covered by the offline data underlying expert $i$, then
$u_t^{(i)}$ is small and the expert is more trustworthy; if the current
context points in directions poorly represented in $V^{(i)}$, then
$u_t^{(i)}$ is large and the expert should be trusted less. We treat
LinUCB itself as an additional candidate, indexed by $0$, with proposal
$a_t^{(0)}\in\arg\max_a\langle\widehat\theta_t,x_{t,a}\rangle+\alpha_t\|x_{t,a}\|_{V_t^{-1}}$
and uncertainty score $u_t^{(0)} := \|x_{t,a_t^{(0)}}\|_{V_t^{-1}}$.

\paragraph{Deterministic confidence-gated selection.}
The natural rule is to prefer the frozen expert whose proposed action
has the smallest uncertainty score, unless the online learner is
already at least as confident. Concretely, with
$i_t^\star\in\arg\min_{i\in[m]}u_t^{(i)}$,
\[
I_t \;=\;
\begin{cases}
i_t^\star, & \text{if }u_t^{(i_t^\star)}<u_t^{(0)},\\[2pt]
0, & \text{otherwise},
\end{cases}
\qquad A_t \;:=\; a_t^{(I_t)}.
\]
The master is entirely deterministic. Early in learning, LinUCB
typically has a large confidence radius, so a frozen expert may be
preferred; as data accumulate, $V_t$ grows and the LinUCB uncertainty
shrinks, and eventually the adaptive learner becomes the most confident
option and takes over automatically.

\paragraph{Optional permanent handover to LinUCB.}
To formalize the idea that LinUCB should eventually dominate, we include
a one-way switching criterion. With upper-confidence scores
$U_t^{(i)} = \langle\widehat\theta_i,x_{t,a_t^{(i)}}\rangle + \alpha_i\|x_{t,a_t^{(i)}}\|_{(V^{(i)})^{-1}}$
for each frozen expert and lower-confidence score
$L_t^{(0)} = \langle\widehat\theta_t,x_{t,a_t^{(0)}}\rangle - \alpha_t\|x_{t,a_t^{(0)}}\|_{V_t^{-1}}$
for LinUCB, if $L_t^{(0)} \ge \max_{i\in[m]} U_t^{(i)}$ for $K$
consecutive rounds we permanently hand control to LinUCB.

\paragraph{The instance used in our experiments.}
All experiments of Section~\ref{sec:experiments} use a single frozen expert
($m=1$), the LLM, with $V^{(1)}=c_{\mathrm{LLM}}I_d$, and do not use the permanent handover:
the gate is evaluated on every round.

\begin{algorithm2e}[!htb]
\caption{Confidence-Gated Master over Frozen Policies and LinUCB}
\label{alg:confidence-gated}
\DontPrintSemicolon
\LinesNumbered
\KwIn{Frozen estimators $\{(\widehat{\theta}_i,V^{(i)})\}_{i=1}^m$, regularization $\lambda>0$, confidence parameters $\{\alpha_t\}_{t=1}^T$}
Initialize LinUCB state: $V_1=\lambda I_d$, $c_1=0$, $\widehat{\theta}_1=0$\;
Set $\textsc{handover}\gets \mathrm{False}$\;
\For{$t=1,\dots,T$}{
  Observe action set $\mathcal{A}_t$\;
  Compute LinUCB proposal
  $a_t^{(0)} \in \arg\max_{a\in\mathcal{A}_t}\bigl\{\langle \widehat{\theta}_t,x_{t,a}\rangle+ \alpha_t \|x_{t,a}\|_{V_t^{-1}}\bigr\}$\;
  Set $u_t^{(0)} \gets \|x_{t,a_t^{(0)}}\|_{V_t^{-1}}$\;
  \For{$i=1,\dots,m$}{
    Compute frozen proposal $a_t^{(i)} \in \arg\max_{a\in\mathcal{A}_t}\langle \widehat{\theta}_i,x_{t,a}\rangle$\;
    Set $u_t^{(i)} \gets \|x_{t,a_t^{(i)}}\|_{(V^{(i)})^{-1}}$\;
  }
  \eIf{$\textsc{handover}=\mathrm{True}$}{
    $A_t \gets a_t^{(0)}$\;
  }{
    Choose $i_t^\star \in \arg\min_{i\in[m]} u_t^{(i)}$\;
    \eIf{$u_t^{(i_t^\star)} < u_t^{(0)}$}{
      $A_t \gets a_t^{(i_t^\star)}$\;
    }{
      $A_t \gets a_t^{(0)}$\;
    }
  }
  Observe reward $r_t$\;
  Update LinUCB:
  $V_{t+1}\gets V_t + x_{t,A_t}x_{t,A_t}^\top$,
  $c_{t+1}\gets c_t + r_t x_{t,A_t}$,
  $\widehat{\theta}_{t+1}\gets V_{t+1}^{-1}c_{t+1}$\;
  Compute $U_t^{(i)}$ for all $i\in[m]$ and $L_t^{(0)}$\;
  \If{$L_t^{(0)} \ge \max_{i\in[m]} U_t^{(i)}$ \textup{for $K$ consecutive rounds}}{
    $\textsc{handover}\gets \mathrm{True}$\;
  }
}
\end{algorithm2e}

\paragraph{Distinction from randomized model selection.}
The deterministic, sample-by-sample geometric routing distinguishes
Algorithm~\ref{alg:confidence-gated} from CORRAL-style randomized
masters in two ways. First, it can exploit context-level information
about which estimator is reliable on \emph{this particular} action
set; second, because it is deterministic, the choice of which rounds to route
is a function of the current widths rather than of a sampled index. Both
rules feed the played action and the observed reward into $(V_t,c_t)$, so
both produce the online warm start $(M_\tau,c_\tau)$ whose value is
$I_\tau$ in Proposition~\ref{prop:decomp}; the decomposition itself does not
distinguish them, and any difference between them shows up in which rounds
are routed, i.e., in $A_\tau$ and in the composition of $M_\tau$.


\section{Pre-Runtime Estimates}
\label{app:diagnostics}

We tried two ways of anticipating the outcome of a handoff from a labeled
held-out sample of $n$ contexts and the expert's proposals on it, before any
online round is played. The first is built only from the expert's held-out
Gram matrix
$\widehat M_{\Exp}=\sum_{s\le n}x_{s,a_s^{\Exp}}x_{s,a_s^{\Exp}}^\top$, that is,
from the coverage its answers would provide; the second also uses whether
those answers are right.

\paragraph{Width-limited regret.} The held-out average
$\widehat\Phi(V)=\frac1n\sum_s\max_a\min\{\widehat\Delta_{s,a},2\beta\|x_{s,a}\|_{V^{-1}}\}$
of the quantity in Lemma~\ref{lem:width-limited}, compared at
$V=\lambda I$ and $V=\lambda I+(\tau/n)\widehat M_{\Exp}$, estimates how much
the expert's data would lower the regret LinUCB can still incur. On AmazonCat
it cannot move: with $13{,}330$ actions in $768$ dimensions the width term
exceeds the unit gap for every held-out context, so $\widehat\Phi$ equals
$1.0000$ under both designs, for every expert and every
$\tau\in\{1000,2000,5000,10000\}$.

\paragraph{Dry run.} Replaying the held-out contexts (resampled to
$T'=50{,}000$ rounds, ten seeds) and simulating cold LinUCB and LinUCB
warm-started from $\tau=1{,}600$ expert rounds ranks the three Flan-T5 experts
correctly (regret ratios $1.021$, $1.057$ and $1.084$; large over small,
paired $t=2.91$) but underestimates the size of the effect by a factor of
about $3.2$ on the regret scale (excesses over one of $2.2\%$, $5.7\%$ and
$8.4\%$ against observed $7.0\%$, $18.7\%$ and $26.6\%$).

\paragraph{What each estimate can see.} The width-limited estimate sees only
coverage, and on this data the experts cover the same directions to within
$5\%$ (Section~\ref{sec:amazon}); what separates them is the transfer, which
depends on the residuals at the actions the expert chooses and is not a
function of $\widehat M_{\Exp}$. The dry run does see the transfer, because
its warm start adds the expert's observed rewards to the response sum, and
that is why it ranks the experts correctly. It still underestimates the size
of the effect, and it needs labeled data and the expert's answers on them,
which is most of the cost of running the experiment. Neither estimate
therefore predicts the size of the multiplier cheaply, before the expert is
run.

\end{document}
